\documentclass[11pt,a4paper]{article}
\usepackage[utf8]{inputenc}
\usepackage[T1]{fontenc}
\usepackage[spanish,es-nodecimaldot,es-noshorthands]{babel}
\usepackage{amsmath,amssymb,amsthm,amsfonts,mathrsfs}
\usepackage{geometry}
\geometry{top=2.5cm,bottom=2.5cm,left=2.5cm,right=2.5cm}
\usepackage{setspace}
\onehalfspacing
\usepackage{xcolor}
\usepackage{tikz}
\usetikzlibrary{arrows.meta,calc,decorations.pathmorphing,patterns}
\tikzset{>=Stealth}
\usepackage{hyperref}
\hypersetup{colorlinks=true,linkcolor=blue!75!black,citecolor=blue!75!black,urlcolor=blue!75!black}
\usepackage{cleveref}

\theoremstyle{definition}
\newtheorem{definicion}{Definición}[section]
\newtheorem{teorema}{Teorema}[section]
\newtheorem{proposicion}{Proposición}[section]
\newtheorem{lema}{Lema}[section]
\newtheorem{corolario}{Corolario}[section]
\newtheorem{ejemplo}{Ejemplo}[section]

\theoremstyle{remark}
\newtheorem{observacion}{Comentario}[section]

\crefname{definicion}{definición}{definiciones}
\Crefname{definicion}{Definición}{Definiciones}
\crefname{teorema}{teorema}{teoremas}
\Crefname{teorema}{Teorema}{Teoremas}
\crefname{proposicion}{proposición}{proposiciones}
\Crefname{proposicion}{Proposición}{Proposiciones}
\crefname{lema}{lema}{lemas}
\Crefname{lema}{Lema}{Lemas}
\crefname{corolario}{corolario}{corolarios}
\Crefname{corolario}{Corolario}{Corolarios}
\crefname{observacion}{comentario}{comentarios}
\Crefname{observacion}{Comentario}{Comentarios}
\crefname{ejemplo}{ejemplo}{ejemplos}
\Crefname{ejemplo}{Ejemplo}{Ejemplos}

\numberwithin{equation}{section}

% Comandos matemáticos útiles
\newcommand{\CC}{\mathbb{C}}
\newcommand{\RR}{\mathbb{R}}
\newcommand{\NN}{\mathbb{N}}
\newcommand{\ZZ}{\mathbb{Z}}
\newcommand{\HH}{\mathcal{H}}
\newcommand{\OO}{\mathcal{O}}
\newcommand{\mm}{\mathfrak{m}}
\newcommand{\pd}[2]{\frac{\partial #1}{\partial #2}}
\newcommand{\pdd}[2]{\frac{\partial^2 #1}{\partial #2^2}}
\newcommand{\dz}{\mathrm{d}z}
\newcommand{\dw}{\mathrm{d}w}
\newcommand{\dt}{\mathrm{d}t}
\DeclareMathOperator{\diam}{diam}
\DeclareMathOperator{\ord}{ord}
\DeclareMathOperator{\dist}{dist}
\DeclareMathOperator{\Int}{Int}
\DeclareMathOperator{\supp}{supp}

\setlength{\parindent}{0pt}
\setlength{\parskip}{6pt}

\title{\textbf{Notas de Estudio: Análisis Complejo en una Variable}\\
\large Basadas en la obra de Raghavan Narasimhan e Yves Nievergelt (2da Edición, 2001)\\[0.2cm]
\textbf{Capítulos 1 y 2: Teoría Elemental y Espacios Recubridores}}
\author{\textbf{Leonardo Montalvo Osores}}
\date{\today}

\begin{document}

\maketitle

\begin{abstract}
El presente documento contiene las notas de estudio sistemáticas de los Capítulos 1 y 2 del texto clásico \emph{Complex Analysis in One Variable} de Raghavan Narasimhan e Yves Nievergelt (2da edición, Birkhäuser, 2001). Se recopilan con exhaustividad y rigor matemático todas las definiciones, proposiciones, teoremas, lemas y corolarios con sus demostraciones completas y detalladas, junto con todos los comentarios conceptuales e históricos íntegros, traducidos fielmente al español. Las demostraciones se presentan de forma rigurosa, justificando cada paso, y los ejemplos van acompañados de diagramas TikZ para facilitar la comprensión geométrica.
\end{abstract}

\tableofcontents
\newpage


\section{Capítulo 1: Teoría Elemental de Funciones Holomorfas}

En este capítulo se desarrolla la teoría clásica de las funciones holomorfas de una variable compleja. Se establece la equivalencia fundamental entre la diferenciabilidad en sentido complejo ($\mathbb{C}$-diferenciabilidad) y la analiticidad (representabilidad local mediante series de potencias), culminando en el Teorema de Cauchy-Goursat y sus vastas consecuencias. A lo largo de todo el capítulo, $\Omega$ denotará un conjunto abierto en $\mathbb{C}$.

\begin{definicion}[$\mathbb{C}$-diferenciabilidad]
Sea $f$ una función con valores complejos definida sobre un abierto $\Omega \subset \mathbb{C}$, y sea $a \in \Omega$. Se dice que $f$ es \textbf{$\mathbb{C}$-diferenciable} en el punto $a$ si el límite
\begin{equation}
\lim_{z \to a} \frac{f(z) - f(a)}{z - a}
\end{equation}
existe. Cuando este límite existe, se denota por $f'(a)$ (o por $\frac{\mathrm{d}f}{\mathrm{d}z}(a)$) y se denomina la \textbf{derivada} de $f$ en $a$. 

Se dice que $f$ es $\mathbb{C}$-diferenciable sobre $\Omega$ si es $\mathbb{C}$-diferenciable en cada punto $a \in \Omega$. En tal caso, la función $a \mapsto f'(a)$, denotada por $f'$, se llama la derivada de $f$ sobre $\Omega$.
\end{definicion}

\begin{definicion}[Función holomorfa]
Sea $f$ una función con valores complejos definida sobre un abierto $\Omega \subset \mathbb{C}$. Se dice que $f$ es \textbf{holomorfa} sobre $\Omega$ si para cada punto $a \in \Omega$ existe un entorno $U \subset \Omega$ de $a$ y una sucesión de números complejos $\{c_n\}_{n=0}^\infty$ tales que, para todo $z \in U$, la serie de potencias
\begin{equation}
\sum_{n=0}^\infty c_n (z - a)^n
\end{equation}
converge a $f(z)$.
\end{definicion}

\begin{observacion}[Equivalencia fundamental]
El resultado vertebral de las primeras secciones es demostrar que una función es holomorfa sobre un abierto $\Omega$ si y sólo si es $\mathbb{C}$-diferenciable sobre $\Omega$ (Teorema de Cauchy-Goursat). Así, la existencia de la primera derivada compleja implica automáticamente la analiticidad y la infinitud de derivadas.
\end{observacion}

\subsection{1.1. Propiedades básicas de funciones $\mathbb{C}$-diferenciables y holomorfas}

\begin{proposicion}
Sea $\Omega \subset \mathbb{C}$ abierto y $f: \Omega \to \mathbb{C}$. Si $f$ es $\mathbb{C}$-diferenciable en $a \in \Omega$, entonces $f$ es continua en $a$.
\end{proposicion}
\begin{proof}
Para $z \in \Omega \setminus \{a\}$, se tiene la identidad $f(z) - f(a) = \frac{f(z) - f(a)}{z - a} (z - a)$. Tomando el límite cuando $z \to a$:
\[
\lim_{z \to a} (f(z) - f(a)) = f'(a) \cdot 0 = 0,
\]
de modo que $\lim_{z \to a} f(z) = f(a)$, demostrando la continuidad de $f$ en $a$.
\end{proof}

\begin{definicion}[Operadores diferenciales de Wirtinger]
Sea $f$ una función con valores complejos sobre $\Omega$, diferenciable en sentido real ($f(x,y) = u(x,y) + iv(x,y)$ con $x,y \in \mathbb{R}$). Se definen formalmente los operadores diferenciales:
\begin{equation}
\frac{\partial}{\partial z} = \frac{1}{2} \left( \frac{\partial}{\partial x} - i \frac{\partial}{\partial y} \right), \qquad \frac{\partial}{\partial \bar{z}} = \frac{1}{2} \left( \frac{\partial}{\partial x} + i \frac{\partial}{\partial y} \right).
\end{equation}
\end{definicion}

\begin{observacion}[Cálculo formal con $z$ y $\bar{z}$]
Identificando $x = \frac{1}{2}(z + \bar{z})$ e $y = \frac{1}{2i}(z - \bar{z})$, la diferencial de una función se descompone como $\mathrm{d}f = \frac{\partial f}{\partial z} \mathrm{d}z + \frac{\partial f}{\partial \bar{z}} \mathrm{d}\bar{z}$. En este sentido formal, una función es independiente de la variable conjugada $\bar{z}$ precisamente cuando $\frac{\partial f}{\partial \bar{z}} = 0$.
\end{observacion}

\begin{proposicion}
Si $f$ es $\mathbb{C}$-diferenciable en $a \in \Omega$, entonces
\begin{equation}
\frac{\partial f}{\partial \bar{z}}(a) = 0 \quad \text{y} \quad f'(a) = \frac{\partial f}{\partial z}(a) = \frac{\partial f}{\partial x}(a) = -i\frac{\partial f}{\partial y}(a).
\end{equation}
\end{proposicion}
\begin{proof}
Calculando el límite incremental a lo largo de las direcciones real y puramente imaginaria:
\begin{enumerate}
\item Con $h \in \mathbb{R}$, $h \to 0$:
\[
f'(a) = \lim_{h \to 0} \frac{f(a+h) - f(a)}{h} = \frac{\partial f}{\partial x}(a).
\]
\item Con $ih \in i\mathbb{R}$, $h \to 0$:
\[
f'(a) = \lim_{h \to 0} \frac{f(a+ih) - f(a)}{ih} = \frac{1}{i} \frac{\partial f}{\partial y}(a) = -i \frac{\partial f}{\partial y}(a).
\]
\end{enumerate}
Igualando ambas derivadas direccionales: $\frac{\partial f}{\partial x}(a) = -i\frac{\partial f}{\partial y}(a)$. En consecuencia:
\[
\frac{\partial f}{\partial \bar{z}}(a) = \frac{1}{2} \left( \frac{\partial f}{\partial x}(a) + i \frac{\partial f}{\partial y}(a) \right) = \frac{1}{2} \left( -i\frac{\partial f}{\partial y}(a) + i\frac{\partial f}{\partial y}(a) \right) = 0.
\]
Asimismo, $\frac{\partial f}{\partial z}(a) = \frac{1}{2}\left(\frac{\partial f}{\partial x}(a) - i\frac{\partial f}{\partial y}(a)\right) = \frac{\partial f}{\partial x}(a) = f'(a)$.
\end{proof}

\begin{proposicion}[Ecuaciones de Cauchy-Riemann]
Sea $f: \Omega \to \mathbb{C}$ con $f = u + iv$, donde $u, v: \Omega \to \mathbb{R}$. Si $f$ es $\mathbb{C}$-diferenciable en $a$, entonces las derivadas parciales de primer orden de $u$ y $v$ existen en $a$ y satisfacen:
\begin{equation}
\frac{\partial u}{\partial x}(a) = \frac{\partial v}{\partial y}(a), \qquad \frac{\partial u}{\partial y}(a) = -\frac{\partial v}{\partial x}(a).
\end{equation}
\end{proposicion}
\begin{proof}
Dado que $\frac{\partial f}{\partial x} = \frac{\partial u}{\partial x} + i\frac{\partial v}{\partial x}$ y $\frac{\partial f}{\partial y} = \frac{\partial u}{\partial y} + i\frac{\partial v}{\partial y}$, la relación $\frac{\partial f}{\partial x} = -i\frac{\partial f}{\partial y}$ equivale a:
\[
\frac{\partial u}{\partial x} + i\frac{\partial v}{\partial x} = -i\left( \frac{\partial u}{\partial y} + i\frac{\partial v}{\partial y} \right) = \frac{\partial v}{\partial y} - i\frac{\partial u}{\partial y}.
\]
Igualando las partes reales e imaginarias se obtienen las identidades enunciadas.
\end{proof}

\begin{definicion}[Ecuaciones de Cauchy-Riemann]
Sea $f = u+iv$. El sistema de ecuaciones diferenciales parciales
\begin{equation}
\frac{\partial u}{\partial x} = \frac{\partial v}{\partial y}, \qquad \frac{\partial u}{\partial y} = -\frac{\partial v}{\partial x},
\end{equation}
o sus formulaciones equivalentes $\frac{\partial f}{\partial x} = -i\frac{\partial f}{\partial y}$ y $\frac{\partial f}{\partial \bar{z}} = 0$, recibe el nombre de \textbf{ecuaciones de Cauchy-Riemann}.
\end{definicion}

\begin{proposicion}[$\mathbb{C}$-linealidad de la diferencial tangente]
Sea $f = u+iv$ diferenciable en sentido real en $a \in \Omega$. Sea $\mu: \mathbb{C} \to \mathbb{R}^2$ el isomorfismo canónico $\mu(x+iy) = (x,y)$. La aplicación $\mathbb{R}$-lineal $\mathrm{d}f = \mu^{-1} \circ \mathrm{d}(u,v) \circ \mu: \mathbb{C} \to \mathbb{C}$ es $\mathbb{C}$-lineal si y sólo si $\frac{\partial f}{\partial \bar{z}}(a) = 0$. En tal caso, para todo $\zeta \in \mathbb{C}$:
\begin{equation}
\mathrm{d}f(\zeta) = f'(a)\zeta.
\end{equation}
\end{proposicion}
\begin{proof}
Una aplicación $\mathbb{R}$-lineal de $\mathbb{R}^2$ representada por la matriz $\begin{pmatrix} A & B \\ C & D \end{pmatrix}$ induce un operador $\mathbb{C}$-lineal en $\mathbb{C}$ si y sólo si conmuta con la multiplicación por $i$ ($L(i\zeta) = i L(\zeta)$), lo cual exige $A = D$ y $B = -C$. Para la diferencial de $(u,v)$, estas condiciones son exactamente las ecuaciones de Cauchy-Riemann. Cuando se cumplen,
\[
\mathrm{d}f(\xi + i\eta) = (A\xi + B\eta) + i(C\xi + D\eta) = (A - iB)(\xi + i\eta) = \left(\frac{\partial u}{\partial x} - i\frac{\partial u}{\partial y}\right)\zeta = f'(a)\zeta.
\]
\end{proof}

\begin{observacion}[Propiedades algebraicas y regla de la cadena]
De la definición de derivada compleja se deducen directamente:
\begin{enumerate}
\item Si $f, g$ son $\mathbb{C}$-diferenciables en $\Omega$ y $\lambda \in \mathbb{C}$, entonces $f+g$, $f \cdot g$ y $\lambda f$ son $\mathbb{C}$-diferenciables en $\Omega$. Si $g(a) \neq 0$, $f/g$ es $\mathbb{C}$-diferenciable en $a$.
\item \textbf{Regla de la cadena}: Si $f: U \to V$ y $g: V \to \mathbb{C}$ son $\mathbb{C}$-diferenciables con $f(U) \subset V$, entonces $g \circ f$ es $\mathbb{C}$-diferenciable sobre $U$, con $(g \circ f)'(a) = g'(f(a)) f'(a)$.
\end{enumerate}
\end{observacion}

\begin{proposicion}
Sea $f: \Omega \to \mathbb{C}$. Supongamos que las derivadas parciales $\frac{\partial f}{\partial x}$ y $\frac{\partial f}{\partial y}$ existen y son continuas sobre $\Omega$. Si además satisfacen $\frac{\partial f}{\partial x} = -i\frac{\partial f}{\partial y}$ (es decir, $\frac{\partial f}{\partial \bar{z}} = 0$) en todo punto de $\Omega$, entonces $f$ es $\mathbb{C}$-diferenciable sobre $\Omega$.
\end{proposicion}
\begin{proof}
Sean $a = \alpha + i\beta \in \Omega$ y $\zeta = \xi + i\eta \in \mathbb{C}$ suficientemente pequeño. Al ser $u, v \in C^1$, por el teorema de Taylor de primer orden en dos variables reales:
\[
u(\alpha+\xi, \beta+\eta) - u(\alpha,\beta) = \frac{\partial u}{\partial x}(a)\xi + \frac{\partial u}{\partial y}(a)\eta + \varepsilon_1(\zeta),
\]
\[
v(\alpha+\xi, \beta+\eta) - v(\alpha,\beta) = \frac{\partial v}{\partial x}(a)\xi + \frac{\partial v}{\partial y}(a)\eta + \varepsilon_2(\zeta),
\]
donde $\lim_{\zeta \to 0} \frac{\varepsilon_j(\zeta)}{|\zeta|} = 0$ ($j=1,2$). Multiplicando la segunda ecuación por $i$, sumando ambas y utilizando que $\frac{\partial f}{\partial y} = i\frac{\partial f}{\partial x}$:
\[
f(a+\zeta) - f(a) = \frac{\partial f}{\partial x}(a)\xi + \frac{\partial f}{\partial y}(a)\eta + \varepsilon(\zeta) = \frac{\partial f}{\partial x}(a)(\xi + i\eta) + \varepsilon(\zeta) = \frac{\partial f}{\partial x}(a)\zeta + \varepsilon(\zeta),
\]
con $\lim_{\zeta \to 0} \frac{\varepsilon(\zeta)}{|\zeta|} = 0$. Por consiguiente, existe $\lim_{\zeta \to 0} \frac{f(a+\zeta)-f(a)}{\zeta} = \frac{\partial f}{\partial x}(a)$, lo que establece que $f$ es $\mathbb{C}$-diferenciable.
\end{proof}

\begin{observacion}[Teorema de Looman-Menchoff]
Nótese que la hipótesis de continuidad de las derivadas parciales es en realidad superflua: basta exigir que existan en cada punto y satisfagan las ecuaciones de Cauchy-Riemann (junto con la continuidad de $f$). Este resultado profundo, debido a Looman y Menchoff, se demostrará rigurosamente en la Sección 1.6.
\end{observacion}

\begin{lema}[Lema de Abel]
Dada una sucesión $\{c_n\}_{n \ge 0}$ de números complejos, existe un número $R \in [0, +\infty]$ tal que la serie de potencias
\begin{equation}
\sum_{n=0}^\infty c_n z^n
\end{equation}
converge absolutamente para $|z| < R$ y diverge para $|z| > R$. Además, la serie converge uniformemente sobre cualquier subconjunto compacto del disco $D(0,R) = \{z \in \mathbb{C} : |z| < R\}$.
\end{lema}
\begin{proof}
Definamos $R = \sup\{r \ge 0 : \exists M_r > 0 \text{ tal que } |c_n| r^n \le M_r, \; \forall n \ge 0\}$.
\begin{itemize}
\item Si $|z| > R$, la sucesión $\{|c_n| |z|^n\}$ no es acotada; por tanto, el término general no converge a cero y la serie diverge.
\item Sea $K \subset D(0,R)$ compacto. Por compacidad, existe $\rho < R$ tal que $|z| \le \rho$ para todo $z \in K$. Escojamos $r$ tal que $\rho < r < R$. Por definición del supremo, existe $M > 0$ tal que $|c_n| r^n \le M$ para todo $n \ge 0$. Luego, para todo $z \in K$:
\[
|c_n z^n| \le |c_n| \rho^n = |c_n| r^n \left(\frac{\rho}{r}\right)^n \le M \left(\frac{\rho}{r}\right)^n.
\]
Puesto que $\rho/r < 1$, la serie geométrica mayorante converge, y por el criterio $M$ de Weierstrass la serie converge absoluta y uniformemente sobre $K$.
\end{itemize}
\end{proof}

\begin{corolario}
Toda función holomorfa sobre un conjunto abierto $\Omega$ es continua sobre $\Omega$.
\end{corolario}
\begin{proof}
Alrededor de cada punto, $f$ es el límite uniforme de polinomios (que son continuos) sobre discos cerrados. Por ende, $f$ es continua.
\end{proof}

\begin{definicion}[Radio de convergencia]
Dada la serie de potencias $\sum_{n=0}^\infty c_n (z-a)^n$, el valor $R$ dado por el Lema de Abel se denomina el \textbf{radio de convergencia} de la serie.
\end{definicion}

\begin{lema}
El radio de convergencia de la serie formalmente derivada $\sum_{n=1}^\infty n c_n z^{n-1}$ es igual al radio de convergencia de $\sum_{n=0}^\infty c_n z^n$.
\end{lema}
\begin{proof}
Sean $R'$ y $R$ los radios respectivos. Para $n \ge 1$, $|c_n z^n| \le |n c_n z^{n-1}| |z|$; luego, si $\{n |c_n| r^{n-1}\}$ es acotada, $\{|c_n| r^n\}$ también lo es, lo que implica $R' \le R$.

Recíprocamente, si $R > 0$, tomemos $0 < r < \rho < R$. Tenemos:
\[
n |c_n| r^{n-1} = \frac{1}{r} \left( n \left(\frac{r}{\rho}\right)^n \right) (|c_n| \rho^n).
\]
Como $|c_n| \rho^n \le M$ y $n(r/\rho)^n \to 0$ al ser $r/\rho < 1$, la sucesión $\{n |c_n| r^{n-1}\}$ resulta acotada, deduciéndose que $R' \ge r$. Siendo $r < R$ arbitrario, $R' \ge R$, con lo que $R' = R$.
\end{proof}

\begin{lema}
Para cualesquiera $\alpha, \beta \in \mathbb{C}$ y todo $n \ge 1$:
\begin{equation}
|(\alpha + \beta)^n - \alpha^n| \le n |\beta| (|\alpha| + |\beta|)^{n-1}.
\end{equation}
\end{lema}
\begin{proof}
Se utiliza la descomposición algebraica
\[
(\alpha + \beta)^n - \alpha^n = \beta \sum_{k=0}^{n-1} (\alpha + \beta)^{n-1-k} \alpha^k.
\]
Tomando valores absolutos y acotando cada término mediante $|(\alpha+\beta)^{n-1-k}\alpha^k| \le (|\alpha|+|\beta|)^{n-1-k}(|\alpha|+|\beta|)^k = (|\alpha|+|\beta|)^{n-1}$, la suma de $n$ términos idénticos produce la cota deseada.
\end{proof}

\begin{proposicion}
Sea $R > 0$ y supongamos que la serie de potencias $f(z) = \sum_{n=0}^\infty c_n (z-a)^n$ converge en el disco $D(a,R)$. Entonces $f$ es $\mathbb{C}$-diferenciable sobre $D(a,R)$ y su derivada se obtiene diferenciando término a término:
\begin{equation}
f'(z) = \sum_{n=1}^\infty n c_n (z-a)^{n-1}.
\end{equation}
\end{proposicion}
\begin{proof}
Fijemos $z \in D(a,R)$ y consideremos $\zeta \in \mathbb{C}$ no nulo con $|\zeta| < \frac{1}{2}(R - |z-a|)$. Definamos $\rho = \frac{1}{2}(R + |z-a|) < R$. Entonces $|z-a| < \rho$ y $|z+\zeta-a| \le |z-a|+|\zeta| < \rho$.
Calculamos:
\[
\frac{f(z+\zeta) - f(z)}{\zeta} - \sum_{n=1}^\infty n c_n (z-a)^{n-1} = \sum_{n=1}^\infty c_n \left[ \frac{(z+\zeta-a)^n - (z-a)^n}{\zeta} - n(z-a)^{n-1} \right].
\]
Por el lema anterior, para cada $n$:
\[
\left| \frac{(z+\zeta-a)^n - (z-a)^n}{\zeta} - n(z-a)^{n-1} \right| \le n \rho^{n-1} + n |z-a|^{n-1} \le 2n \rho^{n-1}.
\]
Dado $\varepsilon > 0$, dado que $\sum_{n=1}^\infty n |c_n| \rho^{n-1}$ converge, existe $N$ tal que $\sum_{n > N} 2n |c_n| \rho^{n-1} < \varepsilon/2$. Para la suma finita de los primeros $N$ términos (un polinomio), la diferenciabilidad garantiza la existencia de $\delta > 0$ tal que para $0 < |\zeta| < \delta$:
\[
\sum_{n=1}^N |c_n| \left| \frac{(z+\zeta-a)^n - (z-a)^n}{\zeta} - n(z-a)^{n-1} \right| < \frac{\varepsilon}{2}.
\]
Sumando ambas estimaciones, la diferencia total es menor que $\varepsilon$ para todo $|\zeta| < \delta$, lo que prueba que el límite existe y es la serie indicada.
\end{proof}

\begin{observacion}
La demostración exhibe además que la convergencia del cociente incremental es uniforme en $z$ sobre cualquier disco cerrado concentrado $\{|z-a| \le r\}$ con $r < R$.
\end{observacion}

\begin{corolario}
Toda función holomorfa sobre $\Omega$ es $\mathbb{C}$-diferenciable sobre $\Omega$.
\end{corolario}

\begin{corolario}
Toda función holomorfa sobre $\Omega$ es de clase $C^\infty(\Omega)$, y su derivada $f'$ es a su vez una función holomorfa sobre $\Omega$.
\end{corolario}
\begin{proof}
Por la proposición previa, $f'$ viene representada localmente por una serie de potencias convergente; por consiguiente, $f'$ es holomorfa. Reiterando el procedimiento por inducción, todas las derivadas sucesivas $f^{(k)}$ existen y son funciones holomorfas.
\end{proof}

\subsection{1.2. Integración a lo largo de curvas}

\begin{definicion}[Integral sobre un intervalo]
Sea $f: [a,b] \to \mathbb{C}$ una función continua sobre el intervalo compacto $[a,b] \subset \mathbb{R}$, con $f(t) = u(t) + iv(t)$ ($u, v$ funciones reales). Se define la integral de $f$ sobre $[a,b]$ mediante:
\begin{equation}
\int_a^b f(t) \dt = \int_a^b u(t) \dt + i \int_a^b v(t) \dt.
\end{equation}
\end{definicion}

\begin{observacion}[Propiedades operacionales y desigualdad fundamental]
La integral compleja sobre un intervalo satisface la $\mathbb{C}$-linealidad: $\int_a^b c f(t) \dt = c \int_a^b f(t) \dt$ para todo $c \in \mathbb{C}$. Además, cuando $a \le b$, rige la desigualdad fundamental:
\begin{equation}
\left| \int_a^b f(t) \dt \right| \le \int_a^b |f(t)| \dt.
\end{equation}
\textit{Demostración}: Sea $I = \int_a^b f(t) \dt$. Si $I = 0$, la desigualdad es trivial. Si $I \neq 0$, escribamos $I = |I|e^{i\theta}$ con $\theta \in \mathbb{R}$. Entonces:
\[
|I| = e^{-i\theta} \int_a^b f(t) \dt = \int_a^b \mathrm{Re}\left(e^{-i\theta} f(t)\right) \dt \le \int_a^b \left|e^{-i\theta} f(t)\right| \dt = \int_a^b |f(t)| \dt.
\]
\end{observacion}

\begin{definicion}[Integral de funciones continuas a trozos]
Si $f: [a,b] \to \mathbb{C}$ es continua a trozos (es decir, continua salvo en un número finito de puntos donde existen límites laterales finitos), fijada una partición $a = t_0 < t_1 < \dots < t_k = b$ tal que $f|_{(t_{j-1}, t_j)}$ sea continua y extensible con continuidad a $[t_{j-1}, t_j]$, se define:
\begin{equation}
\int_a^b f(t) \dt = \sum_{j=1}^k \int_{t_{j-1}}^{t_j} f(t) \dt.
\end{equation}
\end{definicion}

\begin{definicion}[Curva diferenciable a trozos]
Sea $\Omega \subset \mathbb{C}$ abierto. Una aplicación continua $\gamma: [a,b] \to \Omega$ se llama una \textbf{curva diferenciable a trozos} si existe una partición $a = t_0 < t_1 < \dots < t_k = b$ tal que cada restricción $\gamma|_{[t_{j-1}, t_j]}$ sea continuamente diferenciable (con derivadas laterales existentes en los extremos). 
La curva se dice \textbf{cerrada} (o un \textbf{lazo}) si $\gamma(a) = \gamma(b)$. La imagen $\gamma([a,b])$ se denota por $\mathrm{Im}(\gamma)$.
\end{definicion}

\begin{observacion}[Reparametrizaciones admisibles]
Dos curvas $\gamma_1: [a,b] \to \Omega$ y $\gamma_2: [c,d] \to \Omega$ se dicen equivalentes si existe un difeomorfismo creciente $\tau: [c,d] \to [a,b]$ de clase $C^1$ a trozos tal que $\gamma_2 = \gamma_1 \circ \tau$.
\end{observacion}

\begin{definicion}[Composición y curva opuesta]
\begin{enumerate}
\item Si $\gamma_1: [a,b] \to X$ y $\gamma_2: [b,c] \to X$ cumplen $\gamma_1(b) = \gamma_2(b)$, su \textbf{yuxtaposición} $\gamma_1 \cdot \gamma_2: [a,c] \to X$ es la curva que recorre sucesivamente $\gamma_1$ y $\gamma_2$.
\item La \textbf{curva opuesta} $\gamma^-: [a,b] \to X$ se define por $\gamma^-(t) = \gamma(a + b - t)$.
\end{enumerate}
\end{definicion}

\begin{definicion}[Integral a lo largo de una curva]
Sea $\gamma: [a,b] \to \Omega$ una curva diferenciable a trozos y sea $f: \Omega \to \mathbb{C}$ continua. La integral de línea de $f$ a lo largo de $\gamma$ se define como:
\begin{equation}
\int_\gamma f(z) \dz = \sum_{j=1}^k \int_{t_{j-1}}^{t_j} f(\gamma(t)) \gamma'(t) \dt.
\end{equation}
\end{definicion}

\begin{lema}[Invarianza por reparametrización]
Si $\gamma_1$ y $\gamma_2$ son curvas diferenciables a trozos equivalentes mediante una reparametrización estrictamente creciente, entonces $\int_{\gamma_1} f(z) \dz = \int_{\gamma_2} f(z) \dz$.
\end{lema}
\begin{proof}
Escribiendo $\gamma_2(s) = \gamma_1(\tau(s))$ con $\tau'(s) > 0$, el teorema de cambio de variables para integrales de Riemann da:
\[
\int_c^d f(\gamma_2(s))\gamma_2'(s)\mathrm{d}s = \int_c^d f(\gamma_1(\tau(s)))\gamma_1'(\tau(s))\tau'(s)\mathrm{d}s = \int_a^b f(\gamma_1(t))\gamma_1'(t)\dt.
\]
\end{proof}

\begin{definicion}[Longitud de una curva]
La \textbf{longitud} de una curva diferenciable a trozos $\gamma: [a,b] \to \mathbb{C}$ se define por:
\begin{equation}
L(\gamma) = \int_a^b |\gamma'(t)| \dt.
\end{equation}
\end{definicion}

\begin{lema}[Desigualdad de acotación estándar / Desigualdad ML]
Sea $f$ continua sobre $\Omega$ y $\gamma: [a,b] \to \Omega$ diferenciable a trozos. Entonces:
\begin{equation}
\left| \int_\gamma f(z) \dz \right| \le \sup_{z \in \mathrm{Im}(\gamma)} |f(z)| \cdot L(\gamma).
\end{equation}
\end{lema}
\begin{proof}
Aplicando la desigualdad fundamental de integrales:
\[
\left| \int_\gamma f(z) \dz \right| \le \int_a^b |f(\gamma(t))| |\gamma'(t)| \dt \le \left(\sup_{z \in \mathrm{Im}(\gamma)} |f(z)|\right) \int_a^b |\gamma'(t)| \dt = \sup_{z \in \mathrm{Im}(\gamma)} |f(z)| \cdot L(\gamma).
\]
\end{proof}

\begin{lema}[Propiedades de la integral curvilínea]
\begin{enumerate}
\item $\int_{\gamma_1 \cdot \gamma_2} f(z) \dz = \int_{\gamma_1} f(z) \dz + \int_{\gamma_2} f(z) \dz$.
\item $\int_{\gamma^-} f(z) \dz = - \int_\gamma f(z) \dz$.
\item La aplicación $f \mapsto \int_\gamma f(z) \dz$ es $\mathbb{C}$-lineal.
\end{enumerate}
\end{lema}
\begin{proof}
Consecuencia directa de las propiedades de la integral de Riemann y la regla de la cadena.
\end{proof}

\begin{definicion}[Rectángulo cerrado y frontera orientada]
Dados números reales $a < b$ y $c < d$, el conjunto $R = [a,b] \times [c,d] = \{z = x+iy \in \mathbb{C} : a \le x \le b, \; c \le y \le d\}$ se denomina un \textbf{rectángulo cerrado}. Su frontera $\partial R$ se considera como una curva cerrada diferenciable a trozos recorrida en sentido antihorario a través de los cuatro vértices: $v_1 = a+ic \to v_2 = b+ic \to v_3 = b+id \to v_4 = a+id \to v_1$.
\end{definicion}

\begin{proposicion}[Fórmula de Green para rectángulos]
Sea $\Omega \subset \mathbb{C}$ abierto, $f \in C^1(\Omega)$ y $R \subset \Omega$ un rectángulo cerrado. Entonces:
\begin{equation}
\int_{\partial R} f(z) \dz = 2i \iint_R \frac{\partial f}{\partial \bar{z}} \mathrm{d}x\mathrm{d}y.
\end{equation}
\end{proposicion}
\begin{proof}
Puesto que $\dz = \mathrm{d}x + i\mathrm{d}y$, escribimos $\int_{\partial R} f \dz = \int_{\partial R} f \mathrm{d}x + i f \mathrm{d}y$. Aplicando el teorema clásico de Green en el plano:
\[
\int_{\partial R} f \mathrm{d}x + if \mathrm{d}y = \iint_R \left( i\frac{\partial f}{\partial x} - \frac{\partial f}{\partial y} \right) \mathrm{d}x\mathrm{d}y.
\]
Recordando que $\frac{\partial f}{\partial \bar{z}} = \frac{1}{2}\left(\frac{\partial f}{\partial x} + i\frac{\partial f}{\partial y}\right)$, observamos que $2i \frac{\partial f}{\partial \bar{z}} = 2i \cdot \frac{1}{2}\left(\frac{\partial f}{\partial x} + i\frac{\partial f}{\partial y}\right) = i\frac{\partial f}{\partial x} - \frac{\partial f}{\partial y}$. La sustitución inmediata concluye la demostración.
\end{proof}

\begin{corolario}
Si $f \in C^1(\Omega)$ es $\mathbb{C}$-diferenciable, entonces para todo rectángulo cerrado $R \subset \Omega$:
\begin{equation}
\int_{\partial R} f(z) \dz = 0.
\end{equation}
\end{corolario}
\begin{proof}
Al ser $f$ $\mathbb{C}$-diferenciable, $\frac{\partial f}{\partial \bar{z}} = 0$ idénticamente sobre $\Omega$.
\end{proof}

\begin{lema}[Existencia de primitiva e integrales de línea]
Sea $f: \Omega \to \mathbb{C}$ continua. Si existe una función holomorfa $F: \Omega \to \mathbb{C}$ tal que $F' = f$, entonces para cualquier curva diferenciable a trozos $\gamma: [a,b] \to \Omega$:
\begin{equation}
\int_\gamma f(z) \dz = F(\gamma(b)) - F(\gamma(a)).
\end{equation}
En particular, si $\gamma$ es una curva cerrada, $\int_\gamma f(z) \dz = 0$.
\end{lema}
\begin{proof}
Por la regla de la cadena, en cada subintervalo de diferenciabilidad de $\gamma$, $\frac{\mathrm{d}}{\dt} [F(\gamma(t))] = F'(\gamma(t))\gamma'(t) = f(\gamma(t))\gamma'(t)$. Integrando e invocando el Teorema Fundamental del Cálculo en cada tramo se obtiene el resultado.
\end{proof}

\begin{corolario}
Si $P(z)$ es un polinomio complejo y $\gamma$ es cualquier curva cerrada diferenciable a trozos en $\mathbb{C}$, entonces $\int_\gamma P(z) \dz = 0$.
\end{corolario}
\begin{proof}
Todo polinomio admite como primitiva formal otro polinomio $Q(z)$ tal que $Q'(z) = P(z)$.
\end{proof}

\begin{teorema}[Teorema de Cauchy-Goursat]
Sea $\Omega \subset \mathbb{C}$ un conjunto abierto y sea $f: \Omega \to \mathbb{C}$ una función $\mathbb{C}$-diferenciable sobre $\Omega$. Entonces, para todo rectángulo cerrado $R \subset \Omega$:
\begin{equation}
\int_{\partial R} f(z) \dz = 0.
\end{equation}
\end{teorema}
\begin{proof}
Sea $A = |\int_{\partial R} f(z) \dz|$. Dividamos el rectángulo $R$ en cuatro rectángulos congruentes $R_1, R_2, R_3, R_4$ uniendo los puntos medios de lados opuestos. La orientación antihoraria en cada frontera cancela las integrales sobre los segmentos interiores comunes, por lo que:
\[
\int_{\partial R} f \dz = \sum_{j=1}^4 \int_{\partial R_j} f \dz.
\]
Por la desigualdad triangular, existe al menos un subrectángulo, digamos $R^{(1)}$, tal que:
\[
\left| \int_{\partial R^{(1)}} f \dz \right| \ge \frac{1}{4} A.
\]
Nótese que $\mathrm{diam}(R^{(1)}) = \frac{1}{2}\mathrm{diam}(R)$ y $L(\partial R^{(1)}) = \frac{1}{2}L(\partial R)$. 
Iterando este proceso indefinidamente por bisección cuádruple, se obtiene una sucesión encajada de rectángulos compactos $\{R^{(k)}\}_{k=1}^\infty$ tales que:
\begin{enumerate}
\item $R^{(k+1)} \subset R^{(k)}$,
\item $\mathrm{diam}(R^{(k)}) = 2^{-k} \mathrm{diam}(R)$,
\item $L(\partial R^{(k)}) = 2^{-k} L(\partial R)$,
\item $\left| \int_{\partial R^{(k)}} f \dz \right| \ge 4^{-k} A$.
\end{enumerate}
Por el teorema de Cantor de conjuntos compactos encajados, la intersección $\bigcap_{k=1}^\infty R^{(k)} = \{a\}$ consiste en un único punto $a \in R \subset \Omega$.

Puesto que $f$ es $\mathbb{C}$-diferenciable en $a$, definimos el residuo $\varepsilon(z)$ por:
\[
f(z) = f(a) + f'(a)(z-a) + \varepsilon(z),
\]
donde $\lim_{z \to a} \frac{\varepsilon(z)}{z-a} = 0$. Así, dado $\delta > 0$, existe $\eta > 0$ tal que $|\varepsilon(z)| \le \delta |z-a|$ siempre que $|z-a| < \eta$.

Puesto que $z \mapsto f(a) + f'(a)(z-a)$ es un polinomio de grado 1, su integral a lo largo de cualquier curva cerrada es cero. Por consiguiente:
\[
\int_{\partial R^{(k)}} f(z) \dz = \int_{\partial R^{(k)}} \varepsilon(z) \dz.
\]
Elijamos $k$ suficientemente grande tal que $\mathrm{diam}(R^{(k)}) < \eta$. Para todo $z \in R^{(k)}$, se tiene $|z-a| \le \mathrm{diam}(R^{(k)})$, de donde $|\varepsilon(z)| \le \delta \cdot \mathrm{diam}(R^{(k)})$. Aplicando la desigualdad ML:
\[
4^{-k} A \le \left| \int_{\partial R^{(k)}} \varepsilon(z) \dz \right| \le \delta \cdot \mathrm{diam}(R^{(k)}) \cdot L(\partial R^{(k)}) = \delta \cdot 2^{-k}\mathrm{diam}(R) \cdot 2^{-k}L(\partial R) = 4^{-k} \delta \cdot \mathrm{diam}(R) L(\partial R).
\]
Multiplicando por $4^k$, se deduce que $A \le \delta \cdot \mathrm{diam}(R) L(\partial R)$. Como $\delta > 0$ es arbitrario, concluimos que $A = 0$, lo que demuestra el teorema.
\end{proof}

\begin{teorema}[Cauchy-Goursat con punto excepcional]
Sea $\Omega \subset \mathbb{C}$ abierto, $a \in \Omega$, y $f: \Omega \to \mathbb{C}$ continua sobre $\Omega$ y $\mathbb{C}$-diferenciable sobre $\Omega \setminus \{a\}$. Entonces, para todo rectángulo cerrado $R \subset \Omega$:
\begin{equation}
\int_{\partial R} f(z) \dz = 0.
\end{equation}
\end{teorema}
\begin{proof}
Se distinguen tres casos según la posición de $a$ respecto a $R$:
\begin{enumerate}
\item Si $a \notin R$, $f$ es $\mathbb{C}$-diferenciable en un entorno abierto de $R$ y el Teorema 1.2.1 aplica directamente.
\item Si $a \in \partial R$, tomemos rectángulos cerrados $R_\varepsilon \subset \mathrm{Int}(R)$ cuyos vértices convergen a los de $R$ cuando $\varepsilon \to 0$. Como $a \notin R_\varepsilon$, $\int_{\partial R_\varepsilon} f \dz = 0$. Por la continuidad uniforme de $f$ sobre el compacto $R$, $\int_{\partial R_\varepsilon} f \dz \to \int_{\partial R} f \dz$, de modo que $\int_{\partial R} f \dz = 0$.
\item Si $a \in \mathrm{Int}(R)$, tracemos rectas horizontal y vertical que pasan por $a$, dividiendo $R$ en cuatro subrectángulos que tienen al punto $a$ como vértice común. Por el caso anterior, la integral sobre la frontera de cada subrectángulo se anula. Sumando las cuatro integrales, los segmentos interiores se cancelan mutuamente y se obtiene $\int_{\partial R} f \dz = 0$.
\end{enumerate}
\end{proof}

\begin{lema}
Sea $R$ un rectángulo cerrado en $\mathbb{C}$ y sea $a \in \mathrm{Int}(R)$. Entonces:
\begin{equation}
\int_{\partial R} \frac{\dz}{z-a} = 2\pi i.
\end{equation}
\end{lema}
\begin{proof}
Para cada $t \in [0,1]$, existe un único número real $\rho(t) > 0$ tal que $\gamma(t) = a + \rho(t)e^{2\pi i t} \in \partial R$. La función $t \mapsto \rho(t)$ es continua y diferenciable a trozos (salvo en los cuatro instantes correspondientes a los vértices de $R$), y $\gamma$ es una reparametrización admisible de $\partial R$. Calculando:
\[
\gamma'(t) = \rho'(t)e^{2\pi i t} + 2\pi i \rho(t)e^{2\pi i t}, \qquad \frac{\gamma'(t)}{\gamma(t)-a} = \frac{\rho'(t)}{\rho(t)} + 2\pi i.
\]
Integrando de $0$ a $1$:
\[
\int_{\partial R} \frac{\dz}{z-a} = \int_0^1 \frac{\rho'(t)}{\rho(t)}\dt + 2\pi i \int_0^1 \dt = \ln\left(\frac{\rho(1)}{\rho(0)}\right) + 2\pi i = 2\pi i,
\]
dado que $\rho(1) = \rho(0)$ por ser la curva cerrada.
\end{proof}

\begin{teorema}[Fórmula integral de Cauchy para un rectángulo]
Sea $f$ una función $\mathbb{C}$-diferenciable sobre $\Omega$, $R \subset \Omega$ un rectángulo cerrado y $a \in \mathrm{Int}(R)$. Entonces:
\begin{equation}
f(a) = \frac{1}{2\pi i} \int_{\partial R} \frac{f(z)}{z-a} \dz.
\end{equation}
\end{teorema}
\begin{proof}
Definamos la función auxiliar $g: \Omega \to \mathbb{C}$ mediante:
\[
g(z) = \begin{cases} \dfrac{f(z) - f(a)}{z-a}, & z \in \Omega \setminus \{a\}, \\ f'(a), & z = a. \end{cases}
\]
La función $g$ es continua sobre $\Omega$ y $\mathbb{C}$-diferenciable sobre $\Omega \setminus \{a\}$. Por el Teorema 1.2.2, $\int_{\partial R} g(z) \dz = 0$. Por lo tanto:
\[
0 = \int_{\partial R} \frac{f(z) - f(a)}{z-a} \dz = \int_{\partial R} \frac{f(z)}{z-a} \dz - f(a) \int_{\partial R} \frac{\dz}{z-a}.
\]
Utilizando el lema anterior, $\int_{\partial R} \frac{\dz}{z-a} = 2\pi i$, despejando se deduce la fórmula deseada.
\end{proof}

\begin{teorema}[Toda función $\mathbb{C}$-diferenciable es holomorfa]
Sea $\Omega \subset \mathbb{C}$ abierto y $f: \Omega \to \mathbb{C}$ una función $\mathbb{C}$-diferenciable sobre $\Omega$. Entonces $f$ es holomorfa sobre $\Omega$. Más precisamente, dado $a \in \Omega$, existe un disco $D(a,r) \subset \Omega$ y una sucesión $\{c_n\}_{n=0}^\infty \subset \mathbb{C}$ tales que $f(z) = \sum_{n=0}^\infty c_n (z-a)^n$ para todo $z \in D(a,r)$.
\end{teorema}
\begin{proof}
Sea $a \in \Omega$. Escojamos un rectángulo cerrado $R \subset \Omega$ tal que $a \in \mathrm{Int}(R)$, y elijamos $r > 0$ tal que $D(a,r) \subset R$. Para todo $w \in D(a,r)$, la fórmula de Cauchy da:
\[
f(w) = \frac{1}{2\pi i} \int_{\partial R} \frac{f(z)}{z-w} \dz.
\]
Escribiendo $\frac{1}{z-w} = \frac{1}{(z-a) - (w-a)} = \frac{1}{z-a} \frac{1}{1 - \frac{w-a}{z-a}}$, y notando que para $z \in \partial R$ y $w \in D(a,r)$ se cumple $\left|\frac{w-a}{z-a}\right| \le \theta < 1$, desarrollamos en serie geométrica uniformemente convergente sobre $\partial R$:
\[
\frac{1}{z-w} = \sum_{n=0}^\infty \frac{(w-a)^n}{(z-a)^{n+1}}.
\]
Intercambiando la suma y la integral:
\[
f(w) = \sum_{n=0}^\infty c_n (w-a)^n, \qquad \text{donde } c_n = \frac{1}{2\pi i} \int_{\partial R} \frac{f(z)}{(z-a)^{n+1}} \dz.
\]
Esto demuestra que $f$ es desarrollable en serie de potencias alrededor de $a$, y por tanto holomorfa.
\end{proof}

\begin{corolario}
Si $f \in C^1(\Omega)$ y $\frac{\partial f}{\partial \bar{z}} = 0$ sobre $\Omega$, entonces $f$ es holomorfa sobre $\Omega$.
\end{corolario}

\begin{corolario}
Toda función $\mathbb{C}$-diferenciable sobre $\Omega$ es indefinidamente diferenciable ($C^\infty$).
\end{corolario}

\begin{teorema}[Teorema de Morera]
Sea $\Omega \subset \mathbb{C}$ abierto y sea $f: \Omega \to \mathbb{C}$ continua. Si para todo rectángulo cerrado $R \subset \Omega$ se cumple
\begin{equation}
\int_{\partial R} f(z) \dz = 0,
\end{equation}
entonces $f$ es holomorfa sobre $\Omega$.
\end{teorema}
\begin{proof}
Basta probar el teorema localmente en un disco $D(a,r) \subset \Omega$. Sea $a = \alpha + i\beta$. Para $z = x+iy \in D(a,r)$, consideremos el camino $\gamma_z$ compuesto por el segmento horizontal de $(\alpha,\beta)$ a $(x,\beta)$ y el segmento vertical de $(x,\beta)$ a $(x,y)$. Definamos:
\[
F(z) = \int_{\gamma_z} f(w) \dw.
\]
Incrementando $z$ en la dirección vertical por $ih$ ($h \in \mathbb{R}$):
\[
\frac{F(z+ih) - F(z)}{h} = \frac{1}{h} \int_y^{y+h} f(x+it) i\dt \xrightarrow{h \to 0} i f(z),
\]
luego $\frac{\partial F}{\partial y}$ existe y vale $i f(z)$.
Por otra parte, usando la hipótesis de que la integral sobre cualquier rectángulo se anula, el camino $\gamma_z$ puede reemplazarse por el camino alternativo que va primero en vertical hasta $(\alpha,y)$ y luego en horizontal hasta $(x,y)$. Calculando el incremento horizontal con este camino se obtiene:
\[
\frac{\partial F}{\partial x}(z) = f(z).
\]
En consecuencia, $\frac{\partial F}{\partial x}$ y $\frac{\partial F}{\partial y}$ son continuas y satisfacen:
\[
\frac{\partial F}{\partial x} = f = -i(if) = -i \frac{\partial F}{\partial y},
\]
es decir, $\frac{\partial F}{\partial \bar{z}} = 0$. Por la Proposición 1.1.5, $F$ es $\mathbb{C}$-diferenciable, y por ende holomorfa con $F' = f$. Al ser la derivada de una función holomorfa, $f$ es a su vez holomorfa.
\end{proof}

\begin{proposicion}[Existencia de primitivas en abiertos conexos]
Sea $\Omega \subset \mathbb{C}$ un abierto conexo y $f: \Omega \to \mathbb{C}$ continua. Existe una función holomorfa $F$ sobre $\Omega$ tal que $F' = f$ si y sólo si
\begin{equation}
\int_\gamma f(z) \dz = 0
\end{equation}
para toda curva cerrada diferenciable a trozos $\gamma$ en $\Omega$. En tal caso, $f$ es holomorfa.
\end{proposicion}
\begin{proof}
La necesidad se demostró en el Lema 1.2.4. Para la suficiencia, fijemos un punto base $a \in \Omega$. Para cualquier $w \in \Omega$, elijamos una curva diferenciable a trozos $\lambda_w$ que une $a$ con $w$, y definamos $F(w) = \int_{\lambda_w} f \dz$. Puesto que la integral sobre curvas cerradas es nula, el valor de $F(w)$ es independiente de la curva elegida. En un disco $D(w_0,r) \subset \Omega$, $F(w) - F(w_0) = \int_{w_0}^w f \dz$. Como en la prueba del Teorema de Morera, esto implica que $F$ es $\mathbb{C}$-diferenciable y $F'(w) = f(w)$, concluyendo la prueba.
\end{proof}

\subsection{1.3. Propiedades fundamentales de las funciones holomorfas}

\begin{definicion}[El álgebra $\mathcal{H}(\Omega)$]
Para un abierto $\Omega \subset \mathbb{C}$, se denota por $\mathcal{H}(\Omega)$ el conjunto de todas las funciones holomorfas sobre $\Omega$. Con las operaciones de suma puntual, producto y multiplicación por escalares complejos, $\mathcal{H}(\Omega)$ constituye una $\mathbb{C}$-álgebra asociativa y conmutativa.
\end{definicion}

\begin{teorema}[Principio de continuación analítica]
Sea $\Omega \subset \mathbb{C}$ un conjunto abierto conexo y sea $f \in \mathcal{H}(\Omega)$. Si existe un abierto no vacío $U \subset \Omega$ tal que $f|_U \equiv 0$, entonces $f \equiv 0$ sobre todo $\Omega$.
\end{teorema}
\begin{proof}
Definamos el conjunto $E = \{z \in \Omega : f^{(n)}(z) = 0, \; \forall n \ge 0\}$.
\begin{enumerate}
\item Al ser cada derivada $f^{(n)}$ continua, $E = \bigcap_{n=0}^\infty (f^{(n)})^{-1}(\{0\})$ es la intersección de conjuntos cerrados, luego $E$ es cerrado en $\Omega$.
\item Sea $a \in E$. Como $f$ es holomorfa, en un disco $D(a,r) \subset \Omega$ se tiene el desarrollo en serie de Taylor $f(z) = \sum_{n=0}^\infty \frac{f^{(n)}(a)}{n!} (z-a)^n$. Al ser $a \in E$, $f^{(n)}(a) = 0$ para todo $n \ge 0$, de donde $f(z) = 0$ en todo $D(a,r)$. Por consiguiente, $f^{(n)}(z) = 0$ para todo $z \in D(a,r)$ y todo $n \ge 0$, demostrando que $D(a,r) \subset E$. Así, $E$ es abierto en $\Omega$.
\end{enumerate}
Si $f|_U \equiv 0$, claramente $U \subset E$, por lo que $E \neq \emptyset$. Siendo $\Omega$ conexo y $E \subset \Omega$ simultáneamente abierto, cerrado y no vacío, se concluye que $E = \Omega$, es decir, $f \equiv 0$ sobre $\Omega$.
\end{proof}

\begin{definicion}[Serie de Taylor]
Sea $\Omega \subset \mathbb{C}$ abierto, $f \in \mathcal{H}(\Omega)$ y $a \in \Omega$. La serie de potencias
\begin{equation}
\sum_{n=0}^\infty \frac{f^{(n)}(a)}{n!} (z-a)^n
\end{equation}
se llama la \textbf{serie de Taylor} de $f$ en el punto $a$.
\end{definicion}

\begin{lema}
Sea $f \in \mathcal{H}(\Omega)$ y $a \in \Omega$. La serie de Taylor de $f$ en $a$ converge a $f$ en algún entorno de $a$, y es la única serie de potencias $\sum c_n (z-a)^n$ con dicha propiedad.
\end{lema}
\begin{proof}
Si $f(z) = \sum c_n (z-a)^n$ en un entorno de $a$, derivando sucesivamente término a término $k$ veces y evaluando en $z = a$, se obtiene $f^{(k)}(a) = k! c_k$, de donde $c_k = \frac{f^{(k)}(a)}{k!}$, garantizando la unicidad de los coeficientes.
\end{proof}

\begin{teorema}[Aislamiento de los ceros]
Sea $\Omega \subset \mathbb{C}$ un abierto conexo y $f \in \mathcal{H}(\Omega)$, con $f \not\equiv 0$. Entonces el conjunto de ceros
\begin{equation}
Z_f = \{z \in \Omega : f(z) = 0\}
\end{equation}
es un conjunto discreto en $\Omega$ (es cerrado y todos sus puntos son aislados).
\end{teorema}
\begin{proof}
Por la continuidad de $f$, $Z_f = f^{-1}(\{0\})$ es cerrado en $\Omega$. Sea $a \in Z_f$. Como $f \not\equiv 0$, por el Principio de Continuación Analítica no todas las derivadas $f^{(n)}(a)$ pueden anularse. Sea $k \ge 1$ el menor entero tal que $f^{(k)}(a) \neq 0$. Desarrollando en serie de Taylor alrededor de $a$:
\[
f(z) = \sum_{n=k}^\infty c_n (z-a)^n = (z-a)^k \sum_{j=0}^\infty c_{k+j} (z-a)^j = (z-a)^k g(z),
\]
donde $g$ es holomorfa en un entorno de $a$ y $g(a) = c_k \neq 0$. Por continuidad de $g$, existe un disco $V$ centrado en $a$ donde $g(z) \neq 0$. Por lo tanto, en $V \setminus \{a\}$, $f(z) = (z-a)^k g(z) \neq 0$, lo que demuestra que $a$ es un punto aislado de $Z_f$.
\end{proof}

\begin{corolario}[Principio de identidad]
Sean $\Omega \subset \mathbb{C}$ conexo y $f, g \in \mathcal{H}(\Omega)$. Si el conjunto $\{z \in \Omega : f(z) = g(z)\}$ tiene un punto de acumulación en $\Omega$, entonces $f \equiv g$ sobre $\Omega$.
\end{corolario}
\begin{proof}
Se aplica el Teorema 1.3.2 a la función holomorfa $h = f - g$. Si $Z_h$ tiene un punto de acumulación en $\Omega$, no es discreto, lo que obliga a que $h \equiv 0$.
\end{proof}

\begin{lema}
Sea $I \subset \mathbb{R}$ un intervalo abierto y $\phi \in C^2(I, \mathbb{R})$. Si existe $t_0 \in I$ tal que $\phi(t) \le \phi(t_0)$ para todo $t \in I$, entonces $\phi''(t_0) \le 0$.
\end{lema}
\begin{proof}
Utilizando la segunda diferencia simétrica:
\[
\phi''(t_0) = \lim_{h \to 0} \frac{\phi(t_0+h) + \phi(t_0-h) - 2\phi(t_0)}{h^2} \le 0,
\]
ya que por hipótesis $\phi(t_0 \pm h) - \phi(t_0) \le 0$.
\end{proof}

\begin{proposicion}
Sea $\Omega \subset \mathbb{R}^2$ abierto y $u \in C^2(\Omega, \mathbb{R})$. Definamos $\Delta u = \frac{\partial^2 u}{\partial x^2} + \frac{\partial^2 u}{\partial y^2}$. Si $(x_0, y_0) \in \Omega$ es un máximo local de $u$, entonces:
\begin{equation}
(\Delta u)(x_0, y_0) \le 0.
\end{equation}
\end{proposicion}
\begin{proof}
Por el lema anterior aplicado a las restricciones de $u$ sobre las rectas paralelas a los ejes coordenados que pasan por $(x_0,y_0)$, se tiene $\frac{\partial^2 u}{\partial x^2}(x_0,y_0) \le 0$ y $\frac{\partial^2 u}{\partial y^2}(x_0,y_0) \le 0$. Sumando ambas desigualdades se obtiene $\Delta u(x_0,y_0) \le 0$.
\end{proof}

\begin{proposicion}[Generalización a $\mathbb{R}^n$]
Sea $\Omega \subset \mathbb{R}^n$ abierto, $u \in C^2(\Omega, \mathbb{R})$ y $x^0 \in \Omega$ tal que $u(x) \le u(x^0)$ para todo $x \in \Omega$. Entonces, para cualquier vector $(c_1, \dots, c_n) \in \mathbb{R}^n$:
\begin{equation}
\sum_{j,k=1}^n c_j c_k \frac{\partial^2 u}{\partial x_j \partial x_k}(x^0) \le 0.
\end{equation}
\end{proposicion}
\begin{proof}
Definamos $\phi(t) = u(x^0 + t c)$ para $t \in \mathbb{R}$ en un entorno de $0$. La función $\phi$ tiene un máximo en $t = 0$. Aplicando la regla de la cadena y el lema anterior: $\phi''(0) = \sum_{j,k=1}^n c_j c_k \frac{\partial^2 u}{\partial x_j \partial x_k}(x^0) \le 0$.
\end{proof}

\begin{definicion}[Laplaciano real y complejo]
Para $u \in C^2(\Omega)$, el \textbf{operador de Laplace} real es $\Delta u = \frac{\partial^2 u}{\partial x^2} + \frac{\partial^2 u}{\partial y^2}$. El \textbf{Laplaciano complejo} es:
\begin{equation}
\Delta_\mathbb{C} u = \frac{\partial^2 u}{\partial z \partial \bar{z}} = \frac{1}{4} \Delta u.
\end{equation}
\end{definicion}

\begin{teorema}[Principio del máximo: forma débil]
Sea $\Omega \subset \mathbb{C}$ abierto y $u \in C^2(\Omega, \mathbb{R})$ tal que $(\Delta u)(z) \ge 0$ para todo $z \in \Omega$. Entonces, para cualquier subconjunto abierto relativamente compacto $U \Subset \Omega$:
\begin{equation}
\sup_{z \in U} u(z) \le \sup_{z \in \partial U} u(z).
\end{equation}
\end{teorema}
\begin{proof}
Supongamos primero que $\Delta u > 0$ en $\Omega$. Como la clausura $\bar{U}$ es compacta, $u$ alcanza su máximo en algún punto $z_0 \in \bar{U}$. Si $z_0 \in U$, por la proposición anterior tendríamos $\Delta u(z_0) \le 0$, contradiciendo la hipótesis $\Delta u(z_0) > 0$. Por tanto, el máximo se alcanza necesariamente en la frontera $\partial U$.

Para el caso general $\Delta u \ge 0$, para cada $\varepsilon > 0$ definamos $u_\varepsilon(z) = u(z) + \varepsilon |z|^2$. Entonces $\Delta u_\varepsilon = \Delta u + 4\varepsilon > 0$. Por el primer caso:
\[
\sup_{z \in U} u(z) \le \sup_{z \in U} u_\varepsilon(z) \le \sup_{z \in \partial U} u_\varepsilon(z) \le \sup_{z \in \partial U} u(z) + \varepsilon \sup_{z \in \partial U} |z|^2.
\]
Haciendo tender $\varepsilon \to 0$ se obtiene la desigualdad.
\end{proof}

\begin{corolario}[Principio del máximo para funciones holomorfas: forma débil]
Sea $f \in \mathcal{H}(\Omega)$ y $U \Subset \Omega$ un abierto relativamente compacto. Entonces:
\begin{equation}
\sup_{z \in U} |f(z)| \le \sup_{z \in \partial U} |f(z)|.
\end{equation}
\end{corolario}
\begin{proof}
Definamos $u(z) = |f(z)|^2 = f(z)\overline{f(z)}$. Puesto que $\frac{\partial f}{\partial \bar{z}} = 0$ y $\frac{\partial \bar{f}}{\partial z} = \overline{\left(\frac{\partial f}{\partial \bar{z}}\right)} = 0$, calculamos:
\[
\Delta_\mathbb{C} u = \frac{\partial}{\partial z} \left( \frac{\partial (f \bar{f})}{\partial \bar{z}} \right) = \frac{\partial}{\partial z} \left( f \frac{\partial \bar{f}}{\partial \bar{z}} \right) = \frac{\partial f}{\partial z} \frac{\partial \bar{f}}{\partial \bar{z}} = |f'(z)|^2 \ge 0.
\]
Así, $\Delta u = 4\Delta_\mathbb{C} u \ge 0$. Aplicando el Teorema 1.3.3 a $u = |f|^2$ se deduce la afirmación.
\end{proof}

\begin{teorema}[Teorema de la aplicación abierta]
Sea $\Omega \subset \mathbb{C}$ un abierto conexo y sea $f \in \mathcal{H}(\Omega)$ no constante. Entonces $f$ es una aplicación abierta; es decir, para todo conjunto abierto $U \subset \Omega$, $f(U)$ es un conjunto abierto en $\mathbb{C}$.
\end{teorema}
\begin{proof}
Sea $a \in \Omega$. Basta demostrar que $f(\Omega)$ contiene un entorno de $f(a)$. Reemplazando $f$ por $f - f(a)$, podemos suponer sin pérdida de generalidad que $f(a) = 0$. Por el aislamiento de los ceros (Teorema 1.3.2), existe un disco $D = D(a,r) \subset \Omega$ tal que $f(z) \neq 0$ para todo $z \in \bar{D} \setminus \{a\}$.
Sea $\delta = \inf_{|z-a|=r} |f(z)| > 0$. 

Afirmamos que $f(D) \supset \{w \in \mathbb{C} : |w| < \frac{1}{2}\delta\}$. En efecto, supongamos que existe $w \in \mathbb{C}$ con $|w| < \delta$ tal que $w \notin f(D)$. Entonces la función $\phi(z) = \frac{1}{f(z)-w}$ es holomorfa sobre $D(a,r)$ y continua en la clausura. Por el principio del máximo (forma débil):
\[
\frac{1}{|w|} = |\phi(a)| \le \sup_{|z-a|=r} |\phi(z)| = \sup_{|z-a|=r} \frac{1}{|f(z)-w|}.
\]
Para $|z-a|=r$, $|f(z)-w| \ge |f(z)| - |w| \ge \delta - |w| > 0$. Por ende:
\[
\frac{1}{|w|} \le \frac{1}{\delta - |w|} \implies \delta - |w| \le |w| \implies |w| \ge \frac{1}{2}\delta.
\]
Esto demuestra que ningún punto del disco $\{|w| < \frac{1}{2}\delta\}$ puede quedar fuera de $f(D)$, probando que $0 \in \mathrm{Int}(f(D))$ y que $f$ es una aplicación abierta.
\end{proof}

\begin{observacion}
En realidad, un análisis más fino (basado en el principio del argumento) permite demostrar que $f(D) \supset \{w \in \mathbb{C} : |w| < \delta\}$.
\end{observacion}

\begin{teorema}[Principio del máximo para funciones holomorfas: forma fuerte]
Sea $\Omega \subset \mathbb{C}$ un conjunto abierto conexo y acotado, y sea $f \in \mathcal{H}(\Omega)$. Definamos
\begin{equation}
M = \sup_{\zeta \in \partial \Omega} \limsup_{z \to \zeta, z \in \Omega} |f(z)|.
\end{equation}
Entonces, o bien $|f(z)| < M$ para todo $z \in \Omega$, o bien $f$ es una función constante sobre $\Omega$.
\end{teorema}
\begin{proof}
Podemos asumir $M < \infty$. Probemos primero que $|f(z)| \le M$ para todo $z \in \Omega$. Dado $\varepsilon > 0$, para cada $\zeta \in \partial \Omega$ existe un entorno $U_\zeta$ tal que $|f(z)| < M + \varepsilon$ para $z \in U_\zeta \cap \Omega$. Sea $V = \bigcup_{\zeta \in \partial \Omega} U_\zeta$. Entonces $K = \Omega \setminus V$ es un subconjunto compacto de $\Omega$. Por el principio del máximo en forma débil, $\sup_{z \in K} |f(z)| \le \sup_{z \in \partial K} |f(z)| \le M+\varepsilon$. Puesto que en $V \cap \Omega$ ya se cumple $|f| < M+\varepsilon$, se tiene $|f(z)| \le M+\varepsilon$ en todo $\Omega$. Haciendo $\varepsilon \to 0$, se obtiene $|f(z)| \le M$.

Si $f$ no es constante, por el Teorema de la Aplicación Abierta $f(\Omega)$ es un abierto de $\mathbb{C}$ contenido en el disco cerrado $\overline{D(0,M)}$. Por lo tanto, $f(\Omega)$ debe estar contenido en el interior de dicho disco, es decir, $f(\Omega) \subset D(0,M)$, lo cual significa que $|f(z)| < M$ para todo $z \in \Omega$.
\end{proof}

\begin{definicion}[Primitiva de una función holomorfa]
Sea $\Omega \subset \mathbb{C}$ abierto y $f \in \mathcal{H}(\Omega)$. Una función $F \in \mathcal{H}(\Omega)$ se llama una \textbf{primitiva} de $f$ sobre $\Omega$ si $F'(z) = f(z)$ para todo $z \in \Omega$.
\end{definicion}

\begin{lema}
\begin{enumerate}
\item[(a)] Si $f \in \mathcal{H}(D(a,r))$, $f$ admite una primitiva en un entorno de $a$.
\item[(b)] Si $\Omega$ es conexo y $F_1, F_2$ son dos primitivas de $f$ sobre $\Omega$, entonces existe una constante $c \in \mathbb{C}$ tal que $F_1(z) = F_2(z) + c$ para todo $z \in \Omega$.
\end{enumerate}
\end{lema}
\begin{proof}
Para (a), si $f(z) = \sum c_n (z-a)^n$, la serie $F(z) = \sum \frac{c_n}{n+1}(z-a)^{n+1}$ tiene el mismo radio de convergencia y satisface $F' = f$. Para (b), $(F_1 - F_2)' = f - f = 0$; al ser $\Omega$ conexo, la Proposición 1.1.2 implica que $F_1 - F_2$ es constante.
\end{proof}

\begin{teorema}[Teorema de Cauchy para un disco]
Sea $D = D(a,r)$ un disco abierto. Toda función $f \in \mathcal{H}(D)$ admite una primitiva sobre $D$.
\end{teorema}
\begin{proof}
Por el lema anterior, $f$ tiene primitiva en $D(a,\rho)$ para $\rho > 0$ suficientemente pequeño. Definamos:
\[
r_0 = \sup\{\rho \in (0,r] : f \text{ admite una primitiva en } D(a,\rho)\}.
\]
Primero, $f$ tiene una primitiva $F$ en $D(a,r_0)$: para cada $\rho < r_0$, sea $F_\rho$ la única primitiva en $D(a,\rho)$ con $F_\rho(a) = 0$. Por la unicidad hasta constante, si $\rho_1 < \rho_2 < r_0$, $F_{\rho_2}|_{D(a,\rho_1)} = F_{\rho_1}$. Así, la asignación $F(z) = F_\rho(z)$ (para $|z-a| < \rho < r_0$) define una primitiva holomorfa en todo $D(a,r_0)$.

Demostremos que $r_0 = r$. Supongamos por contradicción que $r_0 < r$. Para cada punto $w$ en la circunferencia $\{|w-a|=r_0\}$, elijamos un disco $D_w$ centrado en $w$ en el cual $f$ posea una primitiva $G_w$. Ajustando una constante aditiva a $G_w$, podemos lograr que $G_w = F$ en la intersección $D_w \cap D(a,r_0)$ (que es conexa). Además, si dos discos $D_w$ y $D_{w'}$ se intersecan, su intersección común con $D(a,r_0)$ es no vacía, por lo que $G_w$ y $G_{w'}$ coinciden en $D_w \cap D_{w'}$. Por compacidad de la circunferencia, un número finito de tales discos cubre $\{|w-a|=r_0\}$, extendiendo holomórficamente la primitiva $F$ a un disco $D(a,\rho_0)$ con $\rho_0 > r_0$, lo que contradice la definición de supremo de $r_0$. Por lo tanto, $r_0 = r$.
\end{proof}

\begin{definicion}[Cocadenas, cociclos y primer grupo de cohomología]
Sea $\Omega \subset \mathbb{C}$ abierto y $\mathcal{U} = \{U_i\}_{i \in I}$ un recubrimiento abierto de $\Omega$. Sea $J = \{(i,j) \in I \times I : U_i \cap U_j \neq \emptyset\}$.
\begin{enumerate}
\item Para cualquier abierto $V$, denotamos por $\mathcal{C}(V)$ las funciones localmente constantes sobre $V$.
\item Se define el espacio de 1-cocadenas como $C^1(\mathcal{U}, \mathbb{C}) = \prod_{(i,j) \in J} \mathcal{C}(U_i \cap U_j)$.
\item El subespacio de 1-cociclos es $Z^1(\mathcal{U}, \mathbb{C}) = \{(c_{ij}) \in C^1(\mathcal{U}, \mathbb{C}) : c_{ij} + c_{jk} + c_{ki} = 0 \text{ sobre } U_i \cap U_j \cap U_k\}$.
\item Sea $C^0(\mathcal{U}, \mathbb{C}) = \prod_{i \in I} \mathcal{C}(U_i)$. El operador coborde $\delta: C^0(\mathcal{U}, \mathbb{C}) \to C^1(\mathcal{U}, \mathbb{C})$ se define por $(\delta c)_{ij} = c_i|_{U_i \cap U_j} - c_j|_{U_i \cap U_j}$. El espacio de 1-cobordes es $B^1(\mathcal{U}, \mathbb{C}) = \mathrm{Im}(\delta) \subset Z^1(\mathcal{U}, \mathbb{C})$.
\item El cociente
\begin{equation}
H^1(\mathcal{U}, \mathbb{C}) = Z^1(\mathcal{U}, \mathbb{C}) / B^1(\mathcal{U}, \mathbb{C})
\end{equation}
se llama el \textbf{primer grupo de cohomología} del recubrimiento $\mathcal{U}$ con coeficientes complejos.
\end{enumerate}
\end{definicion}

\begin{teorema}[Condición cohomológica para la existencia de primitivas]
Sea $\Omega \subset \mathbb{C}$ abierto y sea $\mathcal{U} = \{U_i\}_{i \in I}$ un recubrimiento de $\Omega$ por discos. Supongamos que $H^1(\mathcal{U}, \mathbb{C}) = 0$. Entonces toda función $f \in \mathcal{H}(\Omega)$ admite una primitiva sobre $\Omega$.
\end{teorema}
\begin{proof}
Puesto que cada $U_i$ es un disco, por el Teorema 1.3.6 $f$ tiene una primitiva $F_i \in \mathcal{H}(U_i)$. Para $(i,j) \in J$, en la intersección $U_i \cap U_j$ (que es conexa al ser intersección de discos convexos), $F_i' - F_j' = f - f = 0$, luego $c_{ij} = F_i - F_j$ es una constante. 
La familia $c = (c_{ij})$ satisface $c_{ij} + c_{jk} + c_{ki} = (F_i - F_j) + (F_j - F_k) + (F_k - F_i) = 0$, por lo que $c \in Z^1(\mathcal{U}, \mathbb{C})$.
Como $H^1(\mathcal{U}, \mathbb{C}) = 0$, existe $(c_i)_{i \in I} \in C^0(\mathcal{U}, \mathbb{C})$ tal que $c_{ij} = c_i - c_j$ en $U_i \cap U_j$.
Definamos $F$ sobre $\Omega$ mediante $F|_{U_i} = F_i - c_i$. En las intersecciones $U_i \cap U_j$:
\[
(F_i - c_i) - (F_j - c_j) = (F_i - F_j) - (c_i - c_j) = c_{ij} - c_{ij} = 0.
\]
Por lo tanto, $F$ está globalmente bien definida y es holomorfa sobre $\Omega$. Finalmente, $F'|_{U_i} = F_i' = f$, probando que $F$ es una primitiva global de $f$.
\end{proof}

\begin{teorema}
Sea $D = D(a,r)$ un disco abierto y $f \in \mathcal{H}(D)$. Para cualquier curva cerrada diferenciable a trozos $\gamma$ en $D$, se tiene:
\begin{equation}
\int_\gamma f(z) \dz = 0.
\end{equation}
\end{teorema}
\begin{proof}
Consecuencia inmediata de que $f$ admite primitiva sobre $D$ (Teorema 1.3.6) y del Lema 1.2.4.
\end{proof}

\begin{teorema}[Fórmula integral de Cauchy para un disco]
Sea $\Omega \subset \mathbb{C}$ abierto, $a \in \Omega$ y $r > 0$ tales que $\overline{D(a,r)} \subset \Omega$. Sea $\gamma(t) = a + r e^{2\pi i t}$ ($0 \le t \le 1$). Entonces, para todo $f \in \mathcal{H}(\Omega)$ y todo $w \in D(a,r)$:
\begin{equation}
f(w) = \frac{1}{2\pi i} \int_\gamma \frac{f(z)}{z-w} \dz.
\end{equation}
\end{teorema}
\begin{proof}
Fijemos $w \in D(a,r)$ y consideremos la función $g: \Omega \to \mathbb{C}$ dada por $g(z) = \frac{f(z)-f(w)}{z-w}$ para $z \neq w$ y $g(w) = f'(w)$. La función $g$ es holomorfa sobre $\Omega$. Aplicando el Teorema 1.3.8 en un disco ligeramente mayor que contiene a $\overline{D(a,r)}$:
\[
0 = \int_\gamma g(z) \dz = \int_\gamma \frac{f(z)}{z-w} \dz - f(w) \int_\gamma \frac{\dz}{z-w}.
\]
Dado que $\int_\gamma \frac{\dz}{z-w} = 2\pi i$ (por la reparametrización polar clásica), se obtiene la igualdad enunciada.
\end{proof}

\begin{teorema}[Convergencia de la serie de Taylor en el disco maximal]
Sea $f \in \mathcal{H}(D(a,r))$. Entonces la serie de Taylor
\begin{equation}
\sum_{n=0}^\infty \frac{f^{(n)}(a)}{n!} (z-a)^n
\end{equation}
converge a $f(z)$ para todo $z \in D(a,r)$.
\end{teorema}
\begin{proof}
Sea $0 < \rho < r$ y sea $w \in D(a,\rho)$. Por la fórmula de Cauchy aplicada al círculo $\gamma_\rho(t) = a + \rho e^{2\pi i t}$:
\[
f(w) = \frac{1}{2\pi i} \int_{\gamma_\rho} \frac{f(z)}{z-w} \dz.
\]
Desarrollando $\frac{1}{z-w} = \sum_{n=0}^\infty \frac{(w-a)^n}{(z-a)^{n+1}}$, cuya convergencia es uniforme sobre la trayectoria $|z-a|=\rho$ para $w$ fijo con $|w-a| < \rho$:
\[
f(w) = \sum_{n=0}^\infty \left( \frac{1}{2\pi i} \int_{\gamma_\rho} \frac{f(z)}{(z-a)^{n+1}} \dz \right) (w-a)^n.
\]
Por el Lema 1.3.1, los coeficientes coinciden necesariamente con $\frac{f^{(n)}(a)}{n!}$. Puesto que $\rho < r$ es arbitrario, la serie converge en todo $D(a,r)$.
\end{proof}

\begin{corolario}
Si $f \in \mathcal{H}(\Omega)$ y $a \in \Omega$, la serie de Taylor de $f$ en $a$ converge a $f$ en el mayor disco abierto $D(a, R)$ contenido en $\Omega$ (es decir, con $R = \mathrm{dist}(a, \partial \Omega)$).
\end{corolario}

\begin{teorema}[Desigualdades de Cauchy]
Sea $f \in \mathcal{H}(D(a,r))$, sea $0 < \rho < r$ y definamos $M(\rho) = \sup_{|z-a|=\rho} |f(z)|$. Entonces, para todo $n \ge 0$:
\begin{equation}
|f^{(n)}(a)| \le \frac{n! M(\rho)}{\rho^n}.
\end{equation}
\end{teorema}
\begin{proof}
En la prueba del teorema anterior se estableció la fórmula integral para las derivadas:
\[
f^{(n)}(a) = \frac{n!}{2\pi i} \int_{|z-a|=\rho} \frac{f(z)}{(z-a)^{n+1}} \dz.
\]
Aplicando la desigualdad ML sobre la circunferencia de longitud $2\pi \rho$:
\[
|f^{(n)}(a)| \le \frac{n!}{2\pi} \cdot \frac{M(\rho)}{\rho^{n+1}} \cdot 2\pi \rho = \frac{n! M(\rho)}{\rho^n}.
\]
\end{proof}

\begin{teorema}[Teorema de Liouville]
Toda función holomorfa sobre todo el plano complejo $\mathbb{C}$ (función entera) que sea acotada es necesariamente constante.
\end{teorema}
\begin{proof}
Sea $f \in \mathcal{H}(\mathbb{C})$ y supongamos que existe $M > 0$ tal que $|f(z)| \le M$ para todo $z \in \mathbb{C}$. Por las desigualdades de Cauchy con centro en el origen y radio $r > 0$ arbitrario, para cualquier $n \ge 1$:
\[
|c_n| = \left| \frac{f^{(n)}(0)}{n!} \right| \le \frac{M}{r^n}.
\]
Haciendo tender $r \to \infty$, se concluye que $c_n = 0$ para todo $n \ge 1$. Por lo tanto, $f(z) = c_0$ es constante.
\end{proof}

\begin{teorema}[Teorema fundamental del álgebra]
Todo polinomio $P(z) = a_n z^n + a_{n-1} z^{n-1} + \dots + a_0$ con coeficientes complejos de grado $n \ge 1$ ($a_n \neq 0$) tiene al menos una raíz en $\mathbb{C}$.
\end{teorema}
\begin{proof}
Se presentan dos demostraciones canónicas:
\begin{enumerate}
\item \textbf{Demostración clásica vía Teorema de Liouville}:
Factorizando $z^n$, para $|z| = r \ge 1$:
\[
|P(z)| \ge |a_n| r^n - |a_{n-1}| r^{n-1} - \dots - |a_0| \ge |a_n| r^n \left( 1 - \frac{|a_{n-1}|}{|a_n|r} - \dots - \frac{|a_0|}{|a_n|r^n} \right).
\]
Por lo tanto, $|P(z)| \to \infty$ cuando $|z| \to \infty$. 
Si $P$ no tuviera ceros en $\mathbb{C}$, la función racional $g(z) = 1/P(z)$ sería entera ($g \in \mathcal{H}(\mathbb{C})$). Al tender $|P(z)| \to \infty$, $g(z) \to 0$ en el infinito, lo que implica que $g$ es acotada en todo $\mathbb{C}$. Por el Teorema de Liouville, $g$ debería ser constante, lo que implicaría que $P$ es constante, contradiciendo que $n \ge 1$.
\item \textbf{Demostración topológica vía el Teorema de la Aplicación Abierta}:
Puesto que $|P(z)| \to \infty$ cuando $|z| \to \infty$, la imagen $P(\mathbb{C})$ es un subconjunto cerrado de $\mathbb{C}$. En efecto, si $w_k = P(z_k) \to w$, la sucesión $\{z_k\}$ es acotada, luego posee una subsucesión convergente $z_{k_j} \to z_0$, de donde $w = P(z_0) \in P(\mathbb{C})$. 
Por otra parte, como $P$ no es constante, por el Teorema 1.3.4 $P$ es una aplicación abierta, de modo que $P(\mathbb{C})$ es también un abierto de $\mathbb{C}$. Al ser $\mathbb{C}$ conexo, los únicos subconjuntos no vacíos simultáneamente abiertos y cerrados son $\emptyset$ y $\mathbb{C}$. En consecuencia, $P(\mathbb{C}) = \mathbb{C}$. En particular, $0 \in P(\mathbb{C})$, de modo que existe $z \in \mathbb{C}$ con $P(z) = 0$.
\end{enumerate}
\end{proof}


\subsection{1.4. Los teoremas de Weierstrass y Montel}

Esta sección aborda dos pilares del análisis complejo que contrastan vivamente con el análisis real: el límite uniforme de funciones holomorfas es holomorfo (Weierstrass), y toda familia localmente acotada es precompacta (Montel), sin imponer condiciones previas sobre las derivadas.

\begin{lema}
Sea $\Omega \subset \mathbb{C}$ abierto y $\gamma$ una curva diferenciable a trozos en $\Omega$. La función $g: \mathbb{C} \setminus \mathrm{Im}(\gamma) \to \mathbb{C}$ definida por
\begin{equation}
g(w) = \int_\gamma \frac{\dz}{z - w}
\end{equation}
es holomorfa sobre $\mathbb{C} \setminus \mathrm{Im}(\gamma)$, y su derivada es
\begin{equation}
g'(w) = \int_\gamma \frac{\dz}{(z - w)^2}.
\end{equation}
\end{lema}
\begin{proof}
Fijemos $w \notin \mathrm{Im}(\gamma)$ y consideremos $\zeta \in \mathbb{C}$ suficientemente pequeño. Calculamos:
\[
\frac{g(w+\zeta) - g(w)}{\zeta} - \int_\gamma \frac{\dz}{(z-w)^2} = \int_\gamma \left[ \frac{1}{\zeta}\left( \frac{1}{z-w-\zeta} - \frac{1}{z-w} \right) - \frac{1}{(z-w)^2} \right] \dz.
\]
Nótese que el integrando es:
\[
\frac{1}{z-w-\zeta} \frac{1}{z-w} - \frac{1}{(z-w)^2} = \frac{\zeta}{(z-w-\zeta)(z-w)^2}.
\]
Dado que $\mathrm{Im}(\gamma)$ es compacto y $w \notin \mathrm{Im}(\gamma)$, $\mathrm{dist}(w, \mathrm{Im}(\gamma)) = d > 0$. Para $|\zeta| \le d/2$, $|z-w-\zeta| \ge d/2$ y $|z-w| \ge d$. Por tanto, el integrando tiende a cero uniformemente para $z \in \mathrm{Im}(\gamma)$, de donde se obtiene la derivada deseada.
\end{proof}

\begin{teorema}[Teorema de Weierstrass / Teorema de la doble serie]
Sea $\Omega \subset \mathbb{C}$ abierto y sea $\{f_n\}_{n=1}^\infty$ una sucesión de funciones holomorfas sobre $\Omega$. Supongamos que $\{f_n\}$ converge uniformemente sobre cada subconjunto compacto de $\Omega$ hacia una función $f$. Entonces:
\begin{enumerate}
\item $f$ es holomorfa sobre $\Omega$.
\item La sucesión de derivadas $\{f_n'\}$ converge uniformemente sobre cada subconjunto compacto de $\Omega$ hacia la derivada $f'$.
\end{enumerate}
\end{teorema}
\begin{proof}
Sea $a \in \Omega$ y sea $r > 0$ tal que $\overline{D(a,r)} \subset \Omega$. Sea $\gamma$ la circunferencia $|z-a|=r$. Para cualquier $w \in D(a,r)$, la fórmula de Cauchy da:
\[
f_n(w) = \frac{1}{2\pi i} \int_\gamma \frac{f_n(z)}{z-w} \dz.
\]
Dado que $f_n \to f$ uniformemente sobre el compacto $\mathrm{Im}(\gamma)$ y para $w$ fijo en un subdisco cerrado $\{|w-a| \le \rho < r\}$ la cantidad $|z-w| \ge r - \rho > 0$ está uniformemente acotada inferiormente, podemos pasar al límite bajo el signo integral:
\[
f(w) = \lim_{n \to \infty} f_n(w) = \frac{1}{2\pi i} \int_\gamma \frac{f(z)}{z-w} \dz.
\]
Por el lema anterior, la función así definida es holomorfa en $D(a,r)$. Siendo $a \in \Omega$ arbitrario, $f \in \mathcal{H}(\Omega)$.

Para la convergencia de las derivadas, diferenciando bajo la integral:
\[
f_n'(w) = \frac{1}{2\pi i} \int_\gamma \frac{f_n(z)}{(z-w)^2} \dz, \qquad f'(w) = \frac{1}{2\pi i} \int_\gamma \frac{f(z)}{(z-w)^2} \dz.
\]
Para $|w-a| \le \rho < r$:
\[
|f_n'(w) - f'(w)| \le \frac{1}{2\pi} \frac{\sup_{|z-a|=r} |f_n(z) - f(z)|}{(r-\rho)^2} \cdot 2\pi r = \frac{r}{(r-\rho)^2} \sup_{|z-a|=r} |f_n(z) - f(z)| \xrightarrow{n \to \infty} 0.
\]
Esto prueba la convergencia uniforme de $\{f_n'\}$ sobre subdiscos cerrados. Dado cualquier compacto $K \subset \Omega$, un número finito de tales discos cubre $K$, concluyendo que $f_n' \to f'$ uniformemente sobre $K$.
\end{proof}

\begin{teorema}[Teorema de Montel]
Sea $\Omega \subset \mathbb{C}$ abierto y sea $\mathcal{F} \subset \mathcal{H}(\Omega)$ una familia de funciones holomorfas que es \textbf{localmente acotada}; es decir, para cada compacto $K \subset \Omega$, existe una constante $M_K > 0$ tal que
\begin{equation}
|f(z)| \le M_K, \qquad \forall z \in K, \; \forall f \in \mathcal{F}.
\end{equation}
Entonces, dada cualquier sucesión $\{f_p\}_{p=1}^\infty \subset \mathcal{F}$, existe una subsucesión $\{f_{p_k}\}_{k=1}^\infty$ que converge uniformemente sobre cada subconjunto compacto de $\Omega$ hacia una función holomorfa.
\end{teorema}
\begin{proof}
Fijemos $a \in \Omega$ y $r > 0$ con $\overline{D(a,r)} \subset \Omega$. Para cada $f \in \mathcal{F}$, denotemos sus coeficientes de Taylor por $c_n(f) = \frac{f^{(n)}(a)}{n!}$. Por las desigualdades de Cauchy (Teorema 1.3.11):
\[
|c_n(f)| \le M r^{-n}, \qquad \forall f \in \mathcal{F}.
\]
Para $n = 0$, la sucesión numérica $\{c_0(f_p)\}$ está acotada en $\mathbb{C}$, luego admite una subsucesión convergente $\{p_k^{(0)}\}$. Inductivamente, para cada $n \ge 1$, extraemos una subsucesión sucesiva $\{p_k^{(n)}\}$ de $\{p_k^{(n-1)}\}$ tal que $\{c_n(f_{p_k^{(n)}})\}$ converge.
Tomando la subsucesión diagonal $q_k = p_k^{(k)}$, la sucesión numérica $\{c_n(f_{q_k})\}_{k=1}^\infty$ converge para todo $n \ge 0$.

Demostremos que $\{f_{q_k}\}$ converge uniformemente en el disco $D(a, r/2)$. Sea $z \in D(a, r/2)$. Para cualquier $N \ge 1$:
\[
|f_{q_k}(z) - f_{q_l}(z)| \le \sum_{n=0}^N |c_n(f_{q_k}) - c_n(f_{q_l})| (r/2)^n + \sum_{n > N} 2M r^{-n} (r/2)^n.
\]
El resto de la serie es $\sum_{n > N} 2M 2^{-n} = 2M 2^{-N}$. Dado $\varepsilon > 0$, fijamos $N$ tal que $2M 2^{-N} < \varepsilon/2$. Fijado $N$, la convergencia de cada una de las sucesiones finitas $\{c_n(f_{q_k})\}$ ($n=0, \dots, N$) garantiza la existencia de $k_0$ tal que para $k, l \ge k_0$, la suma finita es menor que $\varepsilon/2$. Por ende, $\{f_{q_k}\}$ converge uniformemente sobre $D(a, r/2)$.

Dado que $\Omega$ posee una base numerable de abiertos, existe una cantidad numerable de tales discos $\{D_v\}_{v=1}^\infty$ cuya unión es $\Omega$. Aplicando nuevamente un proceso diagonal sobre los discos $D_v$, se obtiene una subsucesión que converge uniformemente sobre cada $D_v$, y por consiguiente sobre cualquier subconjunto compacto $K \subset \Omega$ (que queda cubierto por un número finito de discos).
\end{proof}

\begin{teorema}[Teorema de Vitali]
Sea $\Omega \subset \mathbb{C}$ un abierto conexo y sea $A \subset \Omega$ un subconjunto con al menos un punto de acumulación en $\Omega$. Sea $\{f_n\}_{n=1}^\infty \subset \mathcal{H}(\Omega)$ una sucesión localmente acotada. Si la sucesión numérica $\{f_n(a)\}$ converge para cada $a \in A$, entonces $\{f_n\}$ converge uniformemente sobre cada subconjunto compacto de $\Omega$ hacia una función holomorfa.
\end{teorema}
\begin{proof}
Razonemos por contradicción. Si $\{f_n\}$ no convergiera uniformemente sobre compactos, existirían un compacto $K \subset \Omega$, dos subsucesiones $\{f_{n_k}\}$ y $\{f_{m_k}\}$, puntos $z_k \in K$ y una constante $\delta > 0$ tales que $|f_{n_k}(z_k) - f_{m_k}(z_k)| \ge \delta$ para todo $k$.

Por el Teorema de Montel, existen subsucesiones comunes que convergen uniformemente en compactos de $\Omega$ hacia funciones límites $f, g \in \mathcal{H}(\Omega)$, respectivamente. Pasando a una subsucesión de los puntos, podemos suponer además que $z_k \to z_0 \in K$, de modo que $|f(z_0) - g(z_0)| \ge \delta$.

Por otro lado, para cada $a \in A$, $\{f_n(a)\}$ converge por hipótesis, de modo que $f(a) = \lim f_{n_k}(a) = \lim f_{m_k}(a) = g(a)$. Así, $f - g \in \mathcal{H}(\Omega)$ se anula sobre $A$. Como $A$ posee un punto de acumulación en el abierto conexo $\Omega$, por el Principio de Identidad (Corolario 1.3.1) se deduce que $f \equiv g$ sobre $\Omega$, contradiciendo que $|f(z_0) - g(z_0)| \ge \delta > 0$. Esta contradicción demuestra el teorema.
\end{proof}

\subsection{1.5. Funciones meromorfas}

\begin{lema}
Sean $0 \le r_1 < r_2 \le \infty$ y sea $\Omega = \{z \in \mathbb{C} : r_1 < |z| < r_2\}$ una corona circular. Si $f \in \mathcal{H}(\Omega)$, la integral
\begin{equation}
\int_{|z|=r} f(z) \dz
\end{equation}
es independiente de $r$ para todo $r \in (r_1, r_2)$.
\end{lema}
\begin{proof}
Escribiendo la integral en coordenadas polares con $z = re^{2\pi i t}$:
\[
I(r) = \int_0^1 f(re^{2\pi i t}) r e^{2\pi i t} 2\pi i \dt = 2\pi i \int_0^1 g(r e^{2\pi i t}) \dt,
\]
donde $g(z) = z f(z)$ es holomorfa en $\Omega$. Derivando bajo el signo integral respecto a $r$:
\[
\frac{\mathrm{d}I}{\mathrm{d}r} = 2\pi i \int_0^1 g'(re^{2\pi i t}) e^{2\pi i t} \dt = \int_0^1 \frac{\mathrm{d}}{\dt} [g(re^{2\pi i t})] \frac{1}{r} \dt = \frac{1}{r} [g(re^{2\pi i}) - g(re^0)] = 0.
\]
Por lo tanto, $I(r)$ es constante.
\end{proof}

\begin{lema}
Sea $\Omega = \{z \in \mathbb{C} : r_1 < |z| < r_2\}$ una corona, $f \in \mathcal{H}(\Omega)$ y $a \in \Omega$. Sean $\rho_1, \rho_2$ tales que $r_1 < \rho_1 < |a| < \rho_2 < r_2$. Entonces:
\begin{equation}
f(a) = \frac{1}{2\pi i} \int_{|z|=\rho_2} \frac{f(z)}{z-a} \dz - \frac{1}{2\pi i} \int_{|z|=\rho_1} \frac{f(z)}{z-a} \dz.
\end{equation}
\end{lema}
\begin{proof}
Definamos en $\Omega$ la función holomorfa $g(z) = \frac{f(z)-f(a)}{z-a}$ para $z \neq a$ y $g(a) = f'(a)$. Por el lema anterior:
\[
\int_{|z|=\rho_2} g(z) \dz = \int_{|z|=\rho_1} g(z) \dz.
\]
Expandiendo la definición de $g(z)$:
\[
\int_{|z|=\rho_2} \frac{f(z)}{z-a}\dz - f(a)\int_{|z|=\rho_2}\frac{\dz}{z-a} = \int_{|z|=\rho_1} \frac{f(z)}{z-a}\dz - f(a)\int_{|z|=\rho_1}\frac{\dz}{z-a}.
\]
Dado que $a$ está en el interior del círculo $|z|=\rho_2$, $\int_{|z|=\rho_2} \frac{\dz}{z-a} = 2\pi i$. Puesto que $a$ está en el exterior del círculo $|z|=\rho_1$, la función $z \mapsto \frac{1}{z-a}$ es holomorfa en el disco $|z| \le \rho_1$, por lo que su integral sobre $|z|=\rho_1$ es cero. Sustituyendo estos valores se obtiene la fórmula.
\end{proof}

\begin{teorema}[Desarrollo en serie de Laurent]
Sean $0 \le r_1 < r_2 \le \infty$ y sea $\Omega = \{z \in \mathbb{C} : r_1 < |z| < r_2\}$. Toda función $f \in \mathcal{H}(\Omega)$ admite una única representación en serie bilátera de potencias:
\begin{equation}
f(z) = \sum_{n=-\infty}^\infty c_n z^n,
\end{equation}
donde para cualquier $r \in (r_1, r_2)$ fijado, los coeficientes están dados por:
\begin{equation}
c_n = \frac{1}{2\pi i} \int_{|z|=r} \frac{f(z)}{z^{n+1}} \dz, \qquad n \in \mathbb{Z}.
\end{equation}
La serie converge absoluta y uniformemente sobre cualquier subconjunto compacto de $\Omega$.
\end{teorema}
\begin{proof}
Sea $w \in \Omega$ y escojamos $\rho_1, \rho_2$ tales que $r_1 < \rho_1 < |w| < \rho_2 < r_2$. Por el lema anterior, $f(w) = I_2 - I_1$, con $I_2 = \frac{1}{2\pi i}\int_{|z|=\rho_2}\frac{f(z)}{z-w}\dz$ e $I_1 = \frac{1}{2\pi i}\int_{|z|=\rho_1}\frac{f(z)}{z-w}\dz$.
\begin{enumerate}
\item Para $|z| = \rho_2 > |w|$, $|w/z| < 1$, desarrollamos en serie geométrica:
\[
\frac{1}{z-w} = \sum_{n=0}^\infty \frac{w^n}{z^{n+1}} \implies I_2 = \sum_{n=0}^\infty \left(\frac{1}{2\pi i}\int_{|z|=\rho_2} \frac{f(z)}{z^{n+1}}\dz\right) w^n = \sum_{n=0}^\infty c_n w^n.
\]
\item Para $|z| = \rho_1 < |w|$, $|z/w| < 1$, desarrollamos:
\[
-\frac{1}{z-w} = \frac{1}{w} \frac{1}{1 - z/w} = \sum_{k=0}^\infty \frac{z^k}{w^{k+1}} = \sum_{m=-\infty}^{-1} \frac{w^m}{z^{m+1}} \implies -I_1 = \sum_{m=-\infty}^{-1} c_m w^m.
\]
\end{enumerate}
Por el Lema 1.5.1, la integral que define $c_n$ es independiente de si se integra sobre $\rho_1, \rho_2$ o cualquier $r \in (r_1, r_2)$.

Por el Lema de Abel, $\sum_{n=0}^\infty c_n w^n$ converge absoluta y uniformemente para $|w| \le \rho_2 < r_2$, y $\sum_{n=1}^\infty c_{-n} w^{-n}$ converge absoluta y uniformemente para $|w| \ge \rho_1 > r_1$. Por tanto, la serie total converge sobre compactos de $\Omega$.

Para la unicidad, supongamos que $f(z) = \sum_{n=-\infty}^\infty a_n z^n$ uniformemente sobre el círculo $|z|=r$. Para cada $m \in \mathbb{Z}$:
\[
\frac{1}{2\pi i} \int_{|z|=r} \frac{f(z)}{z^{m+1}}\dz = \sum_{n=-\infty}^\infty a_n \frac{1}{2\pi i} \int_{|z|=r} z^{n-m-1}\dz = a_m,
\]
ya que la integral $\int_{|z|=r} z^k \dz$ vale $2\pi i$ si $k = -1$ y $0$ si $k \neq -1$.
\end{proof}

\begin{definicion}[Serie de Laurent en una singularidad aislada]
Sea $\Omega \subset \mathbb{C}$ abierto, $a \in \Omega$ y $f \in \mathcal{H}(\Omega \setminus \{a\})$. Si $r > 0$ es tal que el disco perforado $D^*(a,r) = \{0 < |z-a| < r\} \subset \Omega$, el desarrollo
\begin{equation}
f(z) = \sum_{n=-\infty}^\infty c_n (z-a)^n
\end{equation}
se llama el \textbf{desarrollo de Laurent} de $f$ en el punto $a$.
\end{definicion}

\begin{teorema}[Teorema de extensión de Riemann]
Sea $a \in \mathbb{C}$, $r > 0$, $D = D(a,r)$ y $D^* = D \setminus \{a\}$. Sea $f \in \mathcal{H}(D^*)$. Si
\begin{equation}
\lim_{z \to a} (z - a) f(z) = 0,
\end{equation}
entonces existe una función holomorfa $F \in \mathcal{H}(D)$ tal que $F|_{D^*} = f$. En particular, esto se cumple si $f$ es acotada en un entorno perforado de $a$.
\end{teorema}
\begin{proof}
Se presentan dos demostraciones:
\begin{enumerate}
\item \textbf{Demostración mediante la serie de Laurent}:
Asumamos sin pérdida de generalidad $a = 0$. En el disco perforado, $f(z) = \sum_{n=-\infty}^\infty c_n z^n$. Para $\rho \in (0,r)$, sea $M(\rho) = \sup_{|z|=\rho} |f(z)|$. Por hipótesis, $\lim_{\rho \to 0} \rho M(\rho) = 0$. 
Para cada $n \in \mathbb{Z}$, el coeficiente de Laurent satisface:
\[
|c_n| = \left| \frac{1}{2\pi i} \int_{|z|=\rho} \frac{f(z)}{z^{n+1}}\dz \right| \le \frac{M(\rho)}{\rho^n} = \rho^{-n-1} (\rho M(\rho)).
\]
Si $n \le -1$, $-n-1 \ge 0$, por lo que $\rho^{-n-1} (\rho M(\rho)) \to 0$ cuando $\rho \to 0$. Como $c_n$ no depende de $\rho$, deducimos que $c_n = 0$ para todo $n \le -1$. Por lo tanto, $f(z) = \sum_{n=0}^\infty c_n z^n$ se extiende holomórficamente definiendo $F(0) = c_0$.
\item \textbf{Demostración directa sin serie de Laurent}:
Definamos $g: D \to \mathbb{C}$ mediante $g(z) = z^2 f(z)$ para $z \neq 0$ y $g(0) = 0$. Calculando su derivada en el origen:
\[
g'(0) = \lim_{z \to 0} \frac{g(z) - g(0)}{z} = \lim_{z \to 0} z f(z) = 0.
\]
Así, $g$ es $\mathbb{C}$-diferenciable sobre todo $D$, luego $g \in \mathcal{H}(D)$. Puesto que $g(0) = 0$ y $g'(0) = 0$, su serie de Taylor centrada en $0$ comienza a partir de grado 2: $g(z) = \sum_{n=2}^\infty a_n z^n$. Por lo tanto, para $z \neq 0$, $f(z) = g(z)/z^2 = \sum_{n=0}^\infty a_{n+2} z^n$, que define la extensión holomorfa requerida.
\end{enumerate}
\end{proof}

\begin{definicion}[Singularidad removible]
Sea $\Omega \subset \mathbb{C}$ abierto, $a \in \Omega$ y $f \in \mathcal{H}(\Omega \setminus \{a\})$. Se dice que $a$ es una \textbf{singularidad removible} (o evitable) de $f$ si existe $F \in \mathcal{H}(\Omega)$ tal que $F|_{\Omega \setminus \{a\}} = f$. Por el teorema de Riemann, esto equivale a que $f$ sea acotada cerca de $a$.
\end{definicion}

\begin{definicion}[Función meromorfa]
Sea $\Omega \subset \mathbb{C}$ abierto y $E \subset \Omega$ un subconjunto discreto de $\Omega$. Una función holomorfa $f \in \mathcal{H}(\Omega \setminus E)$ se dice \textbf{meromorfa} sobre $\Omega$ si para cada $a \in E$, existe un disco $D$ centrado en $a$ y funciones holomorfas $g, h \in \mathcal{H}(D)$ con $h \not\equiv 0$ tales que
\begin{equation}
h \cdot f|_{D \setminus E} = g|_{D \setminus E}.
\end{equation}
\end{definicion}

\begin{lema}
Sea $f \in \mathcal{H}(D^*(a,r))$ con desarrollo de Laurent $f(z) = \sum_{n=-\infty}^\infty c_n (z-a)^n$. Entonces $f$ es meromorfa sobre $D(a,r)$ si y sólo si existe un entero $N \ge 0$ tal que $c_n = 0$ para todo $n < -N$.
\end{lema}
\begin{proof}
Si $f$ es meromorfa, localmente $h(z) f(z) = g(z)$ con $h \not\equiv 0$. Factorizando el cero de $h$ en $a$: $h(z) = (z-a)^N \psi(z)$ con $\psi(a) \neq 0$. En un disco menor donde $\psi$ no se anula, $f(z) = (z-a)^{-N} (g(z)/\psi(z))$. Puesto que $g/\psi$ es holomorfa, su serie de Taylor sólo tiene potencias positivas, de modo que $(z-a)^{-N} (g/\psi)$ tiene a lo sumo potencias de orden $\ge -N$. La unicidad de la serie de Laurent implica que $c_n = 0$ para $n < -N$.

Recíprocamente, si $c_n = 0$ para $n < -N$, $(z-a)^N f(z) = \sum_{k=0}^\infty c_{k-N} (z-a)^k$ es holomorfa sobre $D(a,r)$, lo que permite escribir $f = g/h$ con $h(z) = (z-a)^N$, demostrando que $f$ es meromorfa.
\end{proof}

\begin{teorema}[Caracterización topológica de singularidades meromorfas]
Sea $\Omega \subset \mathbb{C}$ abierto, $E \subset \Omega$ discreto y $f \in \mathcal{H}(\Omega \setminus E)$. Entonces $f$ es meromorfa sobre $\Omega$ si y sólo si para cada $a \in E$ se verifica una de las dos condiciones siguientes:
\begin{enumerate}
\item $f$ es acotada en un entorno perforado de $a$ (singularidad removible).
\item $\lim_{z \to a} |f(z)| = \infty$ (polo).
\end{enumerate}
\end{teorema}
\begin{proof}
\textbf{Suficiencia}: Si $f$ es acotada, por el teorema de Riemann se extiende holomórficamente ($f = g/1$). Si $|f(z)| \to \infty$, existe un disco $V$ centrado en $a$ donde $|f(z)| \ge 1$ para $z \in V \setminus \{a\}$. La función $h(z) = 1/f(z)$ es holomorfa y acotada en $V \setminus \{a\}$. Por el teorema de Riemann, $h$ se extiende a una función holomorfa en $V$ con $h(a) = 0$. Entonces $f = 1/h$ es el cociente de dos holomorfas, luego $f$ es meromorfa.

\textbf{Necesidad}: Si $f$ es meromorfa, cerca de $a$ se tiene $f = g/h$. Escribiendo $g(z) = (z-a)^k \phi(z)$ y $h(z) = (z-a)^l \psi(z)$ con $\phi(a) \neq 0$ y $\psi(a) \neq 0$: $f(z) = (z-a)^{k-l} (\phi(z)/\psi(z))$. Si $k \ge l$, $f$ es acotada cerca de $a$. Si $k < l$, $|f(z)| = |z-a|^{-(l-k)} |\phi(z)/\psi(z)| \to \infty$ cuando $z \to a$.
\end{proof}

\begin{definicion}[Polo y orden de un polo]
Sea $f$ meromorfa sobre $\Omega$ y $a \in \Omega$.
\begin{enumerate}
\item Se dice que $a$ es un \textbf{polo} de $f$ si $\lim_{z \to a} |f(z)| = \infty$.
\item El \textbf{orden del polo} es el menor entero positivo $N \ge 1$ tal que $\lim_{z \to a} (z-a)^N f(z)$ existe y es finito y no nulo. En la serie de Laurent $f(z) = \sum_{n=-N}^\infty c_n (z-a)^n$, el orden del polo es $N$ con $c_{-N} \neq 0$.
\end{enumerate}
\end{definicion}

\begin{definicion}[Orden de una función meromorfa en un punto]
Sea $\Omega \subset \mathbb{C}$ conexo, $f$ meromorfa en $\Omega$, $f \not\equiv 0$, y $a \in \Omega$. Sea $f(z) = \sum_{n=-\infty}^\infty c_n (z-a)^n$ su serie de Laurent en $a$. Se define el \textbf{orden} de $f$ en $a$ como:
\begin{equation}
\mathrm{ord}_a(f) = \inf\{n \in \mathbb{Z} : c_n \neq 0\}.
\end{equation}
Si $f \equiv 0$, se define $\mathrm{ord}_a(0) = +\infty$.
\end{definicion}

\begin{lema}[Propiedades algebraicas del orden]
Sea $\Omega$ conexo, $f, g$ meromorfas no nulas en $\Omega$ y $a \in \Omega$.
\begin{enumerate}
\item $a$ es un polo de $f \iff \mathrm{ord}_a(f) < 0$. El orden del polo es $-\mathrm{ord}_a(f)$.
\item $f$ es holomorfa y tiene un cero en $a \iff \mathrm{ord}_a(f) > 0$.
\item $f$ es holomorfa en $a \iff \mathrm{ord}_a(f) \ge 0$.
\item $f$ es holomorfa y $f(a) \neq 0 \iff \mathrm{ord}_a(f) = 0$.
\item $\mathrm{ord}_a(f \cdot g) = \mathrm{ord}_a(f) + \mathrm{ord}_a(g)$.
\item Para $\lambda \in \mathbb{C} \setminus \{0\}$, $\mathrm{ord}_a(\lambda f) = \mathrm{ord}_a(f)$. Además:
\begin{equation}
\mathrm{ord}_a(f + g) \ge \min(\mathrm{ord}_a(f), \mathrm{ord}_a(g)),
\end{equation}
dándose la igualdad si $\mathrm{ord}_a(f) \neq \mathrm{ord}_a(g)$.
\end{enumerate}
\end{lema}

\begin{definicion}[Ceros y polos simples]
Si $\mathrm{ord}_a(f) = 1$, $a$ es un \textbf{cero simple}. Si $\mathrm{ord}_a(f) = -1$, $a$ es un \textbf{polo simple}.
\end{definicion}

\begin{definicion}[Singularidad esencial]
Sea $f \in \mathcal{H}(D^*(a,r))$. Se dice que $a$ es una \textbf{singularidad esencial} de $f$ si en su desarrollo de Laurent $\sum_{n=-\infty}^\infty c_n (z-a)^n$ existen infinitos coeficientes $c_n \neq 0$ con $n < 0$. Esto equivale a afirmar que $f$ no es meromorfa en $D(a,r)$.
\end{definicion}

\begin{teorema}[Teorema de Casorati-Weierstrass]
Sea $f \in \mathcal{H}(D^*(a,r))$ y supongamos que $a$ es una singularidad esencial de $f$. Entonces, para cualquier entorno perforado $D^* = D^*(a,r)$, la imagen $f(D^*)$ es densa en $\mathbb{C}$.
\end{teorema}
\begin{proof}
Razonemos por el absurdo. Supongamos que $f(D^*)$ no es densa en $\mathbb{C}$. Entonces existen $c \in \mathbb{C}$ y $\delta > 0$ tales que $f(D^*) \cap D(c,\delta) = \emptyset$; es decir, $|f(z) - c| \ge \delta > 0$ para todo $z \in D^*$.
Consideremos la función auxiliar $g(z) = \frac{1}{f(z) - c}$. La función $g$ es holomorfa sobre $D^*$ y está acotada por:
\[
|g(z)| \le \frac{1}{\delta}, \qquad \forall z \in D^*.
\]
Por el Teorema de Extensión de Riemann (Teorema 1.5.2), $g$ admite una extensión holomorfa $G \in \mathcal{H}(D(a,r))$. Puesto que $g$ nunca se anula en $D^*$ ($G \cdot (f-c) = 1$), $G$ no es idénticamente nula. Por lo tanto,
\[
f(z) = c + \frac{1}{G(z)}, \qquad \forall z \in D^*.
\]
Al ser $G$ holomorfa, $1/G$ es meromorfa sobre $D(a,r)$, lo que implicaría que $f$ es meromorfa sobre $D(a,r)$. Esto contradice la hipótesis de que $a$ es una singularidad esencial. En consecuencia, $f(D^*)$ es densa en $\mathbb{C}$.
\end{proof}

\begin{observacion}[Gran Teorema de Picard]
El Teorema de Casorati-Weierstrass afirma la densidad de la imagen. Un resultado notablemente más profundo, el Gran Teorema de Picard (que se probará en el Capítulo 4), establece que en cualquier entorno de una singularidad esencial, $f$ toma todos los valores complejos con la posible excepción de a lo sumo un único punto.
\end{observacion}

\subsection{1.6. El teorema de Looman-Menchoff}

Esta sección se consagra a la demostración rigurosa del Teorema de Looman-Menchoff, el cual prescinde completamente de la hipótesis de continuidad sobre las derivadas parciales en el Teorema de Cauchy-Riemann.

\begin{definicion}[Rectángulo degenerado]
Un conjunto $R = [a,b] \times [c,d] \subset \mathbb{C}$ se llama un \textbf{rectángulo degenerado} si $a = b$ o $c = d$ (reduciéndose a un segmento o a un punto). Su frontera $\partial R$ se orienta formalmente, y si $f$ es continua, $\int_{\partial R} f(z) \dz = 0$.
\end{definicion}

\begin{lema}
Sea $R \subset \mathbb{C}$ un rectángulo cerrado y sea $f$ continua sobre $R$ y holomorfa sobre el interior $\mathrm{Int}(R)$. Entonces $\int_{\partial R} f(z) \dz = 0$.
\end{lema}
\begin{proof}
Para $\varepsilon > 0$, sea $R_\varepsilon \subset \mathrm{Int}(R)$ un rectángulo concéntrico cuyos vértices tienden a los de $R$. Por el Teorema de Cauchy-Goursat ordinario, $\int_{\partial R_\varepsilon} f \dz = 0$. Por la continuidad uniforme de $f$ sobre el compacto $R$, $\int_{\partial R_\varepsilon} f \dz \to \int_{\partial R} f \dz$, concluyendo que $\int_{\partial R} f \dz = 0$.
\end{proof}

\begin{lema}
Sea $R_2$ un rectángulo cerrado y $R_1 \subset R_2$ un rectángulo (posiblemente degenerado). Supongamos que $f$ es continua sobre $R_2 \setminus \mathrm{Int}(R_1)$ y holomorfa sobre $\mathrm{Int}(R_2) \setminus R_1$. Entonces:
\begin{equation}
\int_{\partial R_2} f(z) \dz = \int_{\partial R_1} f(z) \dz.
\end{equation}
En particular, si $R_1$ es degenerado, $\int_{\partial R_2} f(z) \dz = 0$.
\end{lema}
\begin{proof}
Se descompone la región entre $R_2$ y $R_1$ en cuatro rectángulos disjuntos por sus interiores sobre los cuales $f$ es continua y holomorfa en el interior; aplicando el lema anterior a cada uno, las fronteras intermedias se cancelan dejando $\int_{\partial R_2} f \dz - \int_{\partial R_1} f \dz = 0$.
\end{proof}

\begin{lema}
Sea $S \subset \mathbb{R}$ y sea $S_0$ el conjunto de puntos aislados de $S$. Entonces $S_0$ es finito o numerable.
\end{lema}
\begin{proof}
Sea $\{I_v\}_{v=1}^\infty$ una enumeración de los intervalos abiertos con extremos racionales. Para cada $s \in S_0$, existe un menor índice $v(s)$ tal que $I_{v(s)} \cap S = \{s\}$. La aplicación $s \mapsto v(s)$ es inyectiva de $S_0$ en $\mathbb{N}$, luego $S_0$ es a lo más numerable.
\end{proof}

\begin{lema}
Sea $I = [a,b] \subset \mathbb{R}$ y $\phi: I \to \mathbb{C}$ diferenciable en todo punto de $I$. Sea $E \subset I$ un conjunto cerrado no vacío y $M > 0$ tal que
\begin{equation}
|\phi(x) - \phi(y)| \le M |x - y|, \qquad \forall x \in E, \; \forall y \in I.
\end{equation}
Entonces:
\begin{equation}
\left| \phi(b) - \phi(a) - \int_E \phi'(t) \dt \right| \le M m_1(I \setminus E),
\end{equation}
donde $m_1$ es la medida de Lebesgue en $\mathbb{R}$.
\end{lema}
\begin{proof}
Sea $E_0 = E \cup \{a,b\}$. Definamos $\psi: I \to \mathbb{C}$ igual a $\phi$ en $E_0$, y lineal afín en cada intervalo complementario de $I \setminus E_0$. Se verifica que $|\psi(x) - \psi(y)| \le M |x-y|$ para todo $x,y \in I$. En consecuencia, $\psi$ es absolutamente continua sobre $I$, por lo que $\psi(b) - \psi(a) = \int_I \psi'(t) \dt$. Por el Lema 1.6.3, $\psi' = \phi'$ casi en todas partes sobre $E$, y $|\psi'| \le M$ en $I \setminus E$. Restando la integral sobre $E$ se obtiene:
\[
|\phi(b) - \phi(a) - \int_E \phi'(t)\dt| = \left| \int_{I \setminus E} \psi'(t)\dt \right| \le M m_1(I \setminus E).
\]
\end{proof}

\begin{lema}
Sea $\Omega \subset \mathbb{C}$ abierto y $f: \Omega \to \mathbb{C}$ continua. Sea $R = [a,b] \times [c,d] \subset \Omega$ con $A^{-1} \le \frac{d-c}{b-a} \le A$ ($A > 0$). Supongamos que existe un cerrado $E \subset \Omega$ no vacío y $M > 0$ tales que $|f(x',y) - f(x,y)| \le M|x'-x|$ y $|f(x,y') - f(x,y)| \le M|y'-y|$ siempre que $(x,y) \in E$ y los otros puntos estén en $\Omega$.
Supongamos que $\frac{\partial f}{\partial x}$ y $\frac{\partial f}{\partial y}$ existen en todo $\Omega$. Sea $R_0 \subset R$ el menor rectángulo cerrado que contiene a $E \cap R$. Entonces:
\begin{equation}
\left| \int_{\partial R_0} f(z) \dz - 2i \iint_{E \cap R} \frac{\partial f}{\partial \bar{z}} \mathrm{d}x\mathrm{d}y \right| \le 8 A M m_2(R \setminus (E \cap R)).
\end{equation}
\end{lema}
\begin{proof}
Escribiendo $\int_{\partial R_0} f \dz$ a lo largo de los cuatro lados e integrando respecto a $x$ e $y$, se aplica el Lema 1.6.4 a las secciones unidimensionales $E_x = \{y : (x,y) \in E\}$ y $E_y = \{x : (x,y) \in E\}$. Cuando $E_x \neq \emptyset$, el lema unidimensional proporciona la aproximación de la integral de la derivada parcial; cuando $E_x = \emptyset$, el tamaño de la oscilación está acotado por $4AM(d-c)$. Multiplicando por $i$ en la dirección vertical y sumando, se reconstruye $i\frac{\partial f}{\partial x} - \frac{\partial f}{\partial y} = 2i\frac{\partial f}{\partial \bar{z}}$ sobre $E \cap R$, obteniéndose la cota indicada con la medida bidimensional $m_2$.
\end{proof}

\begin{teorema}[Teorema de Looman-Menchoff]
Sea $\Omega \subset \mathbb{C}$ abierto y sea $f: \Omega \to \mathbb{C}$ continua. Supongamos que las derivadas parciales $\frac{\partial f}{\partial x}$ y $\frac{\partial f}{\partial y}$ existen en cada punto de $\Omega$ y satisfacen las ecuaciones de Cauchy-Riemann
\begin{equation}
\frac{\partial f}{\partial x} = -i \frac{\partial f}{\partial y} \qquad \left(\text{es decir, } \frac{\partial f}{\partial \bar{z}} = 0\right).
\end{equation}
Entonces $f$ es holomorfa sobre $\Omega$.
\end{teorema}
\begin{proof}
Sea $\Omega' = \{z \in \Omega : f \text{ es holomorfa en un entorno de } z\}$ y sea $E = \Omega \setminus \Omega'$. El conjunto $E$ es cerrado en $\Omega$. El teorema afirma que $E = \emptyset$.
Supongamos por contradicción que $E \neq \emptyset$.

Para cada entero $k \ge 1$, definamos los conjuntos:
\[
Q_k^{(1)} = \left\{ (x,y) \in \Omega : \left| \frac{f(x',y) - f(x,y)}{x'-x} \right| \le k, \; \forall 0 < |x'-x| \le \frac{1}{k} \right\},
\]
\[
Q_k^{(2)} = \left\{ (x,y) \in \Omega : \left| \frac{f(x,y') - f(x,y)}{y'-y} \right| \le k, \; \forall 0 < |y'-y| \le \frac{1}{k} \right\},
\]
y pongamos $Q_k = Q_k^{(1)} \cap Q_k^{(2)}$. Por la continuidad de $f$, cada $Q_k$ es cerrado en $\Omega$. Como las derivadas parciales existen en todo punto, $\Omega = \bigcup_{k=1}^\infty Q_k$, y por tanto $E = \bigcup_{k=1}^\infty (Q_k \cap E)$.
Por el Teorema de Categorías de Baire aplicado al espacio métrico localmente compacto $E$, existe un índice $k_0$ tal que $Q_{k_0} \cap E$ tiene interior no vacío relativo a $E$. Es decir, existen un abierto relativamente compacto $W \subset \Omega$ y una constante $M > 0$ tales que $E \cap W \neq \emptyset$ y las hipótesis de acotación lipschitziana del Lema 1.6.5 se cumplen sobre $E \cap W$.

Para demostrar que $f$ es holomorfa en $W$ (lo que contradiría que $E \cap W \neq \emptyset$), por el Teorema de Morera basta probar que $\int_{\partial R} f(z) \dz = 0$ para todo rectángulo cerrado $R \subset W$.
Fijemos $\varepsilon > 0$. Como $E$ es medible cerrado, existe un abierto $U \supset E$ tal que $m_2(U \setminus E) < \varepsilon$. Dividamos $R$ en $4^N$ subrectángulos congruentes $R_v$ ($v = 1, \dots, 4^N$) mediante bisección sucesiva. Si $N$ es suficientemente grande, cada $R_v$ que interseca a $E$ está enteramente contenido en $U$. 
Para los rectángulos $R_v$ que no intersecan a $E$, $f$ es holomorfa en ellos y por Cauchy-Goursat $\int_{\partial R_v} f \dz = 0$. 
Para cada $R_v$ que interseca a $E$, sea $R_v^{(0)}$ el menor rectángulo que contiene a $E \cap R_v$. Por el Lema 1.6.2, $\int_{\partial R_v} f \dz = \int_{\partial R_v^{(0)}} f \dz$. Aplicando el Lema 1.6.5 a $R_v^{(0)}$ y teniendo presente que $\frac{\partial f}{\partial \bar{z}} = 0$ idénticamente:
\[
\left| \int_{\partial R_v} f(z) \dz \right| \le 8AM m_2(R_v \setminus (E \cap R_v)).
\]
Sumando sobre todos los subrectángulos que tocan a $E$:
\[
\left| \int_{\partial R} f(z) \dz \right| \le \sum_{R_v \cap E \neq \emptyset} 8AM m_2(R_v \setminus (E \cap R_v)) \le 8AM m_2(U \setminus E) < 8AM \varepsilon.
\]
Siendo $\varepsilon > 0$ arbitrario, deducimos que $\int_{\partial R} f(z) \dz = 0$. Por el Teorema de Morera, $f$ es holomorfa en $W$, lo que contradice la definición de $E$. Esta contradicción demuestra que $E = \emptyset$, completando la prueba.
\end{proof}

\subsection{Notas y comentarios sobre el Capítulo 1}

\begin{observacion}[Perspectiva histórica y metodológica]
El material presentado en este capítulo unifica las tres grandes corrientes fundacionales del análisis complejo:
\begin{enumerate}
\item \textbf{El enfoque de Cauchy (1825)}: Partió de la $\mathbb{C}$-diferenciabilidad asumiendo continuidad en las derivadas y basando la teoría en las integrales curvilíneas y el teorema del residuo.
\item \textbf{El enfoque de Weierstrass}: Fundamentó rigurosamente la teoría a partir de las series de potencias locales, proporcionando un marco algebraico potente que se generaliza con naturalidad a varias variables complejas.
\item \textbf{El enfoque de Riemann}: Priorizó las ecuaciones diferenciales de Cauchy-Riemann y las propiedades geométricas y topológicas de las transformaciones conformes y superficies de Riemann.
\end{enumerate}
El Teorema de Looman-Menchoff representa el punto culminante de la economía de hipótesis en la definición de holomorfía, demostrando que la mera diferenciabilidad real direccional junto a las ecuaciones de Cauchy-Riemann garantizan la analiticidad sin requerir continuidad previa de las derivadas parciales.
\end{observacion}

\newpage

\subsection{2.1. Espacios recubridores y levantamiento de curvas}

\begin{definicion}[Variedad topológica]
Una variedad topológica de dimensión $n$ es un espacio topológico $X$ que es de Hausdorff, segundo numerable y localmente homeomorfo a $\mathbb{R}^n$. La condición de ser localmente homeomorfo a $\mathbb{R}^n$ significa que para cada punto $x \in X$, existe un entorno abierto $U \subset X$ que contiene a $x$ y un homeomorfismo $\varphi: U \to V$, donde $V$ es un subconjunto abierto de $\mathbb{R}^n$. El par $(U, \varphi)$ recibe el nombre de carta local o sistema de coordenadas locales alrededor de $x$.
\end{definicion}

\begin{observacion}[Variedades de dimensión 2]
En el contexto del análisis complejo, nuestro interés principal recae sobre las variedades de dimensión 2 (es decir, $n=2$). En este caso, al identificar de manera natural el espacio euclidiano $\mathbb{R}^2$ con el plano complejo $\mathbb{C}$ mediante la correspondencia $(x,y) \mapsto x + iy$, las cartas locales toman la forma de homeomorfismos $\varphi: U \to V$, donde $V$ es un conjunto abierto en $\mathbb{C}$. Esta identificación es fundamental porque dota a la variedad subyacente de la estructura topológica necesaria sobre la cual se podrá posteriormente definir una estructura analítica compleja (superficie de Riemann).
\end{observacion}

\begin{definicion}[Homeomorfismo local y levantamiento]
Sean $X$ y $X'$ espacios topológicos. Una aplicación continua $p: X' \to X$ se denomina un homeomorfismo local si, para todo punto $x' \in X'$, existe un entorno abierto $U'$ de $x'$ en $X'$ tal que la imagen $p(U')$ es un conjunto abierto en $X$ y la restricción $p|_{U'}: U' \to p(U')$ es un homeomorfismo.

Dada una aplicación continua $f: Y \to X$ desde un espacio topológico $Y$, un levantamiento (o elevación) de $f$ con respecto a $p$ es una aplicación continua $f': Y \to X'$ tal que la composición de $f'$ con $p$ recupera $f$, esto es:
\begin{equation}
p \circ f' = f.
\end{equation}
En términos de un diagrama conmutativo, el triángulo formado por los espacios $Y$, $X'$ y $X$ con las flechas $f': Y \to X'$, $p: X' \to X$ y $f: Y \to X$ conmuta de manera exacta.
\end{definicion}

\begin{observacion}[No existencia y no unicidad de levantamientos]
Es imperativo notar que, sin hipótesis estructurales adicionales sobre los espacios involucrados o sobre la aplicación $p$, un levantamiento no necesariamente existe, y en caso de existir, generalmente no es único. Desarrollamos esto a través de dos ejemplos concretos, fundamentando rigurosamente cada afirmación:

\textbf{1. No existencia:} Consideremos la aplicación continua $p: \mathbb{R} \to S^1$ dada por la función exponencial $p(t) = e^{it}$. Aquí $\mathbb{R}$ está dotado de su topología euclidiana habitual y el círculo $S^1 \subset \mathbb{C}$ de la topología inducida. Es bien sabido que $p$ es un homeomorfismo local (de hecho, es el recubrimiento universal de $S^1$). Sea $f: S^1 \to S^1$ la aplicación identidad, es decir, $f(z) = z$. Supongamos, por reducción al absurdo, que existe un levantamiento $f': S^1 \to \mathbb{R}$ tal que $p \circ f' = f$. Esto implicaría que $p(f'(z)) = z$ para todo $z \in S^1$. Notemos que $f'$ es una aplicación continua desde $S^1$ hacia $\mathbb{R}$. Puesto que $f'$ es inyectiva (ya que $p \circ f'$ es la identidad, una función inyectiva), tendríamos un homeomorfismo continuo e inyectivo de $S^1$ en su imagen en $\mathbb{R}$. Como $S^1$ es compacto y conexo, su imagen $f'(S^1)$ debe ser un intervalo cerrado y acotado $[a, b] \subset \mathbb{R}$. Sin embargo, quitar un punto del interior de un intervalo conexo lo desconecta, mientras que quitar un punto de $S^1$ no desconecta al círculo. Esto constituye una contradicción topológica fundamental (el teorema de la curva de Jordan también impide sumergir el círculo en la recta real). De manera más algebraica, el grupo fundamental $\pi_1(S^1)$ es isomorfo a $\mathbb{Z}$, mientras que $\pi_1(\mathbb{R})$ es trivial. Si $f'$ existiera, induciría un homomorfismo en grupos fundamentales donde la identidad de $\mathbb{Z}$ se factorizaría a través del grupo trivial, lo cual es imposible. Por tanto, el levantamiento $f'$ de la identidad no puede existir.

\textbf{2. No unicidad:} Utilizando el mismo homeomorfismo local $p: \mathbb{R} \to S^1$, $p(t) = e^{it}$. Sea $Y$ un espacio topológico arbitrario y tomemos la aplicación constante $f: Y \to S^1$ dada por $f(y) = 1$ para todo $y \in Y$. Un levantamiento $f'$ de $f$ debe satisfacer $p(f'(y)) = e^{i f'(y)} = 1$. La preimagen de $1$ bajo $p$ es el conjunto discreto $\{ 2k\pi : k \in \mathbb{Z} \}$. Si elegimos la función constante $f'_0(y) = 0$, claramente es un levantamiento porque $e^{i\cdot 0} = 1$. No obstante, para cualquier número entero $k \neq 0$, la función constante $f'_k(y) = 2k\pi$ también es una aplicación continua de $Y$ a $\mathbb{R}$ y cumple que $e^{i (2k\pi)} = 1$. De este modo, existen infinitos levantamientos distintos de la misma función $f$, demostrando categóricamente que la mera noción de homeomorfismo local es insuficiente para garantizar la unicidad del levantamiento sin fijar un punto base.
\end{observacion}

\begin{lema}[Unicidad de levantamientos]
Sean $Y$ un espacio topológico conexo, y $X'$ un espacio de Hausdorff. Sea $p: X' \to X$ un homeomorfismo local. Supongamos que $f: Y \to X$ es una aplicación continua, y sean $f_1, f_2: Y \to X'$ dos levantamientos de $f$. Si existe al menos un punto $y_0 \in Y$ para el cual coinciden, es decir, $f_1(y_0) = f_2(y_0)$, entonces $f_1(y) = f_2(y)$ para todo $y \in Y$ (en otras palabras, $f_1 \equiv f_2$).
\end{lema}
\begin{proof}
Para probar la identidad global de las funciones, utilizaremos un argumento topológico clásico de tipo "abierto y cerrado". Definamos el conjunto de coincidencia:
\begin{equation}
E = \{ y \in Y : f_1(y) = f_2(y) \}.
\end{equation}
Nuestra meta es demostrar que $E$ es no vacío, cerrado y abierto en $Y$. Como $Y$ es por hipótesis un espacio conexo, el único subconjunto no vacío que es simultáneamente abierto y cerrado es el espacio total $Y$, de lo que se deducirá que $E = Y$ y por ende $f_1 \equiv f_2$.

\textbf{Paso 1: $E$ es no vacío.}
Por la hipótesis del lema, sabemos que $f_1(y_0) = f_2(y_0)$. Luego, el punto $y_0 \in E$, por lo que $E \neq \emptyset$.

\textbf{Paso 2: $E$ es cerrado en $Y$.}
Dado que el codominio $X'$ es un espacio de Hausdorff, la diagonal $\Delta = \{ (x', x') \in X' \times X' \}$ es un conjunto cerrado en el espacio producto $X' \times X'$ dotado de la topología producto. Consideremos la aplicación producto $F: Y \to X' \times X'$ definida mediante $F(y) = (f_1(y), f_2(y))$. Dado que tanto $f_1$ como $f_2$ son aplicaciones continuas, por la propiedad universal de la topología producto, la aplicación conjunta $F$ es continua. Observamos que el conjunto $E$ puede ser rescrito exactamente como la preimagen de la diagonal bajo $F$, es decir, $E = F^{-1}(\Delta)$. Como la preimagen de un conjunto cerrado bajo una aplicación continua es un conjunto cerrado, se concluye de forma inmediata y rigurosa que $E$ es cerrado en $Y$.

\textbf{Paso 3: $E$ es abierto en $Y$.}
Sea $y \in E$ un punto arbitrario. Por la definición de $E$, denotemos $x' = f_1(y) = f_2(y) \in X'$. Sabemos por hipótesis que la aplicación $p: X' \to X$ es un homeomorfismo local. En consecuencia, existe un entorno abierto $U'$ de $x'$ en $X'$ tal que la restricción $p|_{U'}: U' \to p(U')$ es un homeomorfismo, y en particular, es una función inyectiva. 
Como $f_1$ y $f_2$ son funciones continuas, las preimágenes de conjuntos abiertos son conjuntos abiertos. Por tanto, existen entornos abiertos $V_1$ y $V_2$ del punto $y$ en $Y$ tales que $f_1(V_1) \subset U'$ y $f_2(V_2) \subset U'$. Definamos la intersección $V = V_1 \cap V_2$. Dado que es la intersección de dos conjuntos abiertos que contienen a $y$, el conjunto $V$ es un entorno abierto de $y$ en $Y$.
Para cualquier punto $v \in V$, tenemos que $f_1(v) \in U'$ y simultáneamente $f_2(v) \in U'$. Ya que $f_1$ y $f_2$ son ambos levantamientos de la misma función $f$, deben satisfacer $p(f_1(v)) = f(v)$ y $p(f_2(v)) = f(v)$. Esto implica directamente que $p(f_1(v)) = p(f_2(v))$. 
Ahora bien, puesto que ambos valores $f_1(v)$ y $f_2(v)$ residen en el dominio de inyectividad $U'$ de $p$, la igualdad de sus imágenes bajo $p$ fuerza ineludiblemente a que sus argumentos sean idénticos, es decir, $f_1(v) = f_2(v)$. Así, hemos probado que todo punto $v \in V$ también pertenece a $E$, lo que se traduce en la inclusión $V \subset E$. Puesto que para todo $y \in E$ hemos hallado un entorno abierto $V$ contenido completamente en $E$, concluimos rigurosamente que $E$ es un conjunto abierto en $Y$.

\textbf{Conclusión:}
Habiendo establecido que $E \subset Y$ es no vacío, cerrado y abierto, la conexidad del espacio topológico $Y$ obliga a que $E = Y$. Esto demuestra de manera definitiva que $f_1(y) = f_2(y)$ para todo $y \in Y$, probando la unicidad del levantamiento cuando se especifica una condición inicial de coincidencia.
\end{proof}

\begin{definicion}[Aplicación de recubrimiento y conjuntos bien cubiertos]
Una aplicación continua y sobreyectiva $p: X' \to X$ entre dos espacios topológicos se denomina una aplicación de recubrimiento (o simplemente recubrimiento) si satisface la siguiente propiedad estructural: para cada punto $x \in X$, existe un entorno abierto $U$ de $x$ en $X$ (al cual llamaremos un entorno bien cubierto) tal que la preimagen completa $p^{-1}(U)$ se descompone topológicamente como una unión disjunta de una familia de conjuntos abiertos $\{U'_j\}_{j \in J}$ en $X'$, esto es:
\begin{equation}
p^{-1}(U) = \bigsqcup_{j \in J} U'_j,
\end{equation}
con la propiedad adicional de que, para cada índice $j \in J$, la restricción de $p$ a $U'_j$, denotada $p|_{U'_j} : U'_j \to U$, es un homeomorfismo exhaustivo sobre $U$. Los subconjuntos abiertos $U'_j \subset X'$ se denominan convencionalmente las "láminas" o "hojas" del recubrimiento sobre el conjunto $U$.
\end{definicion}

\begin{observacion}[Hojas como componentes conexas]
Es crucial destacar que si se elige el entorno abierto $U$ de forma que sea conexo, las láminas $U'_j$ coinciden exacta e inequívocamente con las componentes conexas de la preimagen $p^{-1}(U)$ dotada de su topología de subespacio. 
La demostración de este hecho es la siguiente: Como cada restricción $p|_{U'_j}: U'_j \to U$ es un homeomorfismo, y la propiedad de ser un espacio conexo es un invariante topológico, la conexidad de $U$ implica irremediablemente la conexidad de cada conjunto $U'_j$. Así, hemos particionado $p^{-1}(U)$ en la unión disjunta de los conjuntos $U'_j$, los cuales son simultáneamente abiertos (por la definición de recubrimiento) y cerrados dentro de $p^{-1}(U)$ (el complemento de cualquier $U'_k$ es la unión $\bigcup_{j \neq k} U'_j$, que es una unión de abiertos y por tanto abierta, haciendo de $U'_k$ su complemento cerrado). Puesto que los $U'_j$ son conjuntos conexos disjuntos, abiertos y cerrados en $p^{-1}(U)$, conforman exactamente el conjunto maximal de subespacios conexos, que es la definición de las componentes conexas.
\end{observacion}

\begin{observacion}[Propiedades elementales de los recubrimientos]
Del análisis profundo de la definición se desprenden directamente las siguientes propiedades esenciales:
\begin{enumerate}
    \item \textbf{Toda aplicación de recubrimiento es un homeomorfismo local, una aplicación continua, abierta y sobreyectiva.} 
    \textit{Prueba:} La continuidad es parte de la definición. Para verificar que es un homeomorfismo local, nótese que cualquier punto $x' \in X'$ pertenece a $p^{-1}(U)$ para algún entorno bien cubierto $U$ de su imagen $x = p(x')$. Así, $x'$ reside en alguna lámina única $U'_j$. El abierto $U'_j$ sirve como el entorno requerido para la definición de homeomorfismo local ya que $p|_{U'_j}: U'_j \to U$ es un homeomorfismo.
    La aplicación es abierta: si $W \subset X'$ es cualquier conjunto abierto, su imagen $p(W)$ es la unión sobre todos los puntos $x' \in W$ de las imágenes locales, lo cual, por ser homeomorfismos locales, traza un conjunto abierto en la base, demostrando que $p(W)$ es abierto.
    La sobreyectividad está dada explícitamente en la definición (para cada $x \in X$ se toma un entorno $U$ bien cubierto, y como $\bigcup U'_j \neq \emptyset$ para recubrir $U$, se sigue que $x$ tiene preimagen).
    \item \textbf{Todo subconjunto abierto de un entorno bien cubierto también es un entorno bien cubierto.}
    \textit{Prueba:} Sea $U$ un conjunto bien cubierto, con preimagen $p^{-1}(U) = \bigsqcup_{j \in J} U'_j$, donde cada $p|_{U'_j}: U'_j \to U$ es un homeomorfismo. Sea $V \subset U$ un subconjunto abierto. La preimagen de $V$ bajo $p$ es la restricción de la preimagen de $U$, es decir:
    \begin{equation*}
    p^{-1}(V) = p^{-1}(V) \cap p^{-1}(U) = p^{-1}(V) \cap \left( \bigsqcup_{j \in J} U'_j \right) = \bigsqcup_{j \in J} \left( p^{-1}(V) \cap U'_j \right).
    \end{equation*}
    Definamos $V'_j = p^{-1}(V) \cap U'_j$. Dado que $p|_{U'_j}$ es un homeomorfismo entre $U'_j$ y $U$, y $V$ es abierto en $U$, la preimagen $V'_j = (p|_{U'_j})^{-1}(V)$ es un subconjunto abierto de $U'_j$, y por ende, abierto en la topología de $X'$. Además, la restricción de un homeomorfismo a un subconjunto abierto es de nuevo un homeomorfismo sobre su imagen. Por lo tanto, $p|_{V'_j} : V'_j \to V$ es un homeomorfismo para todo $j \in J$. Al ser los $V'_j$ subconjuntos de los conjuntos disjuntos $U'_j$, también son mutuamente disjuntos. Esto satisface rigurosamente la definición de que $V$ está bien cubierto por las láminas $V'_j$.
\end{enumerate}
\end{observacion}

\begin{proposicion}[Levantamiento de curvas a lo largo de un recubrimiento]
Sea $p: X' \to X$ una aplicación de recubrimiento. Sea $\gamma: [0, 1] \to X$ una curva continua en el espacio base con punto inicial $\gamma(0) = a$. Dado cualquier punto $a' \in X'$ que yazga en la fibra sobre $a$, esto es, $p(a') = a$, existe un \textbf{único} levantamiento continuo $\gamma': [0, 1] \to X'$ de la curva $\gamma$ tal que satisface la condición inicial $\gamma'(0) = a'$.
\end{proposicion}
\begin{proof}
La demostración se estructura en tres etapas profundamente analíticas: partición topológica del dominio, construcción inductiva usando el lema del pegado y afirmación de unicidad.

\textbf{Paso 1: Partición del intervalo usando el Lema del número de Lebesgue.}
Para todo punto $x \in \gamma([0,1]) \subset X$, la definición de recubrimiento nos garantiza la existencia de un entorno abierto bien cubierto $U_x \subset X$. La colección de conjuntos $\{ U_x \}_{x \in \gamma([0,1])}$ constituye un recubrimiento abierto de la imagen $\gamma([0,1])$. Como $\gamma$ es continua, las preimágenes $\mathcal{V} = \{ \gamma^{-1}(U_x) \}_{x \in \gamma([0,1])}$ conforman un recubrimiento por abiertos del intervalo cerrado unitario $[0,1]$.
Recordamos aquí el \textit{Lema del número de Lebesgue}: en un espacio métrico compacto, todo recubrimiento abierto admite un número real $\delta > 0$ (el número de Lebesgue) tal que cualquier subconjunto del espacio con diámetro estrictamente menor a $\delta$ está completamente contenido en al menos un elemento del recubrimiento. Puesto que $[0,1]$ con la métrica euclidiana usual es un espacio métrico compacto (por el Teorema de Heine-Borel), el lema es aplicable. Sea $\delta > 0$ el número de Lebesgue asociado al recubrimiento $\mathcal{V}$.
Elegimos un entero positivo $r$ suficientemente grande de manera que $1/r < \delta$. Particionamos el intervalo $[0,1]$ usando los nodos equi-espaciados:
\begin{equation}
0 = t_0 < t_1 < t_2 < \dots < t_r < t_{r+1} = 1,
\end{equation}
donde $t_k = k/(r+1)$. Cada subintervalo cerrado $[t_v, t_{v+1}]$ tiene longitud $1/(r+1) < \delta$. Por la propiedad del número de Lebesgue, para cada índice $v \in \{0, 1, \dots, r\}$, el intervalo completo $[t_v, t_{v+1}]$ se mapea a través de $\gamma$ dentro de un único conjunto bien cubierto $U_{x_v}$, es decir, $\gamma([t_v, t_{v+1}]) \subset U_{x_v}$.

\textbf{Paso 2: Construcción inductiva del levantamiento y compatibilidad.}
La construcción de $\gamma'$ se realizará de manera secuencial intervalo por intervalo.
\textit{Caso Base ($v=0$):} Consideramos el primer subintervalo $[t_0, t_1] = [0, t_1]$. Tenemos que $\gamma([0, t_1]) \subset U_{x_0}$. Como $U_{x_0}$ está bien cubierto, su preimagen bajo $p$ es una unión disjunta de láminas abiertas homeomorfas $\bigsqcup U'_{0,j}$. El punto inicial especificado es $a' \in X'$ con $p(a') = \gamma(0) = a$. Dado que $a \in U_{x_0}$, el punto $a'$ debe pertenecer exactamente a una de estas láminas; denotémosla por $U'_{0, j_0}$. La restricción $p_{0, j_0} = p|_{U'_{0, j_0}}$ es un homeomorfismo de $U'_{0, j_0}$ hacia $U_{x_0}$. Definimos rigurosamente el levantamiento en este primer segmento como:
\begin{equation}
\gamma'(t) = (p_{0, j_0})^{-1}(\gamma(t)) \quad \text{para todo } t \in [0, t_1].
\end{equation}
Al ser composición de una función continua ($\gamma$) y un homeomorfismo inverso, $\gamma'$ es continua en $[0, t_1]$. Además, satisface la condición inicial $\gamma'(0) = (p_{0, j_0})^{-1}(a) = a'$ y verifica $p(\gamma'(t)) = p( (p_{0, j_0})^{-1}(\gamma(t)) ) = \gamma(t)$.

\textit{Paso Inductivo:} Supongamos que, para un entero $k$ con $0 \leq k < r$, ya hemos construido de manera continua el levantamiento $\gamma': [0, t_{k+1}] \to X'$ que satisface $p \circ \gamma' = \gamma$ en dicho dominio. Queremos extenderlo al intervalo contiguo $[t_{k+1}, t_{k+2}]$.
Sabemos que $\gamma([t_{k+1}, t_{k+2}]) \subset U_{x_{k+1}}$, un entorno bien cubierto. Su preimagen bajo $p$ es una unión disjunta de hojas $\bigsqcup U'_{k+1, j}$. Del paso inductivo anterior, ya conocemos el valor del punto de empalme $\gamma'(t_{k+1}) \in X'$. Evaluando la composición, $p(\gamma'(t_{k+1})) = \gamma(t_{k+1})$, el cual ciertamente pertenece a $U_{x_{k+1}}$. Por lo tanto, el punto evaluado $\gamma'(t_{k+1})$ yace en exactamente una lámina de la descomposición, digamos la hoja $U'_{k+1, j_{k+1}}$.
Definimos la extensión del levantamiento en el nuevo segmento mediante el homeomorfismo inverso de la hoja seleccionada:
\begin{equation}
\gamma'(t) = (p_{k+1, j_{k+1}})^{-1}(\gamma(t)) \quad \text{para todo } t \in [t_{k+1}, t_{k+2}].
\end{equation}
Esta nueva definición es evidentemente continua. Lo fundamental aquí es probar la compatibilidad en el nodo de intersección $t_{k+1}$. Evaluando la nueva definición en $t = t_{k+1}$ arroja $(p_{k+1, j_{k+1}})^{-1}(\gamma(t_{k+1}))$, lo cual por elección de la hoja es exactamente el punto $\gamma'(t_{k+1})$ originado por el paso inductivo previo. Ya que las dos definiciones concuerdan en el conjunto cerrado $\{ t_{k+1} \}$ que es la intersección de los dominios cerrados $[0, t_{k+1}]$ y $[t_{k+1}, t_{k+2}]$, el \textit{Lema del Pegado} (pasting lemma de la topología general) nos garantiza que la función extendida resultante, definida ahora en $[0, t_{k+2}]$, es una función unificada y continua en todo su dominio que sigue satisfaciendo el levantamiento.
Por inducción matemática finita, extendemos esto hasta alcanzar el punto final $t_{r+1} = 1$, completando así la construcción del levantamiento continuo global $\gamma': [0,1] \to X'$.

\textbf{Paso 3: Unicidad del levantamiento construido.}
Es fundamental asegurar que no existe otra forma de haber construido este camino. Supongamos la existencia de un segundo levantamiento continuo $\sigma: [0,1] \to X'$ que cumple $p \circ \sigma = \gamma$ con la misma condición de arranque $\sigma(0) = a'$. Invocamos el Lema de Unicidad de Levantamientos demostrado en la sección anterior. Tomamos $Y = [0,1]$, el cual es paradigmáticamente un espacio topológico conexo. Las funciones $\gamma', \sigma: [0,1] \to X'$ son ambos levantamientos de la misma aplicación continua $\gamma$. Además, verifican la condición de coincidencia inicial $\gamma'(0) = \sigma(0) = a'$. Dado que la aplicación de recubrimiento $p$ es, por la Observación 2.1.2.1, un homeomorfismo local, y el espacio $X'$ es de Hausdorff, el Lema de Unicidad nos dicta perentoriamente que $\gamma'(t) = \sigma(t)$ para la totalidad de $t \in [0,1]$. De este modo, el levantamiento no solo existe, sino que es singular y único.
\end{proof}

\begin{observacion}[Caracterización por levantamiento de curvas]
La propiedad demostrada en la proposición anterior no es tan solo un subproducto de ser un espacio recubridor, sino que constituye una \textit{caracterización} profunda. Como se demostrará de manera rigurosa más adelante en la teoría (específicamente en el Teorema 2.5.1 de la Sección 2.5), si una aplicación $p: X' \to X$ es un homeomorfismo local entre espacios apropiados y posee la propiedad de \textit{levantamiento único de curvas} para cualquier curva en la base y cualquier condición inicial admisible, entonces $p$ es, de manera forzosa y necesaria, un recubrimiento. Este es el pilar central de la correspondencia entre la topología algebraica y la teoría de revestimientos analíticos que motiva el estudio del haz de funciones holomorfas.
\end{observacion}

\begin{observacion}[Conexidad topológica y conexidad por arcos en variedades]
Aclaramos aquí una equivalencia crucial: en el contexto estricto de las variedades topológicas, un espacio es conexo si y sólo si es conexo por arcos. La demostración se apoya en la propiedad de conexidad local de la variedad. 
$\Rightarrow$ (Conexo implica conexo por arcos): Por definición, una variedad es localmente homeomorfa a un espacio euclidiano $\mathbb{R}^n$. Toda bola abierta euclidiana es trivialmente conexa por arcos (el segmento rectilíneo siempre las une). Por consiguiente, todo punto de una variedad admite una base de entornos conexos por arcos; es decir, toda variedad topológica es un espacio \textit{localmente conexo por arcos}. 
En cualquier espacio localmente conexo por arcos, las componentes conexas por arcos son forzosamente conjuntos abiertos: dado un punto $x$ en la componente por arcos $C$, su entorno básico conexo por arcos estará íntegramente contenido en $C$, por lo que $C$ contiene una vecindad de cada uno de sus puntos y es abierto. Como todas las componentes conexas por arcos particionan el espacio en disjuntos, si el espacio es conexo, no puede ser escrito como la unión de múltiples abiertos disjuntos no vacíos. En conclusión, solo puede haber una única componente conexa por arcos, haciendo al espacio enteramente conexo por arcos.
$\Leftarrow$ (Conexo por arcos implica conexo): Este es un teorema clásico de la topología general válido para cualquier espacio. Si existiese una partición del espacio en dos abiertos disjuntos, la imagen continua de cualquier camino (arco) que una un punto del primer abierto con un punto del segundo tendría que saltar entre abiertos disjuntos, violando la conexidad intrínseca del intervalo continuo $[0,1]$.
Esta equivalencia es capital para el análisis complejo, donde invariablemente asumiremos que nuestras superficies conexas permiten trazar caminos continuos entre cualesquiera dos puntos, permitiéndonos realizar procesos de continuación a lo largo de dichas curvas.
\end{observacion}

---

\subsection{2.2. El haz de gérmenes de funciones holomorfas}

\begin{definicion}[Gérmenes de funciones holomorfas]
Fijemos un punto $a \in \mathbb{C}$. Consideremos el conjunto de todos los pares $(U, f)$ donde $U \subset \mathbb{C}$ es un entorno abierto que contiene al punto $a$, y $f \in \mathcal{H}(U)$ (es decir, $f$ es una función holomorfa definida en el dominio $U$).
Sobre este conjunto de pares establecemos una relación denotada por $\sim_a$: diremos que $(U, f) \sim_a (V, g)$ si y solo si existe un entorno abierto más pequeño $W \subset U \cap V$ que también contiene a $a$, tal que la restricción de las funciones coincida idénticamente en $W$, esto es, $f|_W \equiv g|_W$.

Verificamos rigurosamente que $\sim_a$ es una relación de equivalencia:
\begin{itemize}
    \item \textbf{Reflexividad:} $(U, f) \sim_a (U, f)$ se verifica trivialmente eligiendo el entorno común $W = U$, donde evidentemente $f|_U \equiv f|_U$.
    \item \textbf{Simetría:} Si $(U, f) \sim_a (V, g)$, existe $W \subset U \cap V$ con $f|_W = g|_W$. Por la propiedad simétrica de la igualdad funcional, $g|_W = f|_W$, de modo que $(V, g) \sim_a (U, f)$ usando el mismo entorno $W$.
    \item \textbf{Transitividad:} Supongamos que $(U, f) \sim_a (V, g)$ y $(V, g) \sim_a (Z, h)$. Entonces existe $W_1 \subset U \cap V$ conteniendo a $a$ donde $f \equiv g$, y existe $W_2 \subset V \cap Z$ conteniendo a $a$ donde $g \equiv h$. Definimos la intersección $W = W_1 \cap W_2$. Al ser la intersección finita de conjuntos abiertos que contienen a $a$, $W$ es un entorno abierto de $a$. Además, para todo $z \in W$, se tiene $f(z) = g(z)$ (pues $z \in W_1$) y $g(z) = h(z)$ (pues $z \in W_2$). Por transitividad de la igualdad compleja, $f(z) = h(z)$ para todo $z \in W$. Por lo tanto, $(U, f) \sim_a (Z, h)$.
\end{itemize}

A cada clase de equivalencia bajo esta relación se le llama un \textbf{germen} de función holomorfa en el punto $a$. Denotaremos la clase de equivalencia que contiene al representante $(U, f)$ por $f_a$.
El conjunto de la totalidad de gérmenes en el punto $a$ se denomina el \textbf{tallo} en $a$, y se denota mediante el símbolo $\mathcal{O}_a$.

Definimos el \textbf{valor} de un germen $f_a$ en el punto $a$ evaluando cualquiera de sus representantes en $a$. Es imperativo justificar que esta operación es independiente del representante escogido: si tomamos dos representantes diferentes $(U, f)$ y $(V, g)$ de la misma clase $f_a$, por definición de la relación de equivalencia existe un entorno abierto $W \subset U \cap V$ que contiene a $a$ donde $f \equiv g$. Al evaluar específicamente en el punto $a \in W$, se sigue inevitablemente que $f(a) = g(a)$. De manera análoga y por medio del teorema de derivación puntual en análisis complejo, las derivadas de un germen están bien definidas: dado que $f \equiv g$ en un disco abierto $D(a, r) \subset W$, todas las derivadas de orden superior de las funciones deben coincidir en el centro del disco, esto es, $f^{(n)}(a) = g^{(n)}(a)$ para todo $n \geq 0$.
\end{definicion}

\begin{proposicion}[Estructura de anillo local de $\mathcal{O}_a$]
Las operaciones algebraicas puntuales de suma, multiplicación y multiplicación por escalares complejos dotan al tallo $\mathcal{O}_a$ de la estructura algebraica formal de un anillo conmutativo con unidad y, más aún, de una $\mathbb{C}$-álgebra. Además, $\mathcal{O}_a$ es un \textbf{anillo local}, lo cual significa algebraicamente que posee un y solo un ideal maximal, el cual está dado explícitamente por:
\begin{equation}
\mathfrak{m}_a = \{ f_a \in \mathcal{O}_a : f_a(a) = 0 \}.
\end{equation}
\end{proposicion}
\begin{proof}
Procederemos verificando exhaustivamente las afirmaciones algebraicas requeridas para dotar de estas estructuras a los gérmenes.

\textbf{1. Buena definición de las operaciones inducidas:}
Definimos la suma y el producto de gérmenes mediante la suma y el producto de representantes. Sean $f_a, g_a \in \mathcal{O}_a$ representados por $(U, f)$ y $(V, g)$ respectivamente. Definimos la suma como $(f+g)_a$, el germen de la función $f+g$ definida en el dominio intersecado $U \cap V$. Debemos mostrar que esto es independiente de los representantes elegidos. Supongamos que escogemos otros representantes $(U', f')$ y $(V', g')$ de manera que $(U, f) \sim_a (U', f')$ y $(V, g) \sim_a (V', g')$. Por definición, existen entornos de coincidencia $W_1 \subset U \cap U'$ donde $f \equiv f'$, y $W_2 \subset V \cap V'$ donde $g \equiv g'$. Tomemos el abierto intersección $W = W_1 \cap W_2$. En este dominio común $W$, las sumas puntuales son idénticas: para todo $z \in W$, $(f+g)(z) = f(z)+g(z) = f'(z)+g'(z) = (f'+g')(z)$. De modo riguroso, esto implica que $(U \cap V, f+g) \sim_a (U' \cap V', f'+g')$, con lo cual la suma de gérmenes está bien definida unívocamente. El argumento para el producto $f \cdot g$ es enteramente análogo, apoyándose en la distributividad y asociatividad de las operaciones en $\mathbb{C}$.

\textbf{2. $\mathcal{O}_a$ es un anillo conmutativo con identidad y $\mathbb{C}$-álgebra:}
Los axiomas formales de un anillo conmutativo (asociatividad y conmutatividad de suma y producto, y la existencia de elementos neutros y distributividad) se heredan directamente de la correspondiente álgebra de funciones holomorfas punto a punto $\mathcal{H}(W)$.
El cero aditivo de $\mathcal{O}_a$ es el germen $0_a$ de la función idénticamente nula $z \mapsto 0$ sobre todo el plano $\mathbb{C}$.
La unidad o elemento neutro multiplicativo de $\mathcal{O}_a$ es el germen $1_a$ de la función constante unitaria $z \mapsto 1$.
La multiplicación escalar por $\lambda \in \mathbb{C}$ se define multiplicando por el germen de la función constante $\lambda$, proveyendo los axiomas de $\mathbb{C}$-álgebra vectorial.

\textbf{3. Caracterización de las unidades (elementos invertibles) del anillo:}
Afirmamos que un germen $f_a \in \mathcal{O}_a$ es invertible (existe $g_a \in \mathcal{O}_a$ tal que $f_a \cdot g_a = 1_a$) si y solo si su evaluación en el centro no es cero, es decir, $f_a(a) \neq 0$.
Demostramos ambas implicaciones de forma meticulosa:
$\Rightarrow$ Supongamos que $f_a$ es unidad. Entonces, por definición de invertibilidad en un anillo, existe un germen inverso $g_a$ tal que $f_a \cdot g_a = 1_a$. Pasando a los valores evaluados en $a$, que sabemos es una operación bien definida, obtenemos que $(f_a \cdot g_a)(a) = 1_a(a)$. Esto equivale numéricamente a la ecuación en $\mathbb{C}$: $f_a(a) \cdot g_a(a) = 1$. Dado que en el cuerpo de los números complejos el producto es $1$, ninguno de los factores puede ser cero. Consecuentemente, $f_a(a) \neq 0$.
$\Leftarrow$ Recíprocamente, supongamos que el valor del germen no es nulo, $f_a(a) = f(a) \neq 0$. Por el teorema de continuidad de funciones, como $f$ es holomorfa y por tanto continua en algún entorno $U$ de $a$, debe existir un entorno más pequeño $V \subset U$ que contenga a $a$ en el cual la función no alcance el valor cero, esto es, $f(z) \neq 0$ para todo $z \in V$. Sobre este dominio restrictivo, la función recíproca $g(z) = 1/f(z)$ está perfectamente bien definida y, por las reglas de derivación del cociente, es una función holomorfa, $g \in \mathcal{H}(V)$. El germen $g_a$ definido por la clase $(V, g)$ cumple manifiestamente que para todo $z \in V$, $(f \cdot g)(z) = f(z) \cdot (1/f(z)) = 1$. Por definición de equivalencia de gérmenes, el producto de sus gérmenes coincide con la unidad: $f_a \cdot g_a = 1_a$. Así, $f_a$ es completamente invertible en el anillo de gérmenes.

\textbf{4. $\mathfrak{m}_a$ es un ideal:}
Definimos el conjunto de gérmenes que se anulan en $a$: $\mathfrak{m}_a = \{ f_a \in \mathcal{O}_a : f_a(a) = 0 \}$. Verificamos los dos axiomas de un ideal algebraico:
- Cerradura bajo la suma aditiva: si tomamos $f_a, g_a \in \mathfrak{m}_a$, entonces la suma evaluada $(f_a + g_a)(a) = f_a(a) + g_a(a) = 0 + 0 = 0$. Por tanto, el germen suma $f_a + g_a$ también reside en $\mathfrak{m}_a$.
- Absorción bajo el producto por elementos del anillo: sea $f_a \in \mathfrak{m}_a$ un germen que se anula en $a$, y sea un elemento arbitrario cualquiera $h_a \in \mathcal{O}_a$. Evaluamos el producto en $a$: $(h_a \cdot f_a)(a) = h_a(a) \cdot f_a(a) = h_a(a) \cdot 0 = 0$. De este modo, $h_a \cdot f_a \in \mathfrak{m}_a$.
Estos dos pasos fundamentan innegablemente que $\mathfrak{m}_a$ constituye un ideal propio del anillo (es propio ya que la identidad $1_a \notin \mathfrak{m}_a$).

\textbf{5. $\mathfrak{m}_a$ es el único ideal maximal (propiedad de ser Anillo Local):}
Para comprobar que $\mathcal{O}_a$ ostenta la propiedad algebraica de anillo local, basta con demostrar que el conjunto de los elementos no invertibles coincide con un ideal, que resultará forzosamente ser el ideal maximal único del anillo.
Por la caracterización previamente obtenida en el Paso 3, el conjunto de los elementos no invertibles del anillo es precisamente $\mathcal{O}_a \setminus (\mathcal{O}_a)^{\times} = \{ f_a : f_a(a) = 0 \} = \mathfrak{m}_a$.
Supongamos ahora que existe algún otro ideal propio $J \subsetneq \mathcal{O}_a$. Todo ideal propio consiste estrictamente en elementos no invertibles. La justificación de esto es un resultado básico del álgebra: si un ideal contuviera un elemento invertible $u \in J$, la propiedad de absorción forzaría a que el producto del inverso por el elemento, $u^{-1} \cdot u = 1_a$, pertenezca al ideal. Tener a la unidad $1_a \in J$ obligaría a que el ideal absorbiese a todo elemento $r = r \cdot 1_a \in J$, haciendo que $J$ colapsara en el anillo total $\mathcal{O}_a$, contradiciendo su estatus de ideal propio.
Puesto que todos los elementos de cualquier ideal propio $J$ deben ser obligatoriamente no invertibles, deben pertenecer íntegramente al conjunto de no invertibles, que es $\mathfrak{m}_a$. Esto arroja la inclusión ineludible $J \subset \mathfrak{m}_a$.
Dado que todo ideal propio está contenido como subconjunto de $\mathfrak{m}_a$, concluimos irrefutablemente que $\mathfrak{m}_a$ es el supremo de todos los ideales propios y por tanto es el ideal maximal, y que ningún otro ideal maximal puede existir, probando el carácter de anillo local.
\end{proof}

\begin{lema}[Cuerpo residual del tallo holomorfo]
El cociente algebraico del anillo local de gérmenes por su ideal maximal es un cuerpo isomorfo a los números complejos, es decir, el cuerpo residual satisface:
\begin{equation}
\mathcal{O}_a / \mathfrak{m}_a \cong \mathbb{C}.
\end{equation}
\end{lema}
\begin{proof}
Construimos un homomorfismo de evaluación explícito para invocar los teoremas de isomorfía. Definimos la aplicación $\text{ev}_a: \mathcal{O}_a \to \mathbb{C}$ mediante la fórmula de evaluación en el centro $\text{ev}_a(f_a) = f_a(a)$, donde el valor se obtiene tomando el representante como se formalizó anteriormente.
Comprobamos que es un homomorfismo de anillos válido de $\mathbb{C}$-álgebras:
La linealidad y preservación del producto surgen de forma instantánea de la definición de las operaciones puntuales en $\mathcal{O}_a$:
- Suma: $\text{ev}_a(f_a + g_a) = (f+g)_a(a) = f(a) + g(a) = \text{ev}_a(f_a) + \text{ev}_a(g_a)$.
- Producto: $\text{ev}_a(f_a \cdot g_a) = (f \cdot g)_a(a) = f(a) \cdot g(a) = \text{ev}_a(f_a) \cdot \text{ev}_a(g_a)$.
- Elemento identidad: $\text{ev}_a(1_a) = 1_a(a) = 1$.
La aplicación es manifiestamente sobreyectiva. Dado cualquier número complejo arbitrario $c \in \mathbb{C}$, consideremos la función constante $f(z) = c$ definida ubicuamente sobre el plano complejo $\mathbb{C}$. Esta función analítica induce un germen constante $c_a \in \mathcal{O}_a$, cuya evaluación es precisamente $\text{ev}_a(c_a) = c$.
Examinamos ahora el núcleo del homomorfismo. Por definición algebraica clásica, el núcleo se compone de:
\begin{equation}
\ker(\text{ev}_a) = \{ f_a \in \mathcal{O}_a : \text{ev}_a(f_a) = 0 \} = \{ f_a \in \mathcal{O}_a : f_a(a) = 0 \}.
\end{equation}
Reconociendo de la proposición anterior esta definición constructiva, dictaminamos que el núcleo coincide idénticamente con el ideal maximal: $\ker(\text{ev}_a) = \mathfrak{m}_a$.
Haciendo un llamamiento formal al Primer Teorema de Isomorfía de anillos de Noether, el cociente del anillo dominio sobre el núcleo del homomorfismo es canónicamente isomorfo a la imagen del homomorfismo en el codominio. Puesto que el mapa es exhaustivo (sobreyectivo), su imagen es la totalidad de $\mathbb{C}$. Consiguientemente, el cociente adquiere la estructura de cuerpo y se concluye el isomorfismo deseado: $\mathcal{O}_a / \mathfrak{m}_a \cong \mathbb{C}$.
\end{proof}

\begin{definicion}[El haz $\mathcal{O}$ y su construcción topológica]
Procedemos a agrupar todos los tallos locales para conformar un espacio total que encapsule todas las funciones holomorfas locales de manera unificada. Se define el espacio subyacente del haz de gérmenes de funciones holomorfas como la enorme unión completamente disjunta de los tallos asociados a cada punto particular del plano complejo:
\begin{equation}
\mathcal{O} = \bigsqcup_{a \in \mathbb{C}} \mathcal{O}_a.
\end{equation}
Todo elemento de $\mathcal{O}$ es por tanto un germen $f_a$ anclado en algún punto unívoco de $\mathbb{C}$.
Para imbuir este mero conjunto teórico en un espacio topológico coherente, debemos instaurar una colección de conjuntos abiertos. Se definen las colecciones de vecindades básicas de la siguiente manera:
Sea $U$ un conjunto abierto arbitrario en $\mathbb{C}$, y sea $f \in \mathcal{H}(U)$ una función holomorfa sobre dicho dominio. Definimos la sección del germen tubular asociada al par $(U, f)$ como el subconjunto de $\mathcal{O}$:
\begin{equation}
\mathcal{N}(U, f) = \{ f_z \in \mathcal{O} : z \in U \},
\end{equation}
donde el símbolo $f_z$ representa el germen inducido por la función analítica $f$ analizada localmente en el punto base particular $z \in U$.
Se decreta que la familia inmensa de todas las secciones $\mathcal{N}(U, f)$, generada iterando sobre todos los dominios abiertos admisibles $U \subset \mathbb{C}$ y todas sus funciones holomorfas $f \in \mathcal{H}(U)$, forma una base para generar la topología rigurosa del espacio $\mathcal{O}$.

Complementariamente, se introduce la \textbf{proyección natural} (o map de anclaje) $p: \mathcal{O} \to \mathbb{C}$, cuya acción extrae la posición subyacente del germen. Formalmente, $p(f_a) = a$.
\end{definicion}

\begin{observacion}[Verificación exhaustiva de los axiomas de la base topológica]
Para que la decisión de bautizar a la colección $\{ \mathcal{N}(U, f) \}$ como una "base" adquiera legitimidad topológica, es imperativo certificar meticulosamente que satisface los dos axiomas fundacionales de las bases de topologías de conjuntos.

\textbf{Axioma 1: Recubrimiento total.} La unión de todos los elementos básicos propuestos debe reconstituir el espacio entero $\mathcal{O}$.
Comprobación: Sea $f_a \in \mathcal{O}$ cualquier germen arbitrario existente. Por el propio concepto formador de germen, este nace de algún representante explícito consistente de un entorno abierto $U$ con centro en $a$ y una función holomorfa $f \in \mathcal{H}(U)$. El conjunto básico conformado por la dupla de este representante originario es $\mathcal{N}(U, f)$. Y trivialmente, al evaluar la colección de gérmenes en los puntos de este conjunto $U$, el germen en el centro es $f_a$. Así, $f_a \in \mathcal{N}(U, f)$. Como hemos encontrado un conjunto básico para cada punto del espacio total, el axioma de recubrimiento se verifica irrefutablemente.

\textbf{Axioma 2: Intersección básica refinada.} Si un elemento $h_z$ pertenece a la intersección de dos conjuntos de la base $\mathcal{N}(U, f) \cap \mathcal{N}(V, g)$, debe existir un tercer conjunto de la base, llamémoslo $\mathcal{N}(W, h)$, tal que el elemento esté contenido en este, y el conjunto esté enteramente constreñido en la intersección: $h_z \in \mathcal{N}(W, h) \subset \mathcal{N}(U, f) \cap \mathcal{N}(V, g)$.
Comprobación pormenorizada: Si el germen $h_z$ reside en la intersección de marras, esto significa que simultáneamente pertenece a la colección tubular de $f$ y a la de $g$. En otras palabras, la proyección del germen nos asegura que el punto base de anclaje es $z$, requiriendo que el punto $z \in U$ y $z \in V$ forzosamente, con lo que $z$ recae en el entorno abierto intersección $U \cap V$.
A nivel de gérmenes, la equivalencia postula que $h_z$ concuerda con el germen inducido por $f$ en $z$ y, análogamente, con el germen inducido por $g$ en $z$. Ello es: $h_z = f_z = g_z$.
El hecho de que los gérmenes de $f$ y $g$ sean idénticos en el punto $z$ reviste un significado fenomenológico explícito: existe un entorno de compatibilidad, un abierto $W \subset U \cap V$ centrado alrededor del punto $z$, dentro del cual ambas funciones holomorfas evalúan resultados idénticos para cada punto, formalizado como $f|_W \equiv g|_W$.
Definamos la restricción coincidente $h = f|_W = g|_W$, que claramente retiene su condición holomorfa ($h \in \mathcal{H}(W)$). Construimos el nuevo conjunto de base intermedio $\mathcal{N}(W, h)$.
Primeramente comprobamos si el elemento descansa dentro del conjunto. Evaluando la función $h$ localmente en el punto $z \in W$ nos entrega el germen $h_z$, que coincide de pleno con la definición original, probando $h_z \in \mathcal{N}(W, h)$.
En segunda instancia verificamos la contención del conjunto nuevo dentro de los progenitores originales. Escojamos cualquier germen arbitrario $h_w \in \mathcal{N}(W, h)$ anclado en cualquier punto general $w \in W$. Puesto que sobre todo el dominio de $W$ la función satisface universalmente $h(w) = f(w)$, la inducción del germen en el punto evalúa que el germen de $h$ alrededor de $w$ es equivalente al germen de $f$ en su contorno de $w$. Luego $h_w = f_w$. Como el punto $w$ pertenece a $W$ y por herencia $W \subset U$, deducimos que el germen en cuestión $h_w$ ostenta residencia en $\mathcal{N}(U, f)$. Este argumento simétrico es replicable cambiando la literal a $g$, dictando que $h_w = g_w$ de forma que se inserta sin contradicción en $\mathcal{N}(V, g)$. Habiendo probado rigurosamente que $\mathcal{N}(W, h) \subset \mathcal{N}(U, f) \cap \mathcal{N}(V, g)$, ratificamos concluyentemente que la familia constituye una base legítima para la topología descrita.
\end{observacion}

\begin{proposicion}[El espacio $\mathcal{O}$ exhibe la propiedad de Hausdorff]
El espacio topológico construido de gérmenes holomorfos $\mathcal{O}$ satisface incondicionalmente el axioma de separación de Hausdorff ($T_2$): cualesquiera dos gérmenes distintos $f_a \neq g_b$ pueden encapsularse y separarse mediante la utilización de sendos entornos abiertos que sean mutuamente excluyentes (disjuntos).
\end{proposicion}
\begin{proof}
La demostración debe abordar inexcusablemente dos escenarios lógicos que agotan la posibilidad de diferencia entre dos gérmenes: bien sus puntos base geográficos en el plano son distintos, o bien, residiendo en la misma coordenada base, sus contenidos funcionales intrínsecos difieren de forma irreconciliable.

\textbf{Caso A: Los puntos base difieren ($a \neq b$).}
Esta eventualidad es la más geométrica. Puesto que la topología del plano euclidiano $\mathbb{C}$ satisface intrínsecamente la propiedad de Hausdorff, garantizamos la existencia de dos discos abiertos y desvinculados $U \subset \mathbb{C}$ y $V \subset \mathbb{C}$ tales que contengan a los focos correspondientes $a \in U$ y $b \in V$, operando bajo la cláusula estricta de vacuidad de su intersección, $U \cap V = \emptyset$.
Los gérmenes bajo escrutinio devienen de representantes locales; digamos que $f_a$ está instanciado por un par originario $(\widetilde{U}, \widetilde{f})$ y el otro germen $g_b$ por $(\widetilde{V}, \widetilde{g})$. Para adecuar los dominios a nuestra geometría separada, creamos restricciones abiertas encajadas, tomando los conjuntos intersección $U' = U \cap \widetilde{U}$ y $V' = V \cap \widetilde{V}$. Naturalmente, estos preservan la contención central y la condición de separación disjunta, $U' \cap V' \subset U \cap V = \emptyset$.
Convocamos entonces los entornos básicos del haz proyectados sobre estas nuevas regiones: consideramos $\mathcal{N}(U', \widetilde{f}|_{U'})$ como un entorno de $f_a$ y $\mathcal{N}(V', \widetilde{g}|_{V'})$ como entorno para $g_b$. ¿Podrían estos intersectarse en el espacio de gérmenes de alguna forma insidiosa? De existir un germen hipótetico común $h_z$ radicado simultáneamente en las entrañas de ambos conjuntos abiertos constructivos, la aplicación de la proyección cartográfica $p$ develaría irremediablemente su anclaje, sentenciando que la base del germen $p(h_z) = z$ debe cohabitar lícitamente en el conjunto base del primer entorno ($z \in U'$) y también en el substrato del segundo entorno ($z \in V'$). Sin embargo, este suceso conculca el decreto axiomático previo de que $U' \cap V' = \emptyset$. Frente a semejante imposibilidad, deducimos implacablemente que $\mathcal{N}(U', \widetilde{f}|_{U'}) \cap \mathcal{N}(V', \widetilde{g}|_{V'}) = \emptyset$.

\textbf{Caso B: Los puntos base convergen en el mismo epicentro ($a = b$), pero los gérmenes discrepan lógicamente ($f_a \neq g_a$).}
Escojamos representantes locales de validación de los gérmenes $(\widetilde{U}, \widetilde{f})$ y $(\widetilde{V}, \widetilde{g})$ que ambos centran su influencia analítica gravitando en el origen común $a$. Seleccionamos cuidadosamente el epicentro geográfico trazando un disco continuo, de topología convexa y manifiestamente conexa, $D = D(a, \epsilon) \subset \widetilde{U} \cap \widetilde{V}$. Este dominio circular será nuestra arena de prueba.
Asumimos hipotéticamente que, no obstante la discrepancia originaria fundamental en el centro, la construcción perimétrica no logra aislarlos, sucediendo que sus entornos engendrados sobre este disco sufren de intersección espuria: $\mathcal{N}(D, \widetilde{f}|_{D}) \cap \mathcal{N}(D, \widetilde{g}|_{D}) \neq \emptyset$.
El intruso de esta intersección, digamos un escurridizo germen común $h_z$ localizado en un desplazamiento radial $z \in D$, ostenta el requerimiento semántico de que el germen evaluado de $\widetilde{f}$ alrededor de la vecindad infinitesimal de $z$ iguale precisamente el de $\widetilde{g}$, conminado así: $\widetilde{f}_z = \widetilde{g}_z$.
La identidad local de gérmenes obliga ineludiblemente a la concordancia funcional integral de las ecuaciones evaluadas en un pequeño pero finito entorno $W \subset D$ rodeando holgadamente a $z$. Es decir, sobre este paño $W$, la diferencia residual de las funciones es nula: $(\widetilde{f} - \widetilde{g})|_{W} \equiv 0$.
A estas alturas de la argumentación hacemos gravitar todo el peso conceptual sobre el pilar arquitectónico superlativo del análisis holomorfo global: el Teorema del Principio de Identidad Analítica (Continuación Analítica en dominios). Recordamos sus preceptos: Si una función que profesa holomorfía universal en una región que detenta la innegociable condición de ser abierta y \textit{conexa} se subyuga a anularse en un subconjunto con punto de acumulación interior (y un disco abierto innegablemente encierra un mar continuo de acumulación), se declara matemáticamente extinguida y nulificada de punta a cabo en el global de la región.
Invocamos la aplicación explícita: La función sustracción $H = \widetilde{f} - \widetilde{g}$ goza de credenciales holomorfistas incontrovertibles sobre el disco integrador $D$, el cual es topológicamente conexo por excelencia. Se ha corroborado recién su extinción en el entorno abierto local $W$. Con la ley en mano, el postulado del principio analítico sentencia irrevocablemente que $H \equiv 0$ a través de toda la planicie del disco $D$, por lo que la igualdad se traslada a $\widetilde{f}(w) = \widetilde{g}(w)$ para cada átomo matemático del conjunto $\forall w \in D$.
Dado que $D$ actúa lícitamente de entorno de anclaje para el núcleo focal original $a$, y dado que amparado en dicho disco observamos congruencia ubicua y abarcadora en su extensión, la definición basilar de los elements del tallo forzaría irremediablemente a que los gérmenes originarios fuesen isomorfos y permutables, $\widetilde{f}_a = \widetilde{g}_a$. 
Empero, el fundamento hipotético postulado en la entrada de este subcaso B asumía explícitamente y sin matices que los gérmenes gozaban de divergencia formal irreconciliable $f_a \neq g_a$. El cisma lógico derivado y la contradicción patente evidencian que el pecado original consistió en la elucubración de una intersección en sus contornos perimétricos. Consagrados de la imposibilidad de solapamiento, declaramos victoriosos los entornos $\mathcal{N}(D, \widetilde{f}|_{D})$ y $\mathcal{N}(D, \widetilde{g}|_{D})$ disjuntos, separando los gérmenes de forma terminal y ratificando la propiedad global de Hausdorff para el haz entero $\mathcal{O}$.
\end{proof}

\begin{proposicion}[La proyección natural $p$ es un homeomorfismo local]
La aplicación de proyección base $p: \mathcal{O} \to \mathbb{C}$ que asigna a cada ente topológico su coordenada de anclaje territorial cumple el rol formal de un recubrimiento local cartográfico, operando fehacientemente como un homeomorfismo local puro.
\end{proposicion}
\begin{proof}
Desglosamos la rigurosa comprobación en sus componentes topológicos medulares para legitimar la naturaleza de la proyección:

\textbf{1. Continuidad irrestricta de $p$:}
Por el dogma fundamental de la topología general, requerimos atestiguar que la recolección de las preimágenes inversas evaluadas sobre abiertos territoriales desencadenan la reconstitución de conjuntos ostensiblemente abiertos en la superestructura del haz.
Seleccionemos de forma aleatoria y general un conjunto que ostente la condición formal de ser abierto en la planicie compleja inferior, $V \subset \mathbb{C}$. Inspeccionemos la retro-imagen estructural:
\begin{equation}
p^{-1}(V) = \{ f_z \in \mathcal{O} : p(f_z) = z \in V \}.
\end{equation}
Esta es la colección absoluta y maximal de todo germen viable imaginable adscrito a algún punto en $V$. Podemos transcribir esta multiplicidad formalmente mediante una unión exhaustiva barriendo sobre todos los subdominios que caen bajo su manto y de todas sus configuraciones holomorfas admitidas:
\begin{equation}
p^{-1}(V) = \bigcup_{ \substack{U \subset V \\ U \text{ abierto}}} \; \bigcup_{f \in \mathcal{H}(U)} \mathcal{N}(U, f).
\end{equation}
La justificación de la identidad del conjunto es meridiana: si poseemos cualquier elemento interno de la pre-imagen, digamos $f_z \in p^{-1}(V)$, este emana de un vector generador explícito $(\widetilde{U}, f)$. Tomamos la intersección territorial reencausada $U = \widetilde{U} \cap V$, que ciertamente funge como entorno de $z$. Así pues, la inducción de este germen es acogida por el colector tabular $\mathcal{N}(U, f|_U)$, el cual milita lícitamente en el megaconjunto indexado de la derecha. En retrospectiva, todo sumando tabular congrega su punto basamento de sus miembros dentro del confinamiento prescriptivo de $V$.
Como corolario ineluctable, el subespacio $p^{-1}(V)$ es expresable matemáticamente como el ensamblaje o unión masiva aglutinada de conjuntos $\mathcal{N}(U,f)$, los cuales constituyen la sustancia misma de la base fundadora de la topología superior del haz. Al componerse de bloques puramente elementales y abiertos constitutivos, su suma integral resulta un abierto formal e insoslayable en $\mathcal{O}$. Esto instaura el principio general de que las retro-imágenes conservan la apertura, dictaminando la continuidad universal.

\textbf{2. Transformación cartográfica de aplicación abierta:}
Se requiere contrastar si la inducción proyectiva respeta la condición de mapeo topológico abierto de salida, exportando la propiedad topológica hacia adelante a favor del codominio. Para esto será imperiosamente suficiente acreditarlo al nivel primordial de los sillares básicos constituyentes de la arquitectura, dado que la unión y suma se preserva ante la aplicación de funciones directas sobre la generalidad.
Elijamos un sillar abierto $\mathcal{N}(U, f)$. Por el orden de su definición estructural, el barrido geográfico $p(\mathcal{N}(U, f)) = p( \{ f_z : z \in U \} )$ extrae las coordenadas de anclaje de todos sus especímenes, abarcando íntegramente la vastedad del subconjunto territorial. Por ende, la imagen cartográfica arroja un dictamen limpio: $p(\mathcal{N}(U, f)) = U$.
Como el pre-requisito y definición fundadora dicta que $U$ es genéticamente abierto, la imagen recae obligatoriamente dentro del catálogo de los conjuntos abiertos. Una conjunción extensiva consolida que $p$ funciona en la globalidad como mapeo de naturaleza abierta.

\textbf{3. Biyectividad intrínseca de los vecindarios (Homeomorfismo Local):}
El cenit de la demostración es asentar que cada habitante del espacio $\mathcal{O}$ permite ser circundado por alguna aureola topológica cuya recíproca actúe de carta de deformación isomorfa fidedigna sobre el espacio base sin compresión ni desdoblamiento de capas.
Evaluamos para cualquier germen particular $f_a$ el entorno subyacente derivado $\mathcal{N}(U, f)$ amparado en el representante madre. La función constreñida locacionalmente $p|_{\mathcal{N}(U, f)} : \mathcal{N}(U, f) \to U$ retiene sus linajes de aplicación continua e imagen abierta. Comprobemos su rango biyectivo:
- Inyectividad restrictiva: Acotando sobre esta ventana tabular particular propongamos el hipotético fallo $\exists \; f_{z_1}, f_{z_2} \in \mathcal{N}(U, f)$ tales que sufran una compresión colisional en su retro-imagen basal $p(f_{z_1}) = p(f_{z_2})$. Esto desentraña numéricamente la igualación espacial $z_1 = z_2$. Ya que ambos radican sobre el idéntico sustrato posicional $z_1$, y al emerger en común de la evaluación tubular de la misma fórmula directriz originaria $f$, resultan inescrutablemente discernibles e iguales en forma de gérmenes: $f_{z_1} = f_{z_2}$. Ergo, la inyección queda asegurada.
- Sobreyectividad de cobertura: Tomando el dominio abarcado integral $U$, elegimos un punto destinatario prospectivo aleatorio $z \in U$. El constructo mismo de la forma tubular certifica la inmersión apriorística y permanente de un espécimen correlacionado evaluando la matriz generatriz sobre esta posición, concretando a la entidad $f_z \in \mathcal{N}(U, f)$. Así pues, $p(f_z) = z$, asegurando el espectro exhaustivo y sin lagunas.
Disfrutando la bendición del mapeo biyectivo que goza de la propiedad simultánea de continuidad y condición de aplicación abierta, la topología clásica rubrica sin resquicios que el ensamble resulta ser un homeomorfismo puro sobre el restringido escenario geográfico, cimentando y consolidando el título irrevocable de Homeomorfismo Local para el operador universal.
\end{proof}

\begin{corolario}[Estructura Varietal del haz $\mathcal{O}$]
El espacio envolvente de los gérmenes de funciones analíticas con estructura de haz, $\mathcal{O}$, congrega todos los atributos canónicos precisados para constituirse como una rigurosa Variedad Topológica inmersa en la dimensión doble ($n=2$).
\end{corolario}
\begin{proof}
Invocando retrospectivamente la amalgama de axiomas de la Definición 2.1.1 relativas a las exigencias fundacionales para adquirir el linaje varietal de dimensionalidad par, requerimos asegurar que la topología global obedezca el postulado de Hausdorff de separabilidad y, subsecuentemente de forma local, soporte vecindarios distorsionables homeomórficamente sobre celdas canónicas euclidianas abiertas del doblet analítico $\mathbb{C} \cong \mathbb{R}^2$.
La validación axiomática de Hausdorff se extrajo axiomáticamente del decurso analítico exhaustivo plasmado durante la demostración de la Proposición 2.2.3 (donde el Principio de Identidad Analítica se empleó providencialmente sin escatimar para esquivar la imposibilidad aparente).
Por añadidura colateral y solidificada en la Proposición inminentemente anterior (2.2.4), se proveyó la comprobación fidedigna de que para la totalidad de los moradores germinales $f_a \in \mathcal{O}$ emerge infaliblemente su aureola circunscrita de vecindario abierto $\mathcal{N}(U, f)$ pertrechada con la deformación cartográfica $p|_{\mathcal{N}(U, f)}$, la cual actúa como nexo homomórfico hacia un parche territorial estipulado como $U$, que es patentemente un ente abierto domiciliado en la geografía del plano $\mathbb{C}$. Estas entidades ofician majestuosamente de atlas locacional nativo y fungen como las Cartas Locales exigidas, cristalizando irrebatiblemente que el haz detenta una estructura íntima y topología conformada al modo de variedad de la especie dimensional plana.
\end{proof}

\begin{definicion}[Fenómeno de Continuación Analítica guiada a través de un rastro curvilíneo]
Abordaremos la formulación sistemática e integral sobre qué significa continuar las propiedades holomórficas siguiendo la trayectoria delineada por una curva continua. Concedámonos una ruta y curva cartográfica temporal $\gamma: [0, 1] \to \mathbb{C}$ operante sobre la pista de aterrizaje, estipulando su coordenada de partida original en el nudo $\gamma(0) = a$. Postulamos la inserción en este entorno posicional de un germen de inicio local, designado mediante $f_a \in \mathcal{O}_a$.
Bajo la nomenclatura de esta disertación, se dictamina formalmente la \textbf{Continuación Analítica} extendida del germen matriz inicial $f_a$, desplazándose guiada en congruencia con la directriz geométrica impuesta por $\gamma$, si existe la capacidad constructiva y materialización de un trazado continuo sobreelevado de \textbf{levantamiento paramétrico} $\widetilde{\gamma}: [0, 1] \to \mathcal{O}$, cuyas prerrogativas formales estipulen inexcusablemente el respeto íntegro de la coordenada inicial del espectro germinal superior enclavada en la génesis temporal, asumiendo $\widetilde{\gamma}(0) = f_a$, así como prescriban la obediencia subyugada a la proyección cartográfica de que toda su silueta sombreada siga puntillosa el trazo terrenal de la directriz en todo momento de evaluación, expresándose mediante el decálogo de conmutatividad matemática $p \circ \widetilde{\gamma} = \gamma$.

Es mandatorio articular una exégesis clarificadora para cohesionar y armonizar esta conceptualización modernista y global de revestimiento con la ideología procedimental clásica germinada primigeniamente por el patriarca Weierstrass sobre el plano secante de prolongaciones de disco a disco. El decurso analítico asevera que la variable de tránsito temporal $\widetilde{\gamma}(t)$ codifica sin ambages y a cada compás el \textbf{germen de alguna función formal analítica circunscrita y operativa} centrada imperativamente en la demarcación actual de la caravana posicional descrita por $\gamma(t)$. En particularidad asombrosa, este germen constituye la rama holomorfa genéticamente vinculada que perpetúa el patrimonio de identidades sin roturas procedentes de la etapa predecesora transcurrida. El engranaje y sincronización paramétrica del aparato levitante sella indisolublemente el marco y compila magistralmente una formalización unificadora y purista de cómo la prolongación clásica de entidades holomorfistas encarna, con esplendor natural e innato, una instancia axiomática generalizada en la doctrina pura de topología abstracta de \textit{Teoría de los espacios de recubrimiento y levantamiento}.
\end{definicion}

\begin{proposicion}[Ley Universal y Principio de la Permanencia irrenunciable de Relaciones Funcionales]
Consagrémonos a estipular que hemos concretado de forma exitosa y contigua el tránsito unificado y extensión funcional de no meramente uno, sino concurrentemente de dos especies germinales, originando las continuaciones expansivas y continuas en el estrato paramétrico denotadas $\mathcal{F}(t)$ y $\mathcal{G}(t)$, procedentes desde sus estados primitivos de cuna germinal representados en $f_a$ y $g_a$, desplegadas concomitantemente sobre el trayecto demarcatorio de un mismo sendero unificador $\gamma$.
Adicionemos la suposición medular de que en el escenario matricial y temporal inicial rige un nexo, ligadura o conjunción regida funcional o polimórficamente expresada, tal que las entrañas germinales sucumben de origen a una atadura de ecuación modelada como un polinomio o ecuación universal $P$: $P(f_a, g_a) = 0$. (Por ejemplo asúmase una equivalencia como $f_a \cdot g_a - 1_a = 0$ u otro constructo similar).
En tal conjunción de condicionamientos perentorios, se vaticina con rigurosidad impoluta y dogmática que la susodicha adherencia relacional se mantendrá inmarcesible, prolongando invariablemente la ecuación sobre la generalidad paramétrica y para cualquier tiempo de recorrido, forzando la obligatoriedad matemática imperecedera de:
\begin{equation}
P(\mathcal{F}(t), \mathcal{G}(t)) = 0_a \quad \text{para toda y cada una de las coordenadas temporales } t \in [0, 1].
\end{equation}
\end{proposicion}
\begin{proof}
Para el asedio analítico a este contencioso, reclutaremos nuevamente a la élite probatoria de la técnica "Topológica conectiva Abierto-Cerrada" encarnada sobre el espectro paramétrico dimensional real, diseccionando el conjunto subyacente demarcado por la fidelidad a la ley funcional y estipulado como $S = \{ t \in [0, 1] : P(\mathcal{F}(t), \mathcal{G}(t)) = 0 \}$. Nuestro faro directriz será justificar que este recinto se amolda simultáneamente a la naturaleza abierta e invariablemente a su complementaria clausura cerrada en la traza rectilínea.

\textbf{Fase Inicial - Vacuidad Prohibida:} Por el basamento de la postulación y decreto hipotético de condicionante de origen, evaluando cronológicamente al amanecer del trayecto en el t=0 se recaba la evidencia incontestable $P(\mathcal{F}(0), \mathcal{G}(0)) = P(f_a, g_a) = 0$. Esta validación concede a la raíz temporal cero carta de membresía legitima y sitúa imperativamente a $0 \in S$, desterrando ipso facto la noción de nulidad vacua y sentenciando que el ensamble posee cardinalidad positiva, $S \neq \emptyset$.

\textbf{Fase Central - Apertura Esencial:} Convoquemos aleatoria y arbitrariamente al estado probatorio a un interviniente genérico que ostente ciudadanía en el reducto, digamos el instante $t_0 \in S$. Nos adentramos a desenterrar los portadores locales que codifican los gérmenes congelados en esta temporalidad en específico. Extrapolamos dos enviados directos y representantes holomórficos válidos, denominados paramétricamente $\phi$ y $\psi$, con rango de operación delimitado en el interior de los linderos de algún Disco Abierto conexo en la cartografía plana $\Delta$, que alberga como pupila visual a la encrucijada y posición del trayecto descrita por $\gamma(t_0)$, logrando la correspondencia $\mathcal{F}(t_0) = \phi_{\gamma(t_0)}$ e inequívocamente $\mathcal{G}(t_0) = \psi_{\gamma(t_0)}$.
Ensamblando ahora el constructo de fusión de estas funciones madres engarzadas bajo la rúbrica y tutela algorítmica y algebraica prescrita por el dictamen relacional, engendramos la supra-función compuesta holomorfa total $\mathbb{H} = P(\phi, \psi)$. Al provenir de linajes indiscutiblemente analíticos por inducción de sumas y encajes productivos, su fuero holomorfo a través del disco expansivo y continuo territorial de $\Delta$ goza de reconocimiento inalienable, ostentando la investidura de la pertenencia gremial $\mathbb{H} \in \mathcal{H}(\Delta)$.
Por derecho de cuna estipulado por pertenencia del instante $t_0$ a la cofradía y asamblea central $S$, atestiguamos de manera certera e infalible que en este nódulo crucial el mandato se satisface: el germen extraíble del ente analítico naciente acoplado a la raíz germinal padece forzosamente la condena a nulificarse a plenitud en dicho punto y encrucijada temporal, o sea, $\mathbb{H}_{\gamma(t_0)} = 0_{\gamma(t_0)}$.
El estricto código que moldea los principios intrínsecos de gérmenes nos compele inevitablemente a concluir que la equivalencia del valor de anulación del ente funcional germinal hacia el cero absoluto germinal emana estrictamente de una obligación analítica territorial colindante: la función en su estado macroscópico evaluativo se extingue, evapora y colapsa ininterrumpidamente hacia la inanición del cero sobre todo un reducto circundante microscópico o entorno vecindario de validación encasillado próximo al centro $\gamma(t_0)$.
Se nos presenta nuevamente y en toda su magnitud la irrupción incondicional, sublime e inapelable del Principio y Doctrina Dogmática de la Identidad Analítica General para recintos conectivos y convexos como los circulares. Se confronta a una entidad formalmente analítica y holomorfista, operante en una provincia innegablemente conexa e ininterrumpida $\Delta$, que por vicisitudes padece el mal endémico de subyugación y anulación mortuoria en un conjunto de congregación poseedor ineluctablemente de acumulaciones interiores plenas. El laudo dictamina por unísono analítico el desvanecimiento y defunción general de la entidad: se atestigua formalmente y decreta con rigurosidad que $\mathbb{H}$ cede a perecer con identidad equivalente al cero en la completitud irrestricta territorial general y universal abarcada en sus fronteras de $\Delta$; se concluye así irrefutablemente que $\mathbb{H} \equiv 0$ inquebrantablemente a la redonda de este recinto perimetral abarcador de $\Delta$.
Prosiguientemente, dada la naturaleza perenne y continua que asiste en todo momento sobre el devenir cinemático a las trazas levantadas $\mathcal{F}$ y $\mathcal{G}$, en concertación simultánea y hermanada con la proyección directriz $\gamma$, y evaluado colindando proximidades ínfimas adheridas al estallido originario $t_0$, el desplazamiento sutil en las coordenadas de variables libres transita la frontera adyacente para los instantes infinitesimales $t \approx t_0$. Bajo esta vecindad los dictámenes determinan forzosamente $\mathcal{F}(t) = \phi_{\gamma(t)}$ de igual manera indisoluble que lo prescribe $\mathcal{G}(t) = \psi_{\gamma(t)}$, asumiendo su rol como portadores en todos estos nuevos desplazamientos y derivas encuadradas en las vecindades infinitesimales y encajes colindantes en $\Delta$. Bajo estos lineamientos analizados, el corolario desprendible formula: el evaluó del compendio $P(\mathcal{F}(t), \mathcal{G}(t)) = P(\phi, \psi)_{\gamma(t)} = \mathbb{H}_{\gamma(t)} = 0_{\gamma(t)}$ persevera en la inanición bajo toda la comarca temporal rodeante al estatus $t_0$. En retrospectiva se asienta formalmente el cimiento y descubrimiento imperativo: Para un integrante interior del dominio se pudo forjar su envoltura de vecindad y abierta acogedora encapsulada y sumergida perimetralmente en el enclave, asegurando un hábitat topológico dotado y bendecido con el axioma indispensable inherente a la categorización de Topología Abierta.

\textbf{Fase Final de Clausura Condicional y Deducción Global:} Para aseverar y sellar en su totalidad la cualidad inamovible hermética de carácter cerrado requerida, procedemos a discurrir apoyándonos en la inducción continua de la función y operador combinado. La cartografía funcional $t \mapsto P(\mathcal{F}(t), \mathcal{G}(t))$, como integrante asimilado intrínseco del haz compuesto universal $\mathcal{O}$, ostenta linaje continuo sobre el paramétrico. La evaluación temporal proyecta su rúbrica al centro de anclaje base dictaminando que $\gamma(t)$ es función contigua real de variable temporal $t$. Consiguientemente el mapeo pre-estipulado es cerrado invariablemente por el límite funcional de igualaciones continuas y la pre-imagen clausurada limitante temporal, derivándose certeramente y proveyendo un estatus validado de atributo Cerrado.
Acumulando la dualidad de atributos paradójicos simultáneos de Abierto-Cerrado que le embarga al recoveco y set $S$ inserto en la ininterrumpida y unida topología del sendero rectilíneo $[0, 1]$ que enaltece el arquetipo primordial de estructura de Conexidad Total Universal, nos vemos orillados a proclamar el veredicto terminal: la dualidad paradójica se exime solo abarcando la generalidad universal total $S = [0, 1]$. La ley funcional enmarca indómita su perpetuidad y el dictamen global certifica impertérrito y sin resquicios la permanencia atávica de las interrelaciones bajo el fenómeno general transicional continuista expansivo.
\end{proof}

---

\subsection{2.3. Espacios recubridores e integración a lo largo de curvas}

\textbf{El vínculo entre el tejido geométrico recubridor topológico y el axioma de cálculo analítico e integrativo:}
En lo que sigue, el horizonte y el periplo investigativo se orientan a desentramar, explayar y elucidar fehacientemente una simbiosis inusitada y profunda subyacente que amalgama simbióticamente dos esferas diametralmente dispares pero complementarias: el fenómeno de ascenso topológico puro en escalamientos guiados que subyuga el levantamiento cinemático de curvas, conjurado con la doctrina y operatividad rigurosa que enmarca los dominios canónicos del análisis infinitesimal para la cuantificación integrativa e integración compleja de cauces holomorfistas y funcionales a lo largo del sendero terrenal demarcador curvo en recintos.
La exégesis que vertebra esta correspondencia es que si postula que una función evaluativa emana y ostenta el pedigrí y estatus de pre-descendencia, constituyéndose de forma prístina como una originación primitiva o matriz local, formalizada universalmente bajo el canon $f=G'$ para cierta función progenitora $G$, la labor de recolectar integralmente los valores infinitesimales a lo extenso del vector perimetral curvo $\gamma$ se desmitifica y adquiere una transmutación dual paralela.
De facto, la compilación integrativa de $f$ es isomorfa, paralela y rigurosamente inter-traducible al fenómeno procedimental y secuencial puro de realizar un levantamiento, exento de integrales explícitas, de la estela curva sobreyugante bajo la acción subyugadora de una proyección o amparo topológico peculiar; a saber, un operador recubridor que amalgama y extrae directamente el pulso y ratio infinitesimal subyacente asimilado como un haz y operador derivador universal intrínseco formalizado matemáticamente por la entidad recubridora inmersiva designada como $d: \mathcal{O} \to \mathcal{O}$.
Es mandatario, para aclimatar y ambientar el constructo general con bagaje fenomenológico preexistente asimilar que, en los albores, diversos entes de ramificaciones complejas y prolongaciones funcionales multiformes e injeridas, obran en simetría con recubridores primigenios ilustrados en la galería clásica de la doctrina analítica integral, ejemplificados pormenorizadamente en los pasajes descriptivos consecutivos siguientes.

\begin{ejemplo}[Arquitectura Recubridora Universal Exponencial Paramétrica $\exp: \mathbb{C} \to \mathbb{C}^*$]
Demostramos y atestiguamos solemnemente el linaje legítimo y categórico y el atributo pre-otorgado que ratifica que la función general exponencial analítica $\exp(z) = e^z$ se alza como una superestructura recubridora por excelencia y antonomasia, cubriendo un infinito de sábanas enrollables adscritas y sobre-yugantes a un dominio perforado del plano o centro estelar excluyente, denotado $\mathbb{C}^* = \mathbb{C} \setminus \{0\}$.
Se formaliza el enrutamiento metodológico constructivista: Elegimos con perspicacia y fijamos para cualquier anclaje de posición altimétrica y ordenada imaginaria $r \in \mathbb{R}$ un canal perimetral abierto horizontal longitudinal paralelo confinado denotado como franja plana $V = \{ z \in \mathbb{C} : r < \text{Im}(z) < r + 2\pi \}$, el cual innegablemente se consolida y clasifica bajo la rúbrica de ente abiertamente inmerso en $\mathbb{C}$. 
La evaluación y envoltura de esta estela canalizadora horizontal al atravesar los engranajes exponenciales, distorsionadora concéntrica, extrae como imagen mapeada universal todo el espectro anular radiante salvo la barrera o sutura de exclusión generada intrínsecamente. El output se configura explícitamente y formula territorialmente en la planicie inferior como el contorno de exclusión radiado polar: $\exp(V) = \mathbb{C}^* \setminus \{ \rho e^{ir} : \rho \geq 0 \}$, es decir, el plano integral despojado minuciosamente y escindido por el corte raso rectilíneo de desvanecimiento angulado radial provisto en dirección del parámetro. Este escenario o receptáculo codificado es un subespacio categóricamente abierto y garantizado con estatus legal en el interior topológico en $\mathbb{C}^*$.
Procedemos meticulosamente a desmenuzar su antecesora o espectro integral pre-imagen en la estratosfera original, esto es $\exp^{-1}(\exp(V))$. La ciclicidad incesante trigonométrica inserta implícita subyacente en el código y matriz exponencial fuerza indudablemente al desdoble secuenciado repetitivo estratificado universal generatriz, plasmándose irreductible en la segmentación unionista disjunta fragmentada aditiva integral a escala en niveles desglosables formales conformados por láminas y sabanas separadoras $V_n = V + 2\pi i n$ indexadas omnipotentes por los estratos y renglones dictaminados para cada escala entera $n \in \mathbb{Z}$.
Se invoca la escrutación restrictiva analítica selectiva encasillada operante confinando la maquinaria individualmente bajo cada recuadro lámina particular $\exp|_{V_n} : V_n \to \exp(V)$. Este conducto de traspaso se rubrica, enjuicia y corona inapelablemente con el estandarte inmaculado del isomorfismo holomorfo integral totalizante ("biholomorfismo" absoluto), avalado firmemente tras la superación exhaustiva de las pruebas del decálogo de legitimidad para biyectividades complejas analíticas: 
- Fase Subyacente e Integral de la Sobreyectividad: Comprobada sin paliativos ante la copiosa y generosa inundación del codominio provista irrefutablemente en virtud de la densificación ubicua universal logarítmica y fase de argumentos continuos provistos.
- Fase Crítica e Inflexible Analítica de la Inyectividad Intrínsica: Concedamos de buena fe y postulemos bajo conato hipotético una amalgama equitativa forzada $e^{z_1} = e^{z_2}$ para dos inquilinos concurrentes $z_1, z_2$ oriundos de la misma celda encasillada de la sábana estipulada singularmente enmarcada en $V_n$. La estricta igualdad trigonométrica dictamina un requerimiento perentorio inalterable imponiendo una rectificación puramente imaginaria discreta escalar y desfasadora transicional de la diferencia estipulada inamoviblemente en la pauta algorítmica: $z_1 - z_2 = 2\pi i k$ con algún parámetro obligatoriamente indexado de pureza integra real entera $k \in \mathbb{Z}$. No obstante, constreñidos bajo la ley marcial de habitabilidad interna dispuesta inflexible para la contención simultánea al interior de $V_n$, su altitud imaginaria de las ordenadas y separación de ambos entes se subyuga a una frontera represiva forzosa máxima constreñida acotadamente bajo $|\text{Im}(z_1) - \text{Im}(z_2)| < 2\pi$. La confluencia indisoluble e intercesión simultánea de estas directrices condena inexorablemente a subyugarse en el umbral singular donde el escalón transicional no opera salvo nulidad, forzando implacablemente a claudicar el parámetro asumiendo el encierro nulo $k = 0$, cristalizándose inapelablemente e irrefutablemente el cimiento y conclusión inquebrantable subyacente que atestigua una identidad gemelar absoluta irrompible resolutoria $z_1 = z_2$.
- Fase de Inversión Analítica y Verificación Operadora del Inverso Analítico Diferencial: Exigido y consolidado por virtud y laudo supremo conferido bajo los estatutos normativos regios del dictamen irrenunciable Teorema Fundamental Diferenciador de la Aplicación y Transformación Inversa e Invertibilidad Holomorfa Analítica Universal de índole topológico-complejo. Dado que la operación de transferencia derivada primigenia se evalúa y no decrece en inanidad en el curso de todo el escenario territorial ostentando con contundencia y aplomo la condición inextinguible de $\exp'(z) = e^z \neq 0$ siempre, la exención de encallamientos periciales singulares queda anulada y blindada contra fallas estructurales singulares.

Se cristaliza este entramado teórico materializado mediante su visualización diagramada cartográfica y proyección a escala fidedigna delineada topológicamente:
\begin{center}
\begin{tikzpicture}[scale=1.2]
    % Dibujo de franjas y hojas superpuestas con cortes guiados
    \draw (-2, -1) -- (2, -1) node[right] {$V_{-1}$};
    \draw (-2, 0) -- (2, 0) node[right] {$V_0$};
    \draw (-2, 1) -- (2, 1);
    \draw[dashed] (-2, 2) -- (2, 2) node[right] {$V_1$};
    
    % Flecha del operador
    \draw[->] (2.5, 0.5) -- (3.5, 0.5) node[midway, above] {$\exp$};
    
    % Dibujo del codominio proyectado plano perforado estelar
    \draw[->] (4, 0) -- (7, 0);
    \draw[->] (5.5, -1.5) -- (5.5, 1.5);
    \draw[thick, red] (5.5, 0) -- (7, 0) node[right] {corte radial excluido};
\end{tikzpicture}
\end{center}
\end{ejemplo}

\begin{ejemplo}[La Dinámica del Recubrimiento Polinomial Fraccionado Paramétrico de las Potencias Algébricas, $p_n: \mathbb{C}^* \to \mathbb{C}^*$ operante en $z \mapsto z^n$]
Sometemos bajo el tamiz y prisma riguroso de escrutinio estructural la evidencia fidedigna contundente que corona la inducción y proyección analítica iterativa de potencia estipulada como $p_n(z) = z^n$ enmarcándola insoslayablemente bajo el rango e índole categórico calificado de recubrimiento holomorfo y ramificado constituido invariablemente por un despliegue de multi-coberturas apiladas equivalente en su recuento cardinar a las $n$ láminas o sábanas conformadas superpuestas solapadas en helicoide sin centroide singular.
Al amparo de un eje arbitrario radiante selectivo actuando de frontera supresora estipulado geométricamente como haz polar excluyente $\text{rayo} = \{ R e^{i\alpha} : R > 0 \}$, se disecciona retrospectivamente el análisis topológico en reversa rastreando su pre-imagen inversa y retroactiva englobada territorialmente conformante en el constructo de evaluación y acopio de preimágenes formadas a expensas bajo el comando $p_n^{-1}(\mathbb{C}^* \setminus \text{rayo})$. La recolección analítica desgrana y despliega invariablemente con exactitud cristalina un abanico y amalgama separada desvinculante de componentes fragmentadas disjuntas aglutinadas conformando rigurosamente un portafolio encasillado distribuido por $n$ porciones seccionales sectoriales polar-anguladas puras encasilladas individualmente denominadas perimetralmente $S_k$ (para indexación paramétrica abarcante iterante en los umbrales de escala desde $k=0$ hasta su frontera final límite $n-1$). En cada una y recuadro de estas cuñas sectorizadas aislantes $S_k$, el trazador y fluctuación angular perimetral y evaluador del desvío imaginario se congrega encriptado con exactitud constreñida acotado oscilatoriamente transcurriendo inexorablemente a través y por encima del puente limitante abarcando intervalo paramétrico angulado provisto de amplitud constante invariable de valor máximo acotado angular tasado infalible en $2\pi/n$.

Se invoca peritaje e indagación restrictiva circunscribiendo fronteras en una cámara cuña seccional única aislando una encasillada operante unitaria $p_n|_{S_k} : S_k \to \mathbb{C}^* \setminus \text{rayo}$. Procedemos a refrendar con ahínco infalible el laudo pericial que legitima su condición biyectiva absoluta irrefutable: la presunción contenciosa de que dos embajadores concurrentes en el mismo escenario territorial cuña $z_1, z_2 \in S_k$ se mimetizarían al aplicar la ley operativa proyectando igualación de la fuerza proyectiva dictada $z_1^n = z_2^n$ promueve irrevocablemente una coacción resolutiva de cociente derivativa equivalente forzosa constreñida con amalgama $(z_1/z_2)^n = 1$. Este evento desencadenante se dirime acotadamente acogiéndose al refugio y al amparo excluyente del catálogo resolutorio cíclico de raíces y raigambres complejas integradoras de la unidad cíclica, ergo estipulando transicionalmente la inmersión del factor $z_1/z_2 = e^{2\pi i j / n}$ portando un escalar de calibración índice variable inserto restrictivamente enclavado como  $j \in \{0, 1, \dots, n-1\}$. Empero, si propugnamos lógicamente avalar un desliz de desviación y alteración adoptando $j \neq 0$, este disenso o corrimiento en el engranaje arrojaría inexcusablemente un desfase paramétrico rotacional rotundo en el cociente y consecuentemente atestiguaría una apertura angular abismal desfasada tajante entre los orígenes evaluada y computada diferencialmente como $2\pi j / n$. Conducidos empero, a la contradicción inquebrantable de que el enclaustramiento regulatorio topográfico impone severamente el hospedaje concurrente simultáneo imperativo de ambas raíces hospedadas dentro de las angostas y confinantes murallas de contención abarcadas anguladas en sector único con tolerancia extrema reducida máxima permisiva fijada a $2\pi/n$. La colisión irreconciliable estalla dictaminando imposible el espaciamiento distanciado desfasado. En virtud a las pruebas y a la coacción inexcusable, el dictamen subyuga colapsar el índice resolutorio disidente al nulo apaciguante dictando asilo imperativo infalible $j = 0$, atestiguando irreprochablemente que el nudo transaccional arroja unidad y garantizando inviolable preeminencia e intercesión sellando la prueba inyectiva $z_1 = z_2$.
La fluidez sin estridencias holomorfa de este canal de transferencia fluye sin menoscabo o colapso singular oclusivo por cuenta y beneplácito de la operatividad pericial derivadora operante, a saber la inmersión de diferencial explícita manifestada $\frac{d}{dz}(z^n) = n z^{n-1}$, cuyo vector de potencia carece inobjetable de inanición anuladora ($n z^{n-1} \neq 0$) merced a la escisión estelar primordial y blindaje del entorno excluyente abarcador base $\mathbb{C}^*$, otorgando en plenitud universal el título conferido de transformación holomorfista biyectiva e Isomorfismo completo (Biholomorfismo) para esta cuña segmentaria.

El compendio esquemático diagramático a escala de las conformaciones topológicas secciona y exhibe el caso modelizado para paramétrico ternario o tridente seccionador $n=3$, exponiendo los vectores y linderos fronterizos sectoriales delimitados estelares:
\begin{center}
\begin{tikzpicture}[scale=1]
    % Sistema de base coordenadas
    \draw[->] (-2,0) -- (2,0); 
    \draw[->] (0,-2) -- (0,2);
    
    % Delimitación fronteriza perimetral delimitadora con vectores guías sectoriales acotadores para n=3 (tercios cíclicos 120 grados)
    \draw[thick, blue] (0,0) -- (1.5, 2.6); 
    \draw[thick, blue] (0,0) -- (-2.6, 0); 
    \draw[thick, blue] (0,0) -- (1.5, -2.6); 
    
    % Anotación nominal indicativa de los recintos y cuarteles sectoriales habilitados y segregados
    \node at (0.8, 1) {$S_0$};
    \node at (-1.5, 1) {$S_1$};
    \node at (0, -1.5) {$S_2$};
\end{tikzpicture}
\end{center}
\end{ejemplo}

\begin{ejemplo}[Caso disidente o anómalo de Fallo Topográfico de Revestimiento: Proyección Evaluativa y Mapeo Central del anclaje y punto focal basal de los gérmenes de Funciones del espacio total $\mathcal{O}$]
Invertimos la lente investigativa para postular, escudriñar y atestiguar y rubricar un fenómeno de quiebre de simetría fundamental que certifica irrevocablemente que, a despecho de exhibir con excelsa fidelidad características de localismos homeomorfos conformacionales locales continuos, la majestuosa proyección cenital y eje focal del ente subyacente $\mathcal{O}$ hacia su lecho terrenal $\mathbb{C}$ dictaminado como la operación cartográfica universal $p: \mathcal{O} \to \mathbb{C}$ sucumbe abismal e irremediablemente, mancillando y deponiendo su corona e impidiendo con rigor extremo adjudicarse, bajo ninguna indulgencia o subterfugio analítico, el titulo noble o condecoración preclara de encumbrarse bajo la égida o designio supremo de ostentar y ser proclamado Aplicación Global de Recubrimiento.
Conjuramos al estrado acusatorio como prueba incriminatoria tangible a la vía dolorosa descrita en el plano cartesiano base encarnada e instrumentada en la directriz del segmento curvo vector rectilíneo y de tránsito decadente dictaminada temporal e irremediablemente decreciendo y agonizando perimetralmente en caída y colisión guiada inexorable hacia los infiernos del vórtice vacío, formalizada y descrita rigurosamente bajo la ecuación directriz horaria transitoria $\gamma(t) = 1 - t$, constreñida acotadamente en los marcos temporales entre su génesis y origen incipiente de fulgor matutino en el ápice de la marca y partida e inicio inaugural evaluado y estipulado como el faro $\gamma(0) = 1$, hasta la clausura y destino perimétrico estelar apocalíptico que marca su colapso crepuscular y mortandad colisional final enraizado fatalmente en el hundimiento en $\gamma(1) = 0$.
Alistamos en la línea de salida del espacio superior estratificado y sobreyugante a nuestro campeón y germen de inicio designado y erigido como el portaestandarte $f_1 = (1/z)_1$, estandarte y representante germinal de la emblemática inversión holomorfa singular.
Daremos curso inicial avalando e instruyendo el conato de falsación, acunando de manera transitoria la presunción y dogma de la elucubración vana de existencia y factibilidad del requerimiento condicionado postulado exigente bajo un levantamiento continuo global y unificador, bautizado paramétricamente con las siglas sobreyugantes superpuestas $\widetilde{\gamma}: [0, 1] \to \mathcal{O}$, el cual encarna la supuesta ascensión al edén con la obligación dictaminada al alba matutina inquebrantable de sellar el amarre basal e inicio correlacionado $\widetilde{\gamma}(0) = f_1$. Denominaremos y abreviaremos pragmáticamente esta caravana fantasmal levantada en ruta ascensional denotándola nominalmente bajo un estandarte de macro-entidad transicional evolutiva temporal bautizado de acá en delante y acuñado irrevocablemente en toda instancia y foro de registro como paramétricamente $\mathcal{F}(t) = \widetilde{\gamma}(t)$.
Concitamos y traemos al escenario a la par un ente gemelo solidario continuista, evaluador, observante y referenciador perenne incesante imbuido en simetría inalterada provista trivial y benignamente al arrastre sin fisuras de la identidad analítica de prolongamiento y extensión natural continua germinal emanada sin conflicto e injerencia natural de estirpe directa descendiente por la matriz y progenitora holomorfa $g(z) = z$ que fluye sin traumas a cualquier latitud de anclaje, formalizada en su tránsito continuista por la trayectoria designada en sintonía horaria paramétrica análoga y homóloga descrita en cada estación como $\mathcal{G}(t) = g_{\gamma(t)}$.
Ensamblamos el cruce y convergencia cruzada temporal relacional enraizada y forjada inquebrantablemente al despuntar el alba, es decir en las fronteras originarias en $t=0$, escrutando pormenorizadamente los designios evaluados de confluencia germinal multiplicativa inicial de arranque original de las fuerzas: $\mathcal{F}(0) \cdot \mathcal{G}(0) = f_1 \cdot g_1 = (1/z)_1 \cdot (z)_1 = (1/z \cdot z)_1 = (1)_1 = 1_1$. Se sella imperturbable el vínculo perpetuo e imperecedero forjado por la Ley Relacional Continua de Multiplicación e Igualación Constante, consolidada indudablemente por el estatus y postulado consagrado dictaminado inquebrantable y blindado blindadamente por el resguardo inviolable en toda latitud otorgado y avalado explícitamente en el corpus o canon doctrinal del Principio de la Permanencia de Relaciones Funcionales Estipuladas (Proposición 2.2.4 y sus subyugaciones). El edicto derivativo nos fuerza irreversible e irrecusablemente a trasladar a toda hora e instante sin resquicio eximente en toda la senda el sello e imposición condicional en todos los confines paramétricos abarcados e ininterrumpidos y estipulados y encadenados temporalmente dictaminando ciegamente la inmolación en: $\mathcal{F}(t) \cdot \mathcal{G}(t) = 1_{\gamma(t)}$ de manera ininterrumpida y persistente y obligante universal para la entera globalidad temporal incesante y exhaustiva transcurrida a través de la duración de $t$.
La marcha de la ejecución prosigue e inquirimos escrutinio final en el desenlace colisional postrimero, escudriñando minuciosamente e interpelando exigente en la postrimería fúnebre ocaso límite de mortandad y cierre abismal límite del cese temporal final anclado en destierro en la frontera final evaluada y computada en $t=1$. El recuento e impacto del evento produce el desenlace analítico y la ecuación germinal resultante terminal se expone a sí misma forzada implacablemente evaluando y cristalizando el remanente en $\mathcal{F}(1) \cdot \mathcal{G}(1) = \mathcal{F}(1) \cdot (z)_0 = \mathcal{F}(1) \cdot g_0 = 1_0$.
Para certificar la sentencia exhumamos e inferimos rigurosamente evaluando perimétrica y operativamente al interior de los corazones funcionales representados operativamente en los gérmenes de esta catástrofe en el seno mismo del núcleo epicentral de falla y de anclaje concéntrico desastroso $z=0$, colapsando y obteniendo al evaluar en forma directa numérica en el núcleo de anclaje cero un disparate matemático desgarrador insostenible desmoronador de los andamios de la aritmética axiomática del álgebra de cuerpos: $\mathcal{F}(1)(0) \cdot g(0) = \mathcal{F}(1)(0) \cdot 0 = 1$. De esta abominación colisionante desprendemos sin dilación, de la irremediable nulificación absorbente, el advenimiento espurio aberrante en $0 = 1$, el cual erige irremisible un muro infranqueable inquebrantable contrapuesto insostenible absurdo formal resolutorio irreconciliable de base e impasable a la razón e intelecto y sentido común deductivo lógico algebraico o axiomático.
Apresados sin refugio salvador por el dictamen letal de esta ineludible y profunda y terminante refutación y contradicción abrumadora colosal destructiva, se expone descalzo y al desnudo sin ropaje encubridor la vana e ingenua tentativa primigenia hipotética de levantar asilo al conato quimérico y onírico inasequible de pervivencia global transicional ascensional. El dictamen sentencia la negación taxativa absoluta inapelable del postulado asertivo del axioma fundacional propugnador del atributo topológico de poseer y disponer y ostentar y esgrimir y garantizar irrevocablemente con suficiencia para todo rastro curvilíneo sin sesgo la sublime inigualable preclara, irremplazable Propiedad Inquebrantable Infalible Aseguradora del Levantamiento Integrador Ascensional Global Continuo y Sistemático Invariable Total de las Curvas (cimentada rigurosamente probada inhabilitante por vía indirecta a las espaldas conculcando severamente el cimiento de los atributos enumerados restrictivamente vinculantes indisolubles forzosos en la Observación 2.1.3 y correlacionado dictamen en la Proposición 2.1.1). Es así, consecuentemente e irrevocablemente sentenciado como laudo condenatorio postrimero que en estricto rigor innegociable e implacable topológico, al evidenciar contundentemente la desposesión invalidante de esta vital y neurálgica piedra angular axiomática garantizadora, que la directriz natural y proyección cartográfica $p$ es depuesta y despojada ignominiosamente del fuero y condición para aspirar al blasón de constituir una genuina o lícita y avalada integral y auténtica e inalterada Aplicación de Recubrimiento.

El croquis o bosquejo esquemático ilustra y delata de manera punzante y contundente, graficando perentorio y de facto en proyección terrenal asfixiante descendente al pozo nulo abismal del núcleo fatal estelar la ruta decadente de la traza temporal maldita abocada e inmersa sin salvamento e inaplazable al despeñadero insalvable resquebrajador del quiebre y ruptura asimétrica fatal colisional descarrilando en $0$:
\begin{center}
\begin{tikzpicture}[scale=1.5]
    % Bases y marco posicional estructural delimitante de soporte
    \draw[->] (-1,0) -- (2,0);
    \draw[->] (0,-1) -- (0,1);
    
    % Directriz trazadora fatal mortuoria del abismo rectilíneo convergente hacia colapso en estelar cero y fin de la existencia germinal finita sin singularidades
    \draw[thick, red, ->] (1,0) -- (0.1, 0);
    
    % Demarcación y fijación del epicentro matriz original estelar luminoso y el vórtice final colisional atrayente mortal aniquilador resolutorio de quiebre
    \fill (1,0) circle (1.5pt) node[above] {$1 = \gamma(0)$};
    \fill (0,0) circle (1.5pt) node[below right] {$0$};
\end{tikzpicture}
\end{center}
\end{ejemplo}

\begin{definicion}[El magno y preclaro Operador Diferencial Derivativo y analítico $d: \mathcal{O} \to \mathcal{O}$]
Se instituye rigurosa y procedimentalmente en un acto de imposición definitorio formal y solemne el nombramiento del supra-operador generatriz holomorfo derivador e inter-germinal en las altas esferas elevadas encumbradas sobreyugantes del estrato de gérmenes analíticos y funcionales integrales y holomórficos operando sobre los residentes del espacio topológico base del Haz $\mathcal{O}$ estipulado axiomáticamente prescribiendo a cada germen holomorfo genérico y aseverando la correspondencia asignada evaluadora descrita por los preceptos: 
\begin{equation}
d(f_a) = (f')_a.
\end{equation}
Preconizamos y desvelamos que esta estipulación e imposición normativa demanda escrutinio de rigurosidad imperativa corroborando fidedigna y estrictamente si el decreto legislador formal propugnado se afianza en cimiento sólido validado axiomáticamente superando incólume el test y certificación incondicional de una legítima y asertiva Buena Definición analítica estructural operativa invariante e inmune independiente blindada ante cualquier alteridad impuesta por cambio, permutación o reelección exógena y caprichosa de representantes sustitutivos germinales inter-equivalentes homólogos. Sometemos el constructo derivativo a la criba y avalamos que de mediar, concurrir o converger un representante alterno suplente paralelo e interino equivalente avalado con un mandato de identidad subyugante $(U, f) \sim (V, g)$, aflora entonces indubitable inexcusable el dictamen o laudo restrictivo amparador que promulga explícitamente e irrenunciablemente como consecuencia analítica inexorable inmediata que subyace forzosa una colindancia geográfica o enjambre abierto de encuentro unificado común o entorno fronterizo vecinal acotado concéntrico restrictivo interceptado de validación estipulado temporalmente como $W$ encastrado al abrigo anidado e inmerso interiormente encubierto bajo las paredes aluviales y confluencia delimitada al amparo topográfico amurallado de $U \cap V$ donde la identidad sin roturas funcionales y la concordancia analítica original rige dogmática obligante inquebrantable impoluta igualando $f = g$ de forma tajante imperativa irrestricta en la globalidad integradora territorial de $W$. Desatando la cascada y vertiente algorítmica y operatoria analítica en cascada encauzando el vector derivador consecutivo y consecuencial analítico infinitesimal irrenunciable sobre la faz integral analítica evaluada derivando a ambos flancos de la igualdad en estricta conformidad asimétrica y sincronía, el laudo sentencia sin resquicios e inapelable que $f' = g'$ ejerce tiranía idéntica operando análogamente congruente irrestricta totalizante sin sesgo universal en la totalidad abarcada por el mismo reducto y vecindario abierto asilo limítrofe coincidente $W$. Bajo el rigor escrupuloso imperativo de las escrituras de equivalencia inter-germinales se ratifica estelarmente dictaminando y homologando a cabalidad inmaculada cristalina y unificadora que en el seno epicentral anclado en $a$ de forma local germinal estelar coincidente asienten conformando $(f')_a = (g')_a$.
Se atestigua, promulga explícito para los legajos y previene documentar formalmente colateral aclaratorio suplementario insigne que el ente rector actuante magno referenciado universalmente como mandato u operador derivador transicional abstracto nominado y distinguido supra por las siglas $d$, circunscribe su influjo o radio dictatorial intervencionista exclusivamente operando incesante y con exclusividad reclusa interina albergando acotado bajo y encajado forzosamente encasillado en un canal y sendero limitante dictando encauzando o prescribiendo estrictamente sin desviaciones exógenas enrutando la migración interna recluida herméticamente puramente como una función o aplicación de travesía en tránsito inmanente intrínseca encriptada puramente fluyendo de tallo germinal madre subyacente reclusa a sí misma estancando o constreñida análoga internamente acotada recayendo evaluativa al seno y entrañas del mismo tallo anclado, denotando de forma estricta limitante en encerramiento el flujo operante inyectivo confinado asimétrico dictando invariabilidad de estrato basal: $d(\mathcal{O}_a) \subset \mathcal{O}_a$. Todo lo cual corrobora contundente a gran escala y se transfiere avalando inequívocamente y garantizando categórica y explícitamente sin dubitación o ambigüedad preestablecida que la mano del operador inmanente exime de alteración en forma incólume e inviolable a la posición o eje base geográfica basal inhabilitando tajante imposibilitando categóricamente y excluyendo despojar y erradicar y extirpar inhabilitando dislocar y mover o agitar de su asilo al epicentro o núcleo radiante de anclaje territorial basal nativo.
\end{definicion}

\begin{proposicion}[El Magno Operador Diferencial Derivativo $d: \mathcal{O} \to \mathcal{O}$ exhibe atributos innegables y constituye categórica y asertivamente un firme Revestimiento Universal y preclara de Aplicación Global estricta y genuina de Recubrimiento Topológico y Funcional Analítico Integral Absoluto]
\end{proposicion}
\begin{proof}
Para refrendar la contundencia de esta máxima procedemos sin dilación en la incursión analítica encastrando los argumentos lógicos deductivos en un acueducto secuenciado y acorazado e inquebrantable argumentativo dividido en compartimientos lógicos y topológicos disjuntos y sucesivos integrales a fin de conformar incólume el castillo pericial demostrativo asertivo probatorio y garantizador infalible resolutorio.

\textbf{Bastión 1 (El Acopio Topográfico Estelar Cartográfico Geográfico Generador Abierto de las Láminas Disjuntas y Sabanas Primitivas Analíticas encriptadas Matriciales Inmersivas Integrales Locales en Bases):} 
Postulemos inmersivo de base y reclutemos en evaluación y escrutinio genérico y exhaustivo cualquier germen general arbitrario incondicionado e insigne representante $f_a \in \mathcal{O}$ enraizado originario local basal con linaje avalado estelarmente con credenciales representativas acreditadas avalando preexistencia original y portador matricial avalado dotado originario referenciado bajo los emblemas amparadores del par acoplado $(D, f)$ erigido sobre el cimiento territorial cartográfico de disco cóncavo y asilo topológico y abierto asertivo y amurallado circular puro conexo estelar concéntrico, convexamente engarzado estipulado perimetralmente asumiendo como el receptáculo amparante un disco puro euclidiano $\Delta = D$ acordonando y acorazando centrado y nucleado estelar y diametralmente protegiendo en $a$. Arrastrados, propulsados y arropados irrenunciablemente por las insignes prerrogativas e imponentes estamentos del Magno Teorema Integrativo Fundamental Analítico Funcional Resolutorio Universal de Cauchy estipulado para dominios conexos asilos circulares y estrellados (Thm 1.3.6), dictaminamos y extraemos exhumando de las sentencias del edicto analítico integrador de primitivas absolutas con perentoria exigencia y dictamen de preexistencia resolutoria inmaculada asegurando con rigor ineludible dogmático incontrovertible de certidumbre la manifestación real latente o florecimiento y germinación constituyente inexorable imperativa forzosa y la materialidad y posesión indiscutible fidedigna irrebatible e incontestable sin sombra o atisbo o asomo quimérico o de incertidumbre de la materialización y existencia subyacente certera ineludible constatable en la topografía circular amparante del Disco puro $D$ de una y matriz generadora originaria e inmaculada encumbrada e inmarcesible noble y estelar Primigenia Funcional Primitiva Base Analítica Integral Superior Matriz o Supra-función Madre Analítica Universal $\mathbb{F} \in \mathcal{H}(D)$ dictaminando incesante e invariable obediencia resolutoria subyugante y concordante de equivalencia operativa derivativa infinitesimal a nivel basal dictaminando imperante resolutoria y obligatoria obediencia forzosa y sometimiento riguroso evaluativo en todos los enclaves del asilo al dictado inquebrantable de estricta equiparación derivativa $\mathbb{F}' = f$.
A partir de este basamento originario matricial delinearemos e instauraremos sin reservas en torno y con epicentro irradiante de encauzamiento del estelar inicial nativo germinal irradiante originario aludido en principio inicial e investido por $f_a$ un constructo amparador topológico vecinal fronterizo de asilo de cuna o envolvente y matriz cobertora básico abierto asertivo tubular de cimiento natural tubular abierto y estelar dictaminado $V = \mathcal{N}(D, f)$. El devenir demostrativo resolutorio inexcusable ineludible y el encargo fáctico imperativo deductivo subsecuencial perentorio se acota y confina estricto al objetivo deductivo supremo exigiéndose y abocándose denodadamente a inspeccionar minuciosamente desgajar, catalogar encriptar e inventariar en profundidad las entrañas absolutas retroactivas inmanentes analíticas conformando el ente reverso resolutivo o matriz integral espectral retroactiva o sombra y lecho base origen pre-proyectado de la matriz inversa inversa de imagen y acopio general de procedencias englobadas analíticas $\mathcal{W} = d^{-1}(\mathcal{N}(D, f))$ develando la red y matriz del engranaje basal lámina a lámina subyacente latente intrínseco.

\textbf{Bastión 2 (La Deconstrucción Espectral y Despliegue Cartográfico del Desglose Estelar Resolutorio Matricial Cartográfico Analítico Laminar Abierto Exhaustivo Indisoluble Integral Absoluto Cartográfico):} 
Esclarecemos las tinieblas de los anales retroactivos invocando y enumerando a cada candidato y concurrente recluso espectral habitante integral integrante oriundo recluso y subyacente y encasillado albergado y empadronado ineludiblemente bajo la cubierta abarcativa omnímoda receptáculo asilo del espectro pre-imagen reverso resolutivo o matriz origen $\mathcal{W}$. Este asilo amparador cobija irremisible a los inquilinos germinales originarios avalados estelarmente con credencial genérica subyacente de la forma general $G_z$ bajo la tutela y condición inquebrantable acotada evaluando e imponiéndose forzosa constrictivamente bajo el cimiento y lecho que demanda amparado en axioma derivativo: que el evaluo posicional del núcleo caiga o encaje de asiento terrenal cartográfico dentro de las murallas en el dominio interior discoidal puro $\Delta=D$; y que por ende emane correlacionado y anclado avalado simultáneo concurrente y aparejado ineludible avalando la posesión y dotación fidedigna de rango holomorfista matricial originario e inmaculado primigenio avalado encasillado dictaminado con título y representación legal bajo una macro matriz envolvente madre o supra-generadora $\mathbb{G} \in \mathcal{H}(D)$ estipulándose rigurosamente bajo fianza y sumisión de obediencia que el dictamen operatorio diferencial evaluador devuelva integro restitutivo ineludible reproductor del aval original igualativo diferencial de rigor $\mathbb{G}' = f$ operando en completitud exhaustiva de la cobertura universal asilar en el claustro del asilo cuna o disco madre $D$.
Empero, bajo el faro luminoso irradiador omnímodo esclarecedor purista purificante y analítico exhaustivo y depurativo del axioma derivativo primigenio y lema analítico diferenciador unificador del dogma diferencial, atestiguamos solemne refrendar indiscutido unánime que para cualquier y toda concurrencia integral evaluativa e inquilino matriarcal genérico de la asamblea originaria subyacente o matriz y engendradora primigenia integral generalizada $\mathbb{G}$ la cual profese inquebrantable el título y ostente y satisfaga con denodado sometimiento la equiparación y correlación obligatoria de estricta igualdad $\mathbb{G}' = f$ en el interior fronterizo confinado abarcador del espacio y claustro madre $D$, su linaje genético analítico estructural matricial dista tangencial asimétrico apartándose resolutoriamente evaluando el disenso únicamente diferenciándose evaluativamente apartando en rango u oscilación varianza a lo máximo asintótico desfasado limitativo escindido dictaminado de forma tajante o disociado excluyente apartando en la estructura evaluadora con desfase a un trecho equivalente medible disociativo absoluto a lo sumo igualable y equiparable estrictamente a una estricta barrera de desfase constrictiva y aditiva estipulada rígidamente por un cimiento evaluador vector y matriz aditivo contante paramétrico enraizado escalar inmanente complejo universal avalado $c \in \mathbb{C}$ inserto y superpuesto frente a la progenie primitiva matricial originaria canónica avalada $\mathbb{F}$. Sentenciando y erigiendo irrecusable dictamen conclusivo estipulado analíticamente e imponiendo a la luz resolutoria: $\mathbb{G} = \mathbb{F} + c$.
Procediendo ininterrumpidamente bajo y consecuencial consecuencialmente bajo estos preceptos de inducción deductivos unificadores de despiece estructural analítico, orquestamos e inducimos y revelamos transfigurante magistral omnipotente e inobjetable el acopio y la condensación abarcadora y unificación resolutoria descriptiva constituyente aglutinadora asilar y sumatoria espectral constitutiva aglomeradora para la cobertura universal total matricial de asilo original espectral matriz y pre-imagen $\mathcal{W}$, ensamblándola rigurosamente a escala estructurada conformada irrefutablemente estipulándola inalienable unificando integrando fusionando como unión masiva y agregación o sumatoria de todos los enjambres vecinales asilares básicos e injertados aglutinados y engarzados disjuntos avalados paramétricamente, expresable formalmente:
\begin{equation}
d^{-1}(\mathcal{N}(D, f)) = \bigcup_{c \in \mathbb{C}} \mathcal{N}(D, \mathbb{F} + c).
\end{equation}
Acometemos de soslayo resuelto denodadamente a justificar probatoria asertiva blindada y blindadamente certificar, validar corroborando sin tacha fisura mácula imperfección debilidad falencia estructural estipulando que esta vastedad agregativa unificadora gigantesca aglomeración ostenta y milita bajo el sello irrenunciable atributo riguroso disociativo inmaculado puro o partición estricta garantizada blindada dictaminando la innegociable condición evaluadora analítica disjunta, es decir que dos estratos o inquilinos paramétricos evaluados encasillados y amurallados vecinales tubulares y encastrados no solapan ni sufren roces intercepciones o escaramuzas superposicionales intromisiones o intrusiones, ergo formulado de forma asertiva purista: $\mathcal{N}(D, \mathbb{F} + c_1) \cap \mathcal{N}(D, \mathbb{F} + c_2) = \emptyset$ asumiendo condicionado con disenso paramétrico e impedimento restrictivo ineludible $c_1 \neq c_2$. Para apuntalar este aserto procedemos bajo refutación deductiva o reducción probatoria contenciosa asumiendo falsedad hipotética y presumiendo onírica invasiva y espuria la intercesión y confluencia superposicional o intromisión delatora evaluadora originando inquilino residente doble empadronamiento usurpador un elemento intruso $G_z$ escurridizo cobijado infame escurriéndose insertado clandestino subrepticiamente colado en el ceno y abrigo o recoveco interno del interior de ambas circunscripciones. Este evento y desacato conminatorio y catástrofe forzaría e impelería indiscutible de ser así irrefutable y coercitivamente por decreto supremo inquebrantable de inducción e identidad a que los entes representantes funcionales progenitores macro-estructurales evaluativos matrices y estelares $\mathbb{F} + c_1$ y su contrapuesto gemelo paramétrico espurio evaluativo encasillado gemelar y simétrico contencioso matricial $\mathbb{F} + c_2$ sufrieran y padecieran imperioso un hermanamiento o equiparación resolutoria igualación evaluativa asimiladora funcional mimetizante sin reservas idéntica o concordancia solapada idéntica co-operante en y operando encasillada local coincidente encastrada acotadamente y amparada sobre una circunscripción, vecindad o enclave topográfico acordonado amurallado o perímetro concéntrico abierto puro envolvente acotado resguardado y blindado adosado colindante y periférico ceñido contiguo al entorno epicentro vecinal circunscrito anclado y acotado concéntrico a las inmediaciones estelares vecinales limítrofes nucleares en torno de $z$. Invocamos imperturbables nuevamente al dogma pilar y cimiento irrebatible y majestuoso dogmático asertivo del supremo y regio y unificador rector analítico o \textbf{Principio irrenunciable e Infalible rector y ordenador integrador y cimiento de Identidad de extensión Analítica Holomorfa y Funcional} el cual al actuar operante infalible subyugante y omnímodo dominador rector en la integridad conexidad del escenario y recoveco asilar amparante del Disco puro circular amparante madre y asilo $\Delta=D$, elucubraría y precipitaría desencadenando la hecatombe deductiva de colapso de rangos decretando avalando y forzando y ratificando la extensión abarcadora resolutoria global de forma idéntica general impuesta idéntica obligatoria coercitiva universal que abarca unificando imponiendo ineludible igualación evaluativa asimiladora y mimetismo en toda la geografía extendida subyugando general y total: atestiguando o condenando a ser evaluadas por el dictamen irrevocablemente a sufrir paridad universal unificadora impuesta dictando forzosa y general: $\mathbb{F} + c_1 \equiv \mathbb{F} + c_2$ sin salvedades y en su totalidad geográfica para la integridad cartográfica del Disco $D$. Esta igualación o emparejamiento resolutivo a gran escala unificador desata colateral la precipitación resolutoria asimétrica disyuntiva algebraica pura trivial desvelando escindida depurada que subyace de origen imponiendo y destilando la falencia insostenible irreconciliable $c_1 - c_2 \equiv 0$, colapsando y quebrando insostenible y vulnerando en contradicción manifiesta evidente abrumadora el blindaje inicial y la prohibición fundacional axioma paramétrica rectora prescripta tajante restrictiva prohibitiva inquebrantable $c_1 \neq c_2$, demostrando falsedad categórica y concluyendo asertivamente garantizado inalienable la exclusión recíproca absoluta y por ende resolutoria certificando con garantía total disjunta o pureza separativa garantizada y separación amurallada topográfica absoluta para cada sábana y compartimento tubular paramétrico.

\textbf{Bastión 3 (Ascenso al podio del laudo de Biyectividad Homeomórfica y Reconocimiento y Ascensión Estatus de Homeomorfismo Local Invariable del Operador e Imposición Rectora de Ascenso Biyectiva Homeomorfa y Cristalización Universal del título Local de Intercambio Continuo Restrictivo Bi-Evaluado):} 
Fijamos la vista restrictivamente sobre una y enclaustrada particular celda sectorial y parcela acordonada sábana básica representativa genérica de la estratosfera tubular paramétrica acotada restrictivamente en lámina particular estipulada como $\mathcal{N}(D, \mathbb{F} + c)$. La inmersión y restrictiva intromisión operatoria y reducción focal inter-tubular analítica originando y avalando el traspaso operando al frente directriz $\mathcal{D}_c = d|_{\mathcal{N}(D, \mathbb{F} + c)} : \mathcal{N}(D, \mathbb{F} + c) \to \mathcal{N}(D, f)$ se condecora y envuelve avalado infalible con investidura y ropaje intachable incontestable revestida o arropada atestiguando ostensible de pureza Continua, enraizándose heredera y descendencia avalada colateral irreprochable dependiente directamente dependiente por heredad innata deductiva subyacente al ser emanación y dictamen operativo originario matriz estelar proveniente derivada inherente por subyacente origen y correlacionada filialmente a la operación diferencial o proceso integrador y operador diferencial canónico $d$. Se demanda refrendar y testificar la biyectividad subyugadora indisoluble: 
- Evaluando Inyectividad estricta acordonada inter-tubular restrictiva perimétrica: Si presumiésemos la ocurrencia de intromisión e igualdad espuria evaluadora en el reducto o ceno inter-tubular restrictivo o falla simétrica o colisional resolutoria y mimetismo de dos gérmenes resolutoriamente evaluadores $G_{z_1}$ acoplado y equiparable y concurrentes emparejado y colisionando de forma idéntica evaluativa indistinguible equivalente a $H_{z_2}$ enraizados originariamente ambos en origen y albergados con cobijo asilar bajo la tutela del mismo amparo lámina perimetral tubular y aval y enjambre estipulado y amurallado idéntico paramétricamente $\mathcal{N}(D, \mathbb{F} + c)$, siendo estos y ostentando envestidura o portadores ambos poseedores oriundos portando estandartes homólogos idénticos inmaculados portando e inmersos siendo y representando evaluativa sendos evaluadores germinales idénticos subyacentes evaluativos originarios procedentes descendientes de la estricta igual madre engendradora o estirpe idéntica e inmaculada matriz pura representativa holomorfa canónica evaluativa inamovible $\mathbb{F} + c$ evaluada o anclada respectiva en enclaves asimétricos nucleares posicionales geográficos dispares e inter-distantes posicionalmente ubicados subyacentes discrepantes geográficamente basales y asimétricos en $z_1$ y en un $z_2$. Y que su cruce originara identidad evaluada y cruce en destino y amparo proyectado y mapeado estipulado operando idéntico forzado igualado evaluativo diferencial igualador operante ineludible dictando igual colisional diferencial $d(G_{z_1}) = d(H_{z_2})$, es decir unificadoramente originando paridad germinal y laudo equivalente forzando al laudo subyugador $\widetilde{f}_{z_1} = \widetilde{f}_{z_2}$, la preeminencia analítica evaluativa forzosa del dogma de pre-imagen o axioma topológico y proyectivo cartográfico de anclaje base $p$ ratificado previamente atestigua y desencadena implacablemente arrastrando en consecuencia irrevocable a igualador forzoso posicional laudo dictaminando implacable resolutoriamente subyugante y unificando sentenciando que los estamentos geográficos posicionales núcleos radiantes basales anclados estelares $z_1$ encaran igualación forzosa idéntica sin matiz asimiladora y unificadora enlazando inexorablemente con $z_2$, forjando e instituyendo dictaminando el dogma y sentencia $z_1 = z_2$. Dictaminando y sentenciando sin apelación inquebrantable el cimiento inyectivo estricto garantizado impoluto puro excluyente.
- Evaluando pericial y certificando el cimiento e imposición evaluativa de Sobreyectividad exhaustiva integral omnímoda envolvente: Para todo inquilino concurrente u habitante residente espectro matricial o genérico asilado asilar oriundo y huésped representante espectral e inquilino amparado originario portado germinal en el codominio y asilo tubular receptáculo terminal o destino y amparo enjambre matriz estelar $\mathcal{N}(D, f)$ avalado asilar con estatus legal representativo matriz o genérico en $f_z$, concurre irremisible avalado garantizado infalible ineludible preexistente evaluativo primigenio latente en el reducto u enjambre laminar basal y asilo paramétrico $\mathcal{N}(D, \mathbb{F} + c)$ ostentando título y preexistencia real de su retro-portador primigenio encubierto o germen nativo y lecho evaluado o primigenio encasillado u antecedente o cuna germinal estelar encumbrada alzada matriz genérica pura estipulada evaluativa en $H_z$ subyacente y encarnando preexistencia original y atestiguando o portando ser originada inducida evaluativamente o engendrada por la primigenia y lecho matriz evaluadora matricial primigenia base y cuna de origen matriz original $\mathbb{F} + c$, quien, acatando el mandato evaluador diferencial y sumisa al escrutinio evaluativo del operador derivativo, restituye inmaculada evaluando $H'_z = f_z$, coronando atestiguando garantizando colmado y avalando irrefutable y exhaustivamente la exigencia de relleno y colmatación exhaustiva integral integral y asegurando aval incontestable de sobreyectividad irrefutable pura.
El cruce originario bi-evaluado e isomorfismo cartográfico y topológico homeomorfo originador del desdoblamiento inverso analítico generador de la sección pre-evaluadora seccional inversa originaria evaluativa y matriz reversa originadora avaladora cartográfica y restitutiva $s_c: \mathcal{N}(D, f) \to \mathcal{N}(D, \mathbb{F} + c)$ se rubrica evaluadora y certifica atestiguada investida con ropaje de pureza y ostentando aval e incontestable credencial garantizada avaladora de estatus formal purista puro avalado con estatus legal Continuo, cristalizado y avalado operativamente puesto y dado a merced de que esta directriz matriz evaluadora engendra en cada evaluación o vector posicional analítico evaluador o inducción germinal asilar evaluador originador general asignando o transmutando a favor y designando en cada $f_z$ evaluador y amparado al germen colateral derivado correlacionado inducido y paramétricamente amarrado en estricta y con denodado e inter-vinculación indisoluble de lecho germinada y primigenia avalada $\mathbb{F} + c$ evaluada originaria originadora germinal basal evaluativa asilar concéntrica evaluada anclada colateral y nuclear y anidada en anclaje y asilo u envolvente vecinal nuclear en $z$.
Habiendo escalado y rendido satisfactoria incontestable unánime y triunfal superada avalada con gloria analítica todas las lizas periciales del riguroso escrutinio demostrando que cada germen estelar arbitrario e ineludible y espectro del espacio magno originario $\mathcal{O}$ avalado de cuna o portador y habitante u originario y asilar asilo base cuenta sin sesgo fisura y asegurado con cobertura vecindad asilar originaria con un estricto blindado vecinal tubular asilar o matriz de reducto enjambre amparante y cuna base bien resguardado cubierto y acordonado con rigor blindado garantizado topológicamente en lecho evaluado y atestado de capas segregadas, colegimos atestiguando irrecusable irrevocable dictaminando infalible sentenciando que el magno operador matriz $d$ reina majestuoso e imperecedero, avalado con todas las de la ley axiomática consagrado sin cortapisa y sin parangón dictaminando inquebrantable asertivo erigiéndose soberano amparado y encarnando la excelsa noble estirpe de Recubrimiento Topológico y Operativo Formal Analítico Absoluto Puro Verdadero.
\end{proof}

\begin{lema}[El Enlace Trascendental e Inter-traducción e Identidad Resolutoria: El Levantamiento Guiado Ascensional enmarcado bajo tutela restrictiva de $d$ coincide irrefutablemente mimetizándose fundido análogamente encarnando la esencia resolutoria e integración analítica constituyendo y resultando de manera infalible y análoga equivalente en estricto rigor y siendo intrínsecamente igualitario resolutorio a la Integral Clásica e Inquebrantable Riemanniana o del Integrando Complejo Fundamental Vectorial Lineal o Curvilínea o Funcional Analítica Primitiva Funcional Paramétrica Holomorfa Acumuladora o de Contorno Resolutoria Parametrizada u Operadora Evaluadora Direccional Curvilínea]
Convoquemos y apostemos a escena dispuesta y delineemos el encuadre operativo postulando a disposición una curva o trazado vector directriz u oscilador posicional u cinemático direccional paramétrico estipulado inquebrantable rectilíneo en dominio parametrizado temporal o vector temporal limitante de génesis a declinación y acotado y contenido y estipulado como $\gamma: [0, 1] \to \mathbb{C}$, la cual debe por prerrogativas de rigor normativo ostentar el porte y título garantizado irrenunciable que impone o rubrica evaluada avalada estipulada y gozando perentoriamente y calificada estipulada perimetralmente con grado y cualidad de ser o conformarse Diferenciable con escisiones a trozos o paramétricamente seccionada por partes suaves evaluadas a nivel segmental tramos, ostentando amarradero y anclaje inicial originario u originador de partida enclavada o fijada u estipulada matriz evaluada concéntricamente inalterable e impuesta rígidamente fijada inicializando su génesis en $\gamma(0) = a$.
Desplazamos al mismo ruedo reclutando convocando integrando amarrando un vector o inquilino germinal pre-establecido asilar oriundo y genérico o representativo inmaculado $f_a \in \mathcal{O}_a$.
Concedamos e instituyamos invocando y elevando resolutivamente y encauzando operativamente a favor del requerimiento avalado una y pre-existente amparante ascensional o constructo levantador u trazadora o senda de vuelo y levantamiento paramétrico o enrutamiento guiado estelar elevado o vector super-directriz e inducción alada paramétrica e inducida formal asilar y levantada y avalada estipulada y enrutada o encarnada unificadora paramétrica referenciada genérica y bautizada o coronada formalmente o signada paramétricamente e indexada vectorialmente o denotada unánimemente como $\mathcal{F}: [0, 1] \to \mathcal{O}$ orquestando y cumplimentando la función dictada en rigor subyugante o matriz de levantamiento del trayecto rastro directriz $\gamma$ operante o mediada subyugante filtrada conducida inmersa sujeta asilar acoplada a través e inmersa y canalizada operativamente y filtrada tamizada restringida por el colador matriz restrictivo u amparo del filtro u operante magno universal recubridor estelar garantizado avalado y refrendado analítico estipulado rigurosamente y previamente estatuido y consagrado e instaurado soberano magno recubrimiento $d$, asumiendo e imponiendo sumisión de cimentación y arranque perimetral u condicional matriz obligante evaluador conminatorio matriz restrictivo garantizador amarrando u obligante amparante origen matriz resolutorio originario resolutivo o atadura de arranque exigido limitante matriz estipulado o forzoso evaluado en génesis origen evaluado $\mathcal{F}(0) = f_a$.
Bajo la sombra cobijo o amparo de estas murallas de rigor pericial postulatorio axiomático y condicionantes inquebrantables, el laudo y axioma atestigua rubrica sanciona y certifica perentorio irrenunciable e inapelable que para cualquier punto de evaluación y testeo de tránsito y corte horario genérico de muestreo o instante y franja o coordenada arbitraria temporal paramétrica temporal evaluativa aleatoria e inter-estipulada anidada y acordonada delimitada y albergada circunscrita temporal $t \in [0, 1]$ y de manera inexorable:
El Valor Extrahído Destilado Puro Evaluado Cuantificable y Resultante Tangible Numérico del Evaluador germinal o germen en vuelo u evaluador paramétrico matriz enrutado y transicional de vuelo o levantamiento $\mathcal{F}(t)$ constreñido en evaluación aplicativa local puntual y en enclave u evaluación forzada local de anclaje anidado estelar evaluado en y recayendo resolutivamente u amparado o condicionado con amarre evaluador constrictivo colindante en su centro territorial evaluativo u amparado nuclear base u evaluador basamento cartográfico evaluador en y evaluativo en el enclave base referenciado originador núcleo temporal asilar evaluado u epicentro y evaluación o nudo geográfico nuclear base temporal nuclear $\gamma(t)$ satisface sacia con plena certidumbre absoluta u obedece y encarna mimetizado asilar o se amolda resolutorio avala inquebrantable obediencia pura evaluando irrefutable y conforma con igualdad pura estricta analítica inmaculada cristalina y rinde y rinde tributo y asimila en equiparación rigurosa pura o enmarcada asertiva amparada o rubricada asimilada y sentenciada como estipulando atestiguando o revelando en igualdad o idénticamente equivalente o en idéntica ecuación asimilatoria inquebrantable e indisoluble matriz asertiva pura y rigurosa igualativa matriz dictando:
\begin{equation}
\mathcal{F}(t)(\gamma(t)) = f_a(a) + \int_{\gamma|_{[0,t]}} f(z) dz.
\end{equation}
\end{lema}
\begin{proof}
Para cimentar y amurallar de forma monolítica indisoluble perenne indestructible impoluta sin brecha o filtración fisura fisura o fisura lógica o vacilatorio y atestiguar y refrendar a cabalidad plena total inmaculada cristalina exhaustivamente la imposición axioma y ecuación o igualdad igualadora matriz asimilatoria, acometemos a someterla evaluando diseccionando desgajando escrutando segmentando evaluando o inspeccionando pericial o analizando la directriz u vector u horquilla de validación a lo ancho de los eslabones eslabones consecutivos o compartimentos tramos segmentos y peldaños intermedios transicionales franjas franja evaluadora y transicional o eslabón transicional matriz o vector de transición o puente intermedio subyacente de la disección partición original particionadora acotada fraccionadora particionada amparadora transicional inductiva $I_v = [t_v, t_{v+1}]$ proveniente y evocada heredera amparada subyacente y estipulada al albor evaluativo en asilo referenciado y desprendida de la división matricial matriz operadora o fase de despiece operante subyacente dictaminado inductiva matriz (orquestada o dictaminada en albor inicial de particiones de recubrimientos Proposición 2.1.4).
Avalamos aseveramos estipulando estatuimos perentorio pericial testificando que en la intimidad reclusa albergada en el seno de cada uno de los tramos cerrados y resguardos inter-eslabones transicionales y o nichos y subintervalos y recinto transicional o sub-compartimento de confinamiento matricial, el levantamiento y el vector guía evaluador matriz dictador encumbrado transicional vector ascendido dictaminador o vector evaluador o germen de transición paramétrico en ruta o tránsito elevador y guía o rastro y ascensión levantada transicional $\mathcal{F}(t)$ avala y ostenta con credencial porta y enarbola ser sin dubitación o resquicio alguno indiscutible o incuestionablemente equivalente o encarnando o materializando al vector y evaluando el germen u evaluador e inquilino asilar y amparador de ser y fungir evaluador originario idéntico homólogo matriz de o matriz análoga rama matriz o evaluador estelar evaluando e idéntico idéntico y erigiéndose y manifestándose ostentando el rango estelar analítico o rama subyugante analítica estelar de y evaluando como portador representador idéntico subyacente rama germinal u originador matriz estelar subyacente a y representante o matriz analítica e igualando la progenie analítica evaluativa matricial de una pura genuina incontestable o garante auténtica y asertiva función originadora y matriz primitiva primigenia matriz generatriz matriz estelar local o primitiva local genuina subyacente y encarnadora local estelar matriz pre-imagen primitiva o matriz operadora originaria encarnadora local evaluativa matriz $\mathbb{G}_v$ la cual inquebrantable acata sacia evaluadora atestiguando o evaluando obediencia igualitaria o igualadora derivativa idéntica dictaminando estipulada dictando obediencia a la imposición derivativa ineludible dictaminando irrefutable $\mathbb{G}_v' = f$.
Bajo el prisma y lupa de evaluación transicional del vector u operador u dogma y postulado avalado o axioma y teorema irrenunciable integrador resolutorio analítico dogmático imperioso cimiento o dogma o laudo analítico vector evaluador matriz resolutorio Teorema Supremo Resolutorio Operador Integral y Postulado Supremo y Magno Fundamental Evaluador matriz o Resolutorio Cimiento e Integro Supremo Fundamental del decálogo Analítico Fundamental Complejo (El Teorema Fundamental del Cálculo Complejo avalado resolutorio Integrador), atestiguamos u laudo o laudo evaluador certifica atestiguando u rubrica pericial sanciona avala dictando y atestigua inquebrantable garantizado que evaluando el cúmulo o acopio o peso integrador aditivo amontonado de la integral de trazo y recorrido Riemanniana o sumatoria integradora evaluada y acopio perimetral curvilínea sobre este o a lo largo del encuadre y lecho eslabón matriz o trazo vector matriz o compartimento y segmento o trazado de la o del arco subyacente $\gamma|_{I_v}$ es y rinde fiel o equipara y calibra o compensa u resulta rindiendo atestiguando y rinde o concuerda evaluando asimilando resulta y compensando idéntica igualadora en rinde fielmente y cuantifica sin desfasaje el remanente e impacto y fluctuación y balance diferencial u salto cambio originador incremento u desfasaje o diferencial u incremento matriz u cambio salto posicional u diferencial matriz o cambio y fluctuación o variación evaluada paramétrica de esta originadora o o primigenia o primitivo base y matriz originadora evaluada o matriz y evaluada o lecho u matriz generadora evaluada enclavada posicional y operadora primitivo base matriz o originadora local $\mathbb{G}_v$ entre el cierre y albores fronteras y encuadre u inicio evaluador y remate u terminal cota o amparo terminal o polos y linderos u fronteras y extremidades matriz de la ruta o franja u franja matriz o trazo enmarcado u matriz segmento o de la directriz evaluadora y curva.
Consecuentemente amontonando encadenando amalgamando asimilando acoplando hilvanando unificando agrupando adicionando o y engarzando eslabón o o segmento a y segmento con tramo, la suma integradora resolutoria colapsa y rinde y atestigua u resulta desvelando telescopiando evaluadora en aval telescópico colapsando unificador rindiendo evaluador igual y dictando o revelando igualador originador resultando o equiparando matriz pura arrojando y rindiendo fidedigna dictando el acumulador y asilo o saldo originador totalitario amarrado y evaluador puro neto e y lauro puro balance o saldo e incremento puro evaluador y neto resolutoria y desfasaje vector de origen a fin garantizando y consolidando con certidumbre y cristalizando infalible la validez laudo e imposición y certeza y ratificando inmaculada con aval resolutorio impoluto la legitimidad e identidad de la e ecuación u originadora matriz estipulada integradora e formulación originaria matriz de e igualación integradora o identidad analítica matriz resolutoria e igualdad matriz atestiguada o asertiva formulación analítica de la igualdad formulada postulatoria.
\end{proof}

\textbf{Extrapolación Universal Postrera e Impacto Doctrinal Teórico y de Extensión y Ramificación Resolutoria:} 
Esta magistral revelación analítica y cimiento resolutorio impoluto exuda irradia y propaga una onda expansiva doctrinaria de calado telúrico, puesto que desvela sin asomo o atisbo o reserva y promulga aseverando rubricando que el y fenómeno clásico procedimental integrativo paramétrico de integración evaluativa a lo extenso de e integrando evaluando sobre matrices vectoriales estelares curvas adquiere y exhibe un paralelismo o equivalencia dictaminadora u hermandad isomorfa u isomorfismo e y correlato isomorfo y consanguinidad u equivalencia gemelar y asimetría o transmutación idéntica u es y amolda a y e estatus isomorfo u equivalente irrenunciable e igualitario a o e inquebrantable consustancial asimétrica a de naturaleza idéntica y simétrica o análoga de paridad pura u o idéntica evaluadora purista equivalencia irrestricta o isomórfica evaluando igualitaria o equivalente paralela irrestricta con el vector o fenómeno topológico ascensional o de levantamiento de matriz curvas y o transicional de curvas paramétrico. Esto confiere provee faculta y y dota del andamiaje asilar y utillaje permite y dota la herramienta y catapulta u empodera analítico matriz y otorga de plataforma u habilita de facto o matriz faculta para encauzar orquestar reformular y transfigurar o e transmutar el majestuoso axioma y e y pilar estelar y magno y regente axioma o dogma teorema estelar resolutorio o u estandarte y Teorema estelar Cauchyano Magno fundamental Teorema estelar Analítico Universal y y el integrador cauchyano o de Teorema e de laudo de Cauchy de cauchy en el seno o marco y en jerga y glosario o nomenclatura y en lenguaje u albergue y amparo matriz del argot matriz o y seno del prisma e o paradigma dogma o prisma axiomático enmarcado bajo axiomas de la pura matriz doctrina u topología matriz teoría estelar y teoría matriz recubridora o teoría magno recubridora de evaluadores o y doctrina y laudo de recubrimientos de y levantamientos de e la e y pura de espacios y asilo y de recubridores u recubrimientos o teoría abstracta recubridora analítica matricial paramétrica de e recubrimientos, designio y encomienda pericial y y asilo y gesta magno pericial empresa y o y matriz matriz y proyecto matriz analítico estelar y gesta o postulado o encargo analítico matriz resolutorio asilar que y estelar que y matriz que se consumará abordará y u y fraguará materializará cristalizará o y o se desarrollará u y e pericialmente acometerá y desarrollará explayará o e orquestará se e y desmenuzará estipulará en su plenitud u se desgranará pericialmente y e o evaluará se y consagrará desatará abordará a o cristalizará o atestiguará a gran escala asilar asertiva con lujo estelar de y meticulosidad de detalles minucioso en enclaves en y en estipulación y o en la posterior a y o u y en el seno y reducto matriz asilar matriz u asilar u y o y y ulterior matriz en estrado o reducto o asilo analítico y en y apartado y segmento subsecuente u sección matriz y evaluativa estelar u en las lindes estelares de la sección venidera y subsecuente u en el porvenir o recoveco o subsiguiente tramo y venidera o en y en y subsecuente y en y sección ulterior en la u en el pasaje o subsecuente matriz de análisis u venidera evaluadora y correlacionada o sección colindante subsecuente y y sección ulterior y o a o matriz de análisis posterior en y en o de estipulación y en el epígrafe venidero y venidera u y colindante o o colindante o de y ulterior y subsecuente u venidero u sección correlacionada y sección o y parágrafo y venidera \S 2.4.


\section{2.4 El Teorema de Monodromía y la Forma Homotópica del Teorema de Cauchy}

La sección anterior estableció que la integración a lo largo de una curva $\gamma$ de una función holomorfa $f=g'$ (derivada) equivale a levantar $\gamma$ al haz $\mathcal{O}$ respecto al recubrimiento $d$. La pregunta natural que surge a continuación es: ¿cuándo es que el valor final de la integral (o del levantamiento) depende solo de los extremos de la curva y no del camino específico recorrido? Esta es precisamente la pregunta fundamental que responde el Teorema de Monodromía, cuya respuesta completa involucra inexorablemente la noción topológica de homotopía entre curvas.

\begin{definicion}[Homotopía con extremos fijos]
Sea $X$ una variedad topológica y sean $\gamma_0, \gamma_1 : [0,1] \to X$ dos curvas continuas tales que $\gamma_0(0) = \gamma_1(0) = a$ y $\gamma_0(1) = \gamma_1(1) = b$. Decimos que $\gamma_0$ es homotópica a $\gamma_1$ con extremos fijos $a$ y $b$ si existe una aplicación continua $F : I \times I \to X$ (donde $I = [0,1]$) tal que para todo $t, u \in I$ se cumple:
\begin{enumerate}
    \item $F(t,0) = \gamma_0(t)$ (la homotopía inicia en $\gamma_0$)
    \item $F(t,1) = \gamma_1(t)$ (la homotopía termina en $\gamma_1$)
    \item $F(0,u) = a$ (el extremo inicial se mantiene fijo durante toda la deformación)
    \item $F(1,u) = b$ (el extremo final se mantiene fijo durante toda la deformación)
\end{enumerate}
En este caso escribimos $\gamma_0 \sim \gamma_1$ (con extremos fijos $a, b$).
\end{definicion}

\begin{observacion}[Necesidad de fijar ambos extremos e implicaciones topológicas]
\begin{enumerate}
    \item Si $a=b$ y la homotopía $F$ fija el punto base ($F(0,u)=F(1,u)=a$), cada curva intermedia $\gamma_u(t)=F(t,u)$ es un lazo cerrado con base en $a$, puesto que $\gamma_u(0)=F(0,u)=a$ y $\gamma_u(1)=F(1,u)=a$. En este caso hablamos de lazos homotópicos en $X$ con base $a$, lo cual da origen a la definición del grupo fundamental.
    \item \textbf{Es CRUCIAL fijar ambos extremos.} Sin la restricción de fijar los extremos, toda curva sería trivialmente homotópica a la curva constante, lo cual vaciaría de significado topológico al concepto. Para demostrar esto rigurosamente, tomemos cualquier curva $\gamma : [0,1] \to X$. Definamos $H(t,u) = \gamma((1-u)t)$. Esta función es claramente continua al ser composición de funciones continuas. Verificamos las condiciones: para $u=0$, $H(t,0) = \gamma(t)$. Para $u=1$, $H(t,1) = \gamma(0)$, que es la curva constante en el punto inicial. Además, $H(0,u) = \gamma(0)$ para todo $u$. Sin embargo, observemos el extremo final: $H(1,u) = \gamma(1-u)$, lo que significa que el punto final recorre la curva $\gamma$ hacia atrás. Así, sin la condición explícita de $F(1,u)=b$, toda curva se retrae a su punto inicial. La condición de fijar ambos extremos es, por tanto, el ingrediente topológico esencial que captura la verdadera obstaculización geométrica del espacio (los "agujeros").
    \item La relación de homotopía con extremos fijos es una relación de equivalencia en el conjunto de todas las curvas continuas de $a$ a $b$ en $X$. 
    \begin{itemize}
        \item \textbf{Reflexividad:} Toda curva $\gamma$ es homotópica a sí misma tomando la homotopía constante en el tiempo de deformación $F(t,u) = \gamma(t)$. Esta $F$ es continua y trivialmente satisface $F(t,0)=\gamma(t)$, $F(t,1)=\gamma(t)$, $F(0,u)=\gamma(0)=a$ y $F(1,u)=\gamma(1)=b$.
        \item \textbf{Simetría:} Si $F(t,u)$ es una homotopía de $\gamma_0$ a $\gamma_1$, entonces $G(t,u) = F(t,1-u)$ es una homotopía de $\gamma_1$ a $\gamma_0$. Es continua por ser composición con $u \mapsto 1-u$, y satisface $G(t,0)=F(t,1)=\gamma_1(t)$, $G(t,1)=F(t,0)=\gamma_0(t)$, $G(0,u)=F(0,1-u)=a$ y $G(1,u)=F(1,1-u)=b$.
        \item \textbf{Transitividad:} Si $F_1$ deforma $\gamma_0$ en $\gamma_1$ y $F_2$ deforma $\gamma_1$ en $\gamma_2$, concatenamos en la variable $u$: definimos $H(t,u) = F_1(t, 2u)$ si $0 \le u \le 1/2$, y $H(t,u) = F_2(t, 2u-1)$ si $1/2 \le u \le 1$. La continuidad de $H$ está garantizada por el Lema del Pegado (Pasting Lemma), ya que los conjuntos $[0,1] \times [0,1/2]$ y $[0,1] \times [1/2,1]$ son cerrados cuya intersección es $[0,1] \times \{1/2\}$, y en ella $F_1(t,1) = \gamma_1(t) = F_2(t,0)$, por lo que $H$ está bien definida y es continua en todo el cuadrado.
    \end{itemize}
    \item Podemos visualizar el dominio de la homotopía y sus condiciones de frontera en el siguiente diagrama:
\end{enumerate}

\begin{center}
\begin{tikzpicture}[scale=2]
    % Cuadrado
    \draw[thick] (0,0) rectangle (2,2);
    % Etiquetas de vértices
    \node[below left] at (0,0) {$(0,0)$};
    \node[below right] at (2,0) {$(1,0)$};
    \node[above left] at (0,2) {$(0,1)$};
    \node[above right] at (2,2) {$(1,1)$};
    % Condiciones de borde
    \node[below] at (1,0) {$F(t,0) = \gamma_0(t)$};
    \node[above] at (1,2) {$F(t,1) = \gamma_1(t)$};
    \node[left, align=center] at (0,1) {$F(0,u)$\\$= a$};
    \node[right, align=center] at (2,1) {$F(1,u)$\\$= b$};
    % Curva intermedia
    \draw[->, thick, blue, dashed] (0,1) .. controls (0.5,1.5) and (1.5,0.5) .. (2,1);
    \node[above, blue] at (1,1.1) {$F(t,u) = \gamma_u(t)$};
    
    % Ejes
    \draw[->] (-0.5,-0.5) -- (0.5,-0.5) node[right] {$t$};
    \draw[->] (-0.5,-0.5) -- (-0.5,0.5) node[above] {$u$};
\end{tikzpicture}
\end{center}
\end{observacion}

\begin{teorema}[2.4.1 - Teorema de Monodromía (versión general)]
Sea $p : X' \to X$ una aplicación de recubrimiento. Sean $\gamma_0, \gamma_1 : [0,1] \to X$ dos curvas homotópicas con extremos fijos mediante una homotopía continua $F : I \times I \to X$. Sea $a = F(0,0)$ el punto inicial común, y sea $a' \in X'$ un punto en la fibra sobre $a$, es decir, $p(a') = a$. Entonces los levantamientos únicos $\gamma'_0$ y $\gamma'_1$ de $\gamma_0$ y $\gamma_1$ respectivamente, que comienzan en $a'$, tienen el mismo punto final: $\gamma'_0(1) = \gamma'_1(1)$.
\end{teorema}
\begin{proof}[Demostración ultra-rigurosa]
Procederemos sistemáticamente en 4 pasos lógicos:
\paragraph{Paso 1: Construcción del levantamiento de la homotopía $F$.} Para cada $u \in I$ fijo, consideramos la curva $\gamma_u : I \to X$ definida por $\gamma_u(t) = F(t,u)$. Como $p: X' \to X$ es una aplicación de recubrimiento, la Proposición 2.1.1 garantiza que para toda curva continua $\gamma_u$ y punto inicial $a' \in X'$ tal que $p(a') = \gamma_u(0) = a$, existe un único levantamiento $\gamma'_u : I \to X'$ tal que $p \circ \gamma'_u = \gamma_u$ y $\gamma'_u(0) = a'$. 
Definimos la función $F' : I \times I \to X'$ puntualmenete mediante $F'(t,u) = \gamma'_u(t)$. Por construcción, $p(F'(t,u)) = p(\gamma'_u(t)) = \gamma_u(t) = F(t,u)$ para todo $(t,u) \in I \times I$, y $F'(0,u) = \gamma'_u(0) = a'$. 

\paragraph{Paso 2: Continuidad de $F'$.} Debemos demostrar que esta función $F'$ definida punto a punto es continua en todo el compacto $I \times I$. Sea $(t_0, u_0) \in I \times I$. Puesto que $p$ es un recubrimiento, existe una cubierta abierta $\{U_x\}_{x \in X}$ de $X$ formada por entornos bien cubiertos. Al ser $F$ continua, la preimagen $F^{-1}(U_x)$ para cada $x$ da una cubierta abierta de $I \times I$. Puesto que $I \times I$ es un espacio métrico compacto, aplicamos el Lema del Número de Lebesgue: existe $\delta > 0$ tal que cualquier subconjunto de diámetro menor a $\delta$ está contenido enteramente en algún $F^{-1}(U_x)$. 
Esto nos permite subdividir el cuadrado $I \times I$ en una cuadrícula fina definida por particiones $0=t_0 < t_1 < \dots < t_N = 1$ y $0=u_0 < u_1 < \dots < u_M = 1$, de forma que cada rectángulo cerrado $R_{v,w} = [t_v, t_{v+1}] \times [u_w, u_{w+1}]$ se aplica mediante $F$ dentro de un único entorno bien cubierto $U_{v,w}$. 
Sobre $U_{v,w}$, la preimagen $p^{-1}(U_{v,w})$ es unión disjunta de abiertos en $X'$, cada uno homeomorfo a $U_{v,w}$ vía la restricción de $p$. La continuidad de $F'$ se prueba inductivamente sobre las celdas $R_{v,w}$. Comenzando en $R_{0,0}$, el punto inicial $F'(0,0) = a'$ cae en alguna hoja $V_{0,0}$ sobre $U_{0,0}$. Como las hojas son disjuntas y $R_{0,0}$ es conexo, todo el levantamiento $F'(R_{0,0})$ debe residir en $V_{0,0}$, y por unicidad de la inversa local continua, $F'|_{R_{0,0}} = (p|_{V_{0,0}})^{-1} \circ F|_{R_{0,0}}$, lo que prueba que $F'$ es continua en la primera celda. Inductivamente (de izquierda a derecha y de abajo a arriba), el valor de $F'$ en los bordes compartidos entre celdas determina unívocamente en qué hoja continua debe caer el levantamiento de la celda adyacente. Por el Lema del Pegado sobre cerrados que intersecan, $F'$ es globalmente continua sobre $I \times I$.

\begin{center}
\begin{tikzpicture}[scale=3]
    % Cuadrícula de I x I
    \draw[step=0.5, gray, thin] (0,0) grid (2,2);
    \draw[thick] (0,0) rectangle (2,2);
    \node[below left] at (0,0) {$0$};
    \node[below right] at (2,0) {$1$};
    \node[above left] at (0,2) {$1$};
    \node[below] at (1,-0.1) {$t$};
    \node[left] at (-0.1,1) {$u$};
    
    % Flecha hacia X
    \draw[->, thick] (2.2, 1) -- (2.8, 1) node[midway, above] {$F$};
    \draw (3, 0.5) ellipse (0.3 and 0.2) node {$U_{v,w}$};
    
    % Flecha levantada
    \draw[->, thick, red] (2.2, 1.5) -- (2.8, 2.5) node[midway, above left, red] {$F'$};
    \draw[red] (3, 2.5) ellipse (0.3 and 0.2) node {$V_{v,w}$};
    
    % Proyección
    \draw[->, thick] (3, 2.2) -- (3, 0.8) node[midway, right] {$p$};
\end{tikzpicture}
\end{center}

\paragraph{Paso 3: Condiciones de frontera del levantamiento $F'$.} Consideremos la función $\phi(t) = F'(t,0)$. Esta función es un levantamiento continuo de $\gamma_0(t) = F(t,0)$ que empieza en $\phi(0) = F'(0,0) = a'$. Por la rigurosa unicidad del levantamiento de curvas (Proposición 2.1.1), se debe tener $F'(t,0) = \gamma'_0(t)$ para todo $t$. De manera análoga, $F'(t,1) = \gamma'_1(t)$. 
Para la frontera izquierda, fijemos $t=0$. La curva $u \mapsto F'(0,u)$ levanta a la curva constante $u \mapsto F(0,u) = a$. El levantamiento de la curva constante en $a$ que parte de $a'$ es idénticamente la curva constante $a'$. Luego $F'(0,u) = a'$ para todo $u$. Esto concuerda con la definición constructiva dada en el Paso 1.

\paragraph{Paso 4: Igualdad estricta de puntos finales.} Consideremos ahora la frontera derecha del cuadrado, es decir, el conjunto $t=1$. La función continua $u \mapsto F'(1,u)$ representa una curva en $X'$. Su proyección es $p(F'(1,u)) = F(1,u) = b$ para todo $u \in I$. Esto implica que para todo $u \in I$, $F'(1,u)$ pertenece a la fibra $p^{-1}(\{b\})$. 
Ahora, por ser $p$ una aplicación de recubrimiento, todo punto de la fibra sobre $b$ tiene un entorno abierto $U'$ tal que $p|_{U'}$ es un homeomorfismo, y consecuentemente $U' \cap p^{-1}(\{b\}) = \{b'\}$. Es decir, la topología de subespacio heredada por la fibra $p^{-1}(\{b\})$ desde $X'$ es la topología discreta. 
Por un teorema fundamental de topología general, la imagen continua de un conjunto conexo ($I = [0,1]$) debe ser conexa. El único subconjunto conexo no vacío de un espacio discreto es el conjunto formado por un solo punto. Por lo tanto, la aplicación $u \mapsto F'(1,u)$ debe ser constante. En particular, sus valores en $u=0$ y $u=1$ deben coincidir, lo que nos da:
$$ \gamma'_0(1) = F'(1,0) = F'(1,1) = \gamma'_1(1) $$
Esto concluye la demostración rigurosa del Teorema de Monodromía.
\end{proof}

\begin{proposicion}[Invarianza del punto final y monodromía de un germen]
Si $p : X' \to X$ es una aplicación de recubrimiento, $a' \in X'$, y $\gamma_0 \sim \gamma_1$ son curvas homotópicas con extremos fijos en $X$ desde $a = p(a')$ hasta $b$, entonces los levantamientos desde $a'$ llegan al mismo punto en la fibra sobre $b$. 

En particular, aplicando este marco general al haz de gérmenes analíticos $\mathcal{O}$ dotado con la proyección canónica (recubrimiento) $d : \mathcal{O} \to \Omega$: si $f_a \in \mathcal{O}_a$ es un germen holomorfo en $a$, y $\gamma_0 \sim \gamma_1$ son curvas homotópicas con $a$ y $b$ como extremos fijos, y si además $f_a$ admite continuación analítica a lo largo de AMBAS curvas, entonces los gérmenes resultantes en el punto $b$ son idénticos: $\mathcal{F}_0(1) = \mathcal{F}_1(1) \in \mathcal{O}_b$. En otras palabras, la continuación analítica a lo largo de curvas homotópicas produce un resultado final único e invariante bajo homotopía.
\end{proposicion}

\begin{corolario}[Forma Homotópica del Teorema de Cauchy]
Sea $\Omega \subset \mathbb{C}$ un conjunto abierto y $f \in \mathcal{H}(\Omega)$ una función holomorfa. Si $\gamma_0$ y $\gamma_1$ son dos curvas diferenciables a trozos en $\Omega$, desde $a$ hasta $b$, que son homotópicas en $\Omega$ con extremos fijos, entonces:
$$ \int_{\gamma_0} f(z) dz = \int_{\gamma_1} f(z) dz $$
\end{corolario}
\begin{proof}
Procedemos primero con un caso especial: Supongamos que $f$ admite una primitiva global $g \in \mathcal{H}(\Omega)$, es decir, $g' = f$. En este escenario, la integral a lo largo de cualquier curva $\gamma_i$ se reduce a la diferencia de los valores de la función $g$ evaluada en los extremos: $\int_{\gamma_i} f(z) dz = g(b) - g(a)$. Esta igualdad es inmediata y el resultado se sigue trivialmente, ya que la integral no depende del camino en absoluto.

Sin embargo, para el caso general, $f$ puede NO poseer una primitiva global. En este caso apelamos a la estructura local. Como $f$ es holomorfa, en cada disco contenido en $\Omega$, admite una primitiva local. Integramos $f$ a lo largo de las curvas y esta integral se puede interpretar como la continuación analítica de una primitiva local a lo largo de $\gamma$.
Específicamente, consideramos el germen de una primitiva local $G_a$ en $a$. La integración a lo largo de $\gamma_i$ se corresponde con la evaluación final del levantamiento de la curva al haz de gérmenes. Puesto que $\gamma_0 \sim \gamma_1$, el Teorema de Monodromía asegura que el germen obtenido en el punto $b$ al seguir $\gamma_0$ es exactamente el mismo que se obtiene al seguir $\gamma_1$. Por lo tanto, el valor de la integral ---que es simplemente la evaluación de este germen en el punto final menos el germen inicial evaluado en el inicio, ajustado localmente por parches--- dependerá solo de la homotopía. Más formalmente, al subdividir las curvas vía el lema de Lebesgue (como en el paso 2 de la monodromía) y aplicar integrales locales sobre discos donde sí hay primitivas exactas, la suma telescópica de estas integrales produce el mismo valor, asegurando $\int_{\gamma_0} f dz = \int_{\gamma_1} f dz$. Este notable resultado generaliza y sobrepasa al Teorema de Cauchy clásico.
\end{proof}

\begin{definicion}[Dominio Simplemente Conexo]
Un dominio (conjunto abierto y conexo) $\Omega \subset \mathbb{C}$ se dice \textbf{simplemente conexo} si toda curva cerrada (lazo) en $\Omega$ es homotópica (con extremos fijos en su punto base) a la curva constante en $\Omega$. De manera equivalente en el lenguaje de topología algebraica, el grupo fundamental $\pi_1(\Omega, a)$ es trivial para todo punto base $a \in \Omega$.
\end{definicion}

\begin{observacion}[Interpretación topológica rigurosa de simplemente conexo]
\begin{enumerate}
    \item Intuitivamente, se dice que $\Omega$ es simplemente conexo si y sólo si "no tiene agujeros". Todo lazo puede deformarse continuamente, contrayéndose dentro de los confines de $\Omega$, hasta convertirse en un único punto matemático. Los prototipos absolutos son el plano complejo $\mathbb{C}$ y el disco unitario abierto $\mathbb{D}$.
    \item \textbf{Demostración de que el plano perforado $\mathbb{C}^* = \mathbb{C} \setminus \{0\}$ NO es simplemente conexo:} Consideremos el lazo canónico $\gamma(t) = e^{2\pi i t}$ para $t \in [0,1]$, que recorre el círculo unitario una vez en sentido antihorario. Si $\mathbb{C}^*$ fuese simplemente conexo, existiría una homotopía continua $F : I \times I \to \mathbb{C}^*$ tal que $F(t,0) = \gamma(t)$ y $F(t,1) = 1$ (la curva constante). Definamos la función $u \mapsto n(F(\cdot, u), 0)$, es decir, el índice de enrrollamiento o número de vueltas respecto al origen para la curva intermedia parametrizada por $u$. 
    Calculamos explícitamente:
    $$ n(F(\cdot, u), 0) = \frac{1}{2\pi i} \int_{F(\cdot, u)} \frac{dz}{z} $$
    Puesto que el integrando $1/z$ es holomorfo en $\mathbb{C}^*$ y $F$ es continua, la integral sobre el parámetro $u$ varía continuamente. Sin embargo, el índice de enrrollamiento toma por definición valores puramente enteros ($\mathbb{Z}$). Una función continua con dominio conexo ($[0,1]$) y codominio discreto ($\mathbb{Z}$) debe ser forzosamente constante. 
    Para $u=0$, $\frac{1}{2\pi i}\int_{\gamma} \frac{dz}{z} = 1 \neq 0$.
    Para $u=1$, la integral sobre un punto constante es $0$. 
    Tendríamos $1 = 0$, lo cual es una contradicción absurda. Por consiguiente, tal homotopía $F$ no puede existir, y $\mathbb{C}^*$ no es simplemente conexo.
    \item La diferencia abismal entre $\mathbb{C}$ y $\mathbb{C}^*$ desde la perspectiva del Teorema de Cauchy: En $\mathbb{C}$, toda función holomorfa posee una primitiva global absoluta porque toda curva cerrada es homotópicamente nula, garantizando integrales cerradas de valor 0. En $\mathbb{C}^*$, la función $f(z) = 1/z$ es holomorfa en todo el dominio pero carece irreparablemente de primitiva global, siendo la prueba su integral sobre el lazo $\gamma$ que produce $2\pi i \neq 0$.
    \item A continuación una comparativa visual de un dominio simplemente conexo vs uno no simplemente conexo:
\end{enumerate}

\begin{center}
\begin{tikzpicture}[scale=1.5]
    % Dominio Simplemente Conexo
    \draw[thick, fill=blue!10] (0,0) circle (1.5);
    \node[above] at (0, 1.6) {Simplemente Conexo (ej. Disco)};
    \draw[->, thick, blue] (0,0) ellipse (1 and 0.5);
    \filldraw[black] (1,0) circle (1pt) node[right] {$a$};
    \draw[->, dashed, blue!60] (0,0) ellipse (0.5 and 0.25);
    \draw[->, dotted, blue!40] (0,0) ellipse (0.2 and 0.1);
    \filldraw[red] (0,0) circle (1pt) node[below] {\small contraíble};

    % Dominio NO Simplemente Conexo
    \begin{scope}[xshift=4.5cm]
        \draw[thick, fill=orange!10] (0,0) circle (1.5);
        \draw[thick, fill=white] (0,0) circle (0.4);
        \node[above] at (0, 1.6) {No Simplemente Conexo (ej. Anular)};
        \node at (0,0) {$\times$};
        \node[below] at (0,-0.1) {\small agujero};
        \draw[->, thick, orange] (0,0) circle (0.9);
        \filldraw[black] (0.9,0) circle (1pt) node[right] {$b$};
        \node[below right] at (0.6,0.6) {\small Lazo atascado};
    \end{scope}
\end{tikzpicture}
\end{center}
\end{observacion}

\begin{teorema}[2.4.2 - Lazos homotópicos a cero]
En un espacio topológico $X$ simplemente conexo, todo lazo (curva cerrada) es nulo-homotópico, es decir, homotópico a la curva constante con extremos fijos. Esta proposición constituye de hecho la propia definición axiomática de la conexión simple.
\end{teorema}

\begin{teorema}[2.4.3 - Extensión Universal del Teorema de Cauchy]
Si $\Omega$ es un dominio simplemente conexo y $f \in \mathcal{H}(\Omega)$ es holomorfa, entonces para toda curva cerrada rectificable o diferenciable a trozos $\gamma$ contenida completamente en $\Omega$, se tiene invariablemente que:
$$ \int_{\gamma} f(z) dz = 0 $$
\end{teorema}
\begin{proof}
Puesto que $\Omega$ es un dominio simplemente conexo, la curva cerrada $\gamma$ es, por el Teorema 2.4.2, homotópica a la curva constante (denotémosla $c$) en su punto base mediante una homotopía con extremos fijos en $\Omega$. 
Invocando la Forma Homotópica del Teorema de Cauchy, establecemos la igualdad de integrales de camino:
$$ \int_{\gamma} f(z) dz = \int_{c} f(z) dz $$
Trivialmente, la integral de cualquier función sobre un único punto (una curva constante paramétricamente nula) es $0$. Por lo tanto, el resultado queda rigurosamente probado.
\end{proof}

\begin{observacion}[Catálogo exhaustivo de ejemplos topológicos]
\begin{itemize}
    \item \textbf{$\mathbb{C}$ es simplemente conexo:} Siendo convexo, se puede emplear la homotopía lineal explícita $F(t,u) = (1-u)\gamma(t) + ua$, la cual no escapa del plano y contrae cualquier lazo a su base.
    \item \textbf{Todo disco $D(a,r)$ es simplemente conexo:} Idéntico argumento de convexidad lineal estricta.
    \item \textbf{Todo dominio convexo es simplemente conexo:} La homotopía lineal jamás abandona el dominio por la definición geométrica de convexidad.
    \item \textbf{Todo dominio estrellado es simplemente conexo:} Contracción radial hacia el centro de la estrella a través de la homotopía lineal afín; el segmento siempre yace en el dominio.
    \item \textbf{El complemento de un rayo (como el semieje negativo) es simplemente conexo:} Es un dominio estrellado.
    \item \textbf{$\mathbb{C}^*$ y el Anular $r < |z| < R$ NO son simplemente conexos:} Probado rigurosamente arriba mediante el invariante topológico de índice de enrrollamiento nulo.
\end{itemize}
\end{observacion}

\section{2.5 Aplicaciones de la Monodromía}

En esta sección derivaremos y consolidaremos las consecuencias analíticas de mayor peso originadas en el Teorema de Monodromía, a saber: la caracterización topológica del levantamiento de recubrimientos, la existencia irrefutable de primitivas y logaritmos en geometrías simplemente conexas, y la profunda acotación en la dimensionalidad de las ramas multiformes encapsulada en el Teorema de Poincaré-Volterra.

\begin{definicion}[Propiedad de Levantamiento de Curvas - CLP]
Decimos que un homeomorfismo local $p: X' \to X$ ostenta la Propiedad de Levantamiento de Curvas (abreviada CLP) si, y solo si, para toda curva continua $\gamma : [0,1] \to X$ y todo punto $a' \in X'$ que se proyecta al origen de la curva (es decir, $p(a') = \gamma(0)$), existe un ÚNICO levantamiento continuo $\gamma' : [0,1] \to X'$ tal que $p \circ \gamma' = \gamma$ y $\gamma'(0) = a'$.
\end{definicion}

\begin{teorema}[2.5.1 - Caracterización topológica de recubrimientos por CLP]
Sea $p : X' \to X$ un homeomorfismo local entre espacios topológicos que son variedades conexas. Entonces $p$ es una aplicación de recubrimiento global si y sólo si $p$ satisface la Propiedad de Levantamiento de Curvas (CLP).
\end{teorema}
\begin{proof}
\textbf{($\Longrightarrow$)} Esta implicación es precisamente el contenido de la Proposición 2.1.1. Todo recubrimiento garantiza el levantamiento único dividiendo la curva compacta vía Lebesgue y subiendo pedazo a pedazo usando los inversos locales.

\textbf{($\Longleftarrow$)} Asumamos por hipótesis que $p$ es un homeomorfismo local que disfruta de la CLP. Debemos probar que $p$ es un recubrimiento. Elegimos un punto arbitrario $a \in X$. Queremos hallar un entorno bien cubierto de $a$. 
Por ser homeomorfismo local, elegimos un entorno abierto conexo $U$ de $a$ tal que localmente $p$ sea invertible de forma bicontinua. 
Puesto que $p$ posee la CLP, para cualquier componente $U'_j$ de $p^{-1}(U)$, cualquier curva $\gamma$ en $U$ desde $a$ hasta $x \in U$ puede levantarse unívocamente. 
Dado un punto $a'_j \in p^{-1}(a)$, definimos $U'_j$ como el conjunto de todos los puntos finales de los levantamientos de curvas en $U$ que comienzan en $a'_j$. 
Por la conexidad por arcos de $U$, todo punto $x \in U$ es alcanzable por una curva $\gamma$ desde $a$. El levantamiento $\gamma'$ termina en un punto $x' \in U'_j$. Esta construcción muestra no solo que $p$ se restringe a un homeomorfismo biyectivo $p|_{U'_j} : U'_j \to U$, sino que las hojas $U'_j$ y $U'_k$ para puntos base distintos no se intersectan (si lo hicieran, violarían la unicidad del levantamiento o la biyectividad local). Así, $p^{-1}(U)$ es unión disjunta de hojas homeomorfas a $U$, convirtiendo a $U$ en un entorno bien cubierto por definición estricta.
\end{proof}

\begin{proposicion}[Levantamiento generalizado de aplicaciones continuas]
Sea $p: X' \to X$ un recubrimiento y sea $X''$ un espacio simplemente conexo, conexo y localmente conexo por arcos. Dada una aplicación continua $f: X'' \to X$, y puntos elegidos $a'' \in X''$ y $a' \in X'$ sujetos a la compatibilidad $p(a') = f(a'')$, se asegura la existencia de un único levantamiento continuo $f': X'' \to X'$ de $f$ tal que $f'(a'') = a'$.
\end{proposicion}
\begin{proof}
La demostración transcurre rigurosamente en tres pasos.
\paragraph{Paso 1: Construcción puntual de $f'$.} Sea $x'' \in X''$ un punto cualquiera. Por la conexidad por arcos de $X''$, fijamos una curva $\gamma : [0,1] \to X''$ tal que $\gamma(0) = a''$ y $\gamma(1) = x''$. La composición $f \circ \gamma$ es una curva continua en $X$ comenzando en $f(a'')$. Levantamos esta curva a $X'$ empezando en $a'$ usando la CLP de $p$, obteniendo una curva levantada $\phi$. Definimos entonces el valor $f'(x'') = \phi(1)$.

\paragraph{Paso 2: Buena definición y dependencia del camino.} Supongamos que $\gamma_0$ y $\gamma_1$ son dos curvas en $X''$ desde $a''$ hasta $x''$. Como el espacio $X''$ es estrictamente simplemente conexo, ambas curvas son homotópicas con extremos fijos. Sea $F$ dicha homotopía en $X''$. Entonces la composición $f \circ F$ proporciona una homotopía continua explícita en $X$ entre las curvas proyectadas $f \circ \gamma_0$ y $f \circ \gamma_1$, también con extremos fijos. 
Aquí invocamos de lleno el Teorema de Monodromía: dos curvas homotópicas en el espacio base con extremos fijos, al ser levantadas al espacio recubridor desde el mismo punto inicial $a'$, convergerán infaliblemente al mismo punto final en $X'$. Esto demuestra incuestionablemente que $\phi_0(1) = \phi_1(1)$, y por ende $f'(x'')$ está impecablemente bien definida, dependiendo solo del punto final $x''$.

\paragraph{Paso 3: Continuidad de $f'$.} Consideramos un punto arbitrario $x'' \in X''$ y su imagen propuesta $f'(x'') \in X'$. Seleccionamos un entorno abierto $U'$ alrededor de $f'(x'')$ tal que la restricción $p|_{U'} : U' \to U$ (donde $U = p(U')$) es un homeomorfismo perfecto. Puesto que $f$ es continua, podemos extraer un entorno conexo $V''$ de $x''$ en $X''$ que satisfaga el mapeo $f(V'') \subset U$. 
Para cualquier punto vecino $y'' \in V''$, la curva hacia $y''$ puede construirse astutamente concatenando nuestra curva fija a $x''$ con un pequeño segmento enteramente contenido en $V''$. El levantamiento de este segmento final recae exclusivamente en $U'$ por la unicidad en recubrimientos trivializados. Así, $f'(V'') \subset U'$ y localmente $f' = (p|_{U'})^{-1} \circ f$, una composición de dos funciones continuas en una vecindad de $x''$. Esto corona la prueba de continuidad analítica de $f'$.
\end{proof}

\begin{teorema}[2.5.2 - Existencia analítica de primitivas]
Sea $\Omega \subset \mathbb{C}$ un dominio simplemente conexo y sea $f \in \mathcal{H}(\Omega)$ holomorfa. Entonces, $f$ admite una primitiva holomorfa incondicional $F \in \mathcal{H}(\Omega)$ tal que $F'(z) = f(z)$ para todo $z \in \Omega$.
\end{teorema}
\begin{proof}
Presentamos dos abordajes rigurosos en paralelo:
\paragraph{Enfoque 1 (Abstracto vía Haces):} Consideramos el recubrimiento $d : \mathcal{O} \to \Omega$ para el haz de gérmenes con primitiva. La sección local está dada por $i_*(f)$. Aplicamos la Proposición anterior al mapa identidad $i: \Omega \to \Omega$ usando que $\Omega$ es simplemente conexo. Obtenemos una sección continua global que proporciona gérmenes de primitiva en todo el dominio. La holomorfía de estos gérmenes se ensambla en una función primitiva global.

\paragraph{Enfoque 2 (Constructivo de Análisis Real-Complejo):} Fijamos un punto base $a \in \Omega$. Para cualquier $z \in \Omega$, definimos la función $F(z) = \int_{\gamma_z} f(w) dw$, donde $\gamma_z$ es cualquier curva diferenciable a trozos conectando $a$ con $z$. Por la Forma Homotópica del Teorema de Cauchy y la simple conexidad de $\Omega$, el valor de esta integral es absolutamente independiente de la selección de $\gamma_z$, certificando su buena definición estricta. 
Probamos su holomorfía y derivada. Fijemos $z_0 \in \Omega$. Tomamos un disco contenido en $\Omega$ centrado en $z_0$. Para un $z$ en dicho disco, calculamos la tasa de cambio:
$$ \frac{F(z) - F(z_0)}{z - z_0} = \frac{1}{z-z_0} \int_{[z_0, z]} f(w) dw $$
(aprovechando que podemos concatenar y cancelar curvas usando el segmento rectilíneo cerrado $[z_0, z]$ dentro del disco).
Restando $f(z_0)$:
$$ \left| \frac{F(z) - F(z_0)}{z - z_0} - f(z_0) \right| = \left| \frac{1}{z-z_0} \int_{[z_0, z]} (f(w) - f(z_0)) dw \right| \le \max_{w \in [z_0, z]} |f(w) - f(z_0)| $$
Puesto que $f$ es una función holomorfa, y a fortiori continua, el supremo del lado derecho tiende irrefutablemente a $0$ conforme $z \to z_0$. Se concluye con absoluto rigor analítico que $F'(z_0) = f(z_0)$.
\end{proof}

\begin{corolario}[Existencia universal de Logaritmos]
Sea $\Omega \subset \mathbb{C}$ un dominio simplemente conexo que excluye el origen ($0 \notin \Omega$). Entonces se garantiza la existencia de una función holomorfa $h: \Omega \to \mathbb{C}$ que satisface $e^{h(z)} = z$ en todo el dominio $\Omega$. Esta función $h$ es llamada "rama holomorfa del logaritmo". Además, cualesquiera dos de estas ramas difieren por la constante imaginaria pura $2\pi i k$ para algún $k \in \mathbb{Z}$ estático.
\end{corolario}
\begin{proof}
\textbf{Existencia:} La función racional $1/z$ es eminentemente holomorfa en $\Omega$, el cual evade el polo en el origen. Al ser $\Omega$ simplemente conexo, por el magno Teorema 2.5.2 recién probado, la función $1/z$ admite una primitiva holomorfa que denotaremos provisoriamente $g \in \mathcal{H}(\Omega)$, es decir, $g'(z) = 1/z$.
Introducimos una función auxiliar inteligente $\phi(z) = z e^{-g(z)}$. Derivamos usando la regla del producto y la cadena:
$$ \phi'(z) = 1 \cdot e^{-g(z)} + z \cdot e^{-g(z)}(-g'(z)) = e^{-g(z)}(1 - z \cdot \frac{1}{z}) = 0 $$
Una función con derivada nula en un dominio conexo es incuestionablemente constante. Así, $z e^{-g(z)} = c$. Evaluando en un $z_0$ arbitrario, vemos que $c \neq 0$. Por lo tanto, escribimos la constante como $c = e^{c_0}$. Definimos formalmente $h(z) = g(z) + c_0$. Esta $h$ es holomorfa y $z e^{-h(z)} = 1$, o sea $e^{h(z)} = z$.

\textbf{Unicidad modular:} Supongamos dos ramas $h_1$ y $h_2$. Entonces $e^{h_1(z)} = e^{h_2(z)} = z$. Dividiendo da $e^{h_1(z) - h_2(z)} = 1$. De la teoría elemental de la exponencial compleja se extrae que la diferencia en el exponente debe ser $2\pi i k(z)$ para un $k(z) \in \mathbb{Z}$. 
Como $h_1$ y $h_2$ son holomorfas (por ende continuas), la diferencia $h_1 - h_2$ es continua. Una función continua de un conexo a un discreto ($\mathbb{Z}$) está confinada a ser puramente constante. Así, el $k(z) = k \in \mathbb{Z}$ es invariante sobre todo $\Omega$.
\end{proof}

\begin{corolario}[Existencia de Raíces n-ésimas]
Bajo las mismas premisas sobre $\Omega$ ($0 \notin \Omega$ conexo y simplemente conexo), para cualquier $n \in \mathbb{N}$, existe $\phi \in \mathcal{H}(\Omega)$ tal que $\phi(z)^n = z$.
\end{corolario}
\begin{proof}
Simplemente definimos $\phi(z) = e^{h(z)/n}$, donde $h(z)$ es la rama holomorfa del logaritmo asegurada en el corolario predecesor. Evidentemente holomorfa, y mediante el álgebra de exponentes: $\phi(z)^n = (e^{h(z)/n})^n = e^{h(z)} = z$.
\end{proof}

\begin{teorema}[2.5.3 - Teorema de Poincaré-Volterra sobre ramas analíticas]
Sea $\Omega \subset \mathbb{C}$ un dominio abierto y conexo. Considérese un subespacio $X \subset \mathcal{O}$ (del haz de gérmenes analíticos en $\Omega$) asumiendo:
\begin{enumerate}
    \item $X$ es topológicamente conexo bajo la topología inducida de haz.
    \item La restricción de la proyección base $p|_X : X \to \Omega$ actúa como un homeomorfismo local puro.
\end{enumerate}
Bajo estas exigencias, cada fibra singular $X_a = X \cap \mathcal{O}_a$ es a lo sumo un conjunto infinito \textbf{numerable}.
\end{teorema}
\begin{proof}
Presentamos la demostración moderna estructurada:
\paragraph{Parte 1: Análisis topológico discreto de la fibra.} La proyección restringida $p|_X$ es homeomorfismo local. Esto dictamina que para cada germen $f_a \in X$, existe una vecindad $U' \subset X$ proyectada homeomórficamente sobre $U \subset \Omega$. En dicha vecindad recae a lo sumo un germen originado en $a$ (el propio $f_a$), forzando a que la intersección $U' \cap p^{-1}(\{a\}) = \{f_a\}$. 
Esto implica colosalmente que la fibra $X_a$ hereda la topología \textbf{discreta} dentro del subespacio $X$. Adicionalmente, se verifica que la fibra es cerrada.

\paragraph{Parte 2: Argumento dimensional vía axiomas de numerabilidad y series de Taylor.} El plano complejo $\mathbb{C}$ posee una base numerable canónica (discos $D(c, r)$ donde $c$ tiene componentes racionales y $r$ es racional). La topología generada del haz $\mathcal{O}$ hereda este "Segundo Axioma de Numerabilidad", originando que la colección completa de todos los posibles tubos funcionales forme una base topológica numerable.
Alternativamente a la base racional, argumentamos por determinación de Taylor: fijemos $a \in \Omega$. Todo germen en $\mathcal{O}_a$ queda férreamente determinado por su infinita expansión de Taylor local, representada por la sucesión de coeficientes $(c_0, c_1, c_2, \dots) \in \mathbb{C}^\mathbb{N}$. 
El subespacio $X$ es conexo. Cualquier germen $g_a \in X_a$ debe ser alcanzable desde un elemento base $f_a^{(0)} \in X_a$ mediante continuación analítica explícita a lo largo de secuencias finitas de discos sobre curvas compactas (pasos iterativos). Usando que los centros y radios racionales son enumerables, las cadenas topológicas de continuación (rutas de monodromía sobre recubrimientos en la base) configuran una estructura de red (Lema de König en árboles numerablemente generados). 
Acoplando esto al axioma de separabilidad topológica: como el haz total $\mathcal{O}$ (y por tanto su subespacio conexo $X$) detentan el segundo axioma de numerabilidad, todo subespacio discreto encapsulado dentro de él está obligado de facto a ser numerable. Como la fibra $X_a$ es discreta y $X$ tiene base numerable, se extrae concluyentemente que $|X_a| \le \aleph_0$.
\end{proof}

\begin{corolario}[Numerabilidad de las ramas o determinaciones]
Para cualquier función holomorfa multivaluada originada desde un germen $f_a$, la multiplicidad dimensional de sus extensiones (sus diferentes "ramas" de comportamiento) al rodear singularidades y proyectar de regreso a la fibra en cualquier punto $z \in \Omega$, jamás superará el cardinal de los números naturales. El conjunto de ramas es numerable. 
Esto cementa de manera incuestionable la validez de la teoría de funciones multivaluadas.
\end{corolario}

\begin{observacion}[Notas históricas y comentarios finales]
\begin{enumerate}
    \item \textbf{Superficies de Riemann y la motivación original:} El genio matemático Bernhard Riemann (1851) introdujo su magna concepción de "superficies" para inyectar sentido absoluto a funciones endiabladamente multivaluadas, como la raíz $z^{1/2}$ o el logaritmo $\log(z)$. En contraposición a la visión miope de múltiples valores asignados a una única coordenada compleja, Riemann modeló un "nuevo dominio de definición". Al construir una variedad entrelazada, $f(z)$ retomaba su naturaleza pura de función univaluada sobre esta superficie. El haz moderno de gérmenes $\mathcal{O}$ destilado en $\S2.2$ cristaliza y formaliza hasta la última gota esta intuición de Riemann. Los "puntos" topológicos de la superficie son simplemente nuestros gérmenes, y cada componente conexa de $\mathcal{O}$ no es otra cosa que el avatar moderno de una Superficie de Riemann.
    
    \item \textbf{Puntos de ramificación y valor en el infinito:} La noción canónica de "punto de ramificación" (el ilustre \textit{Verzweigungspunkt}) fue la piedra angular del análisis clásico del siglo XIX. En la terminología actual, un punto de ramificación $w_0$ es aquel alrededor del cual, al someter a la función a una continuación analítica a lo largo de un lazo circunscrito (monodromía), el germen final transmuta en uno radicalmente distinto al original. El prototipo inquebrantable es $w_0 = 0$ para $z^{1/2}$. Un lazo $e^{2\pi i t}$ ($t \in [0,1]$) transfiere la rama $+z^{1/2}$ de lleno a $-z^{1/2}$. Nuestro Teorema de Monodromía ($\S2.4$) dicta cuándo \textbf{NO} emergen estos cataclismos: un dominio simplemente conexo exorciza la ramificación. Al erradicar los lazos libres, la continuación de un germen es garantizadamente estable e invariante, colapsando a monodromía trivial.
    
    \item \textbf{Weierstrass y la continuación analítica metódica:} Karl Weierstrass (1842-1880) fue el arquitecto principal de la continuación analítica a través de sus formidables \textit{Funktionselemente} (elementos de función), idénticos al concepto subyacente de "germen" empleado aquí. Su genial tesis era que la "función analítica completa" surgía orgánicamente como la clase de equivalencia colosal reuniendo todos los elementos alcanzables mediante prolongación encadenada. Hoy reconocemos esa totalidad como la componente conexa del haz $\mathcal{O}$. Irónicamente, Weierstrass operó sin una formalización topológica de conjunto para su ente; la topología inducida sobre su colección la debemos históricamente a las revoluciones posteriores de Leray (1945) y la máquina de Haces de H. Cartan.
    
    \item \textbf{El grupo fundamental y el paradigma de Poincaré:} Henri Poincaré (1895), cimentando los pilares de la topología algebraica, dio a luz al "grupo fundamental" $\pi_1(X, a)$. El Teorema de Monodromía delinea el maridaje profundo entre el álgebra y el análisis complejo: la falta de trivialidad en $\pi_1(\Omega)$ cifra la verdadera obstrucción topológica para el monovalorismo de las funciones. La correspondencia unívoca donde la monodromía dictamina las proyecciones globales se celebra bajo la "Correspondencia de Riemann-Hilbert". Monodromía plana impera si y solo si el lazo pertenece a la identidad de $\pi_1$.
    
    \item \textbf{El asombroso Teorema de Poincaré-Volterra:} Atribuido independientemente a Poincaré y Vito Volterra entre 1888 y 1895, el Teorema 2.5.3 impone una asombrosa "finitud estructural" o de numerabilidad. Resulta verdaderamente desconcertante el hecho de que, en el incontable continuo subyacente de todo el plano complejo $\mathbb{C}$, la proliferación de determinaciones mutantes en un mismo punto no pueda jamás explotar hacia densidades incontables ($\aleph_1$ o mayores). La prueba extrae su magia del esqueleto racional, apalancándose en la dimensión finita por componente de la serie de Taylor, demostrando que el análisis se encuentra sometido a un corsé algebraico y separabilidad topológica.
    
    \item \textbf{Referencia formativa externa (Conway):} J. B. Conway, en su texto \textit{Functions of One Complex Variable} (GTM 11), despacha el material de monodromía y recubrimientos (capítulo IX) gravitando notablemente hacia el andamiaje puro de la topología algebraica. Nuestra exposición de las notas, tributo a la senda de Narasimhan, es abrumadoramente más directa, encastrada en el análisis y enarbolando el haz de gérmenes con maestría inigualable. Para el purista que anhele dominar el grupo fundamental subyacente o las construcciones del recubrimiento universal abstracto, la lectura de Conway se perfila como un excelente contraste a esta destilación analítica.
    
    \item \textbf{El legado de Cartan y los Haces Multivariables:} H. Cartan y J.P. Serre (1950-1955) inyectaron esteroides a la teoría de haces para acometer la hidra de múltiples variables complejas ($\mathbb{C}^n$). Nuestro haz $\mathcal{O}$ modesto sirve como entrenamiento para los formidables "Haces Coherentes", donde los supremos Teoremas A y B de Cartan suplantan al mero Teorema de Cauchy, dando respuesta a problemas de interpelación topológica y cohomología dimensional infinita. Este capítulo debe visualizarse como el rito de paso ineludible al análisis pluricomplejo de la alta escuela francesa.
\end{enumerate}
\end{observacion}


\end{document}
